% Leçon 15 — Vocabulaire et compréhension de texte : protection de l’environnement — Chinois 1ère A — Cours signé OnBuch+
% Programme officiel MINESEC. Compiler avec : tectonic ZHA15-vocabulaire-protection-de-lenvironnement.tex
\newcommand{\DOCMATIERE}{Chinois}
\newcommand{\DOCNIVEAU}{1ère A}
\newcommand{\DOCMODULE}{Module 2 : Environnement et santé}
\newcommand{\DOCLECON}{ZHA15}
\newcommand{\DOCTITRE}{Leçon 15 — Vocabulaire et compréhension de texte : protection de l’environnement}
% OnBuch+ — préambule « cours signés » ENRICHI (compilé avec tectonic / XeLaTeX)
% Système visuel « Pop » d'OnBuch+ : fond crème (jamais blanc), bordures encre
% épaisses, ombres « dures » décalées, titres Archivo Black en capitales.
% Version MATHÉMATIQUES : graphes (pgfplots/tikz), boîtes pédagogiques
% (définition, propriété, méthode, attention, le savais-tu, exercice résolu,
% exercice, TP, bilan), page de garde et signature OnBuch+ ancrée en bas.
%
% Macros attendues dans main.tex AVANT \input{preamble} :
%   \DOCMATIERE  \DOCNIVEAU  \DOCMODULE  \DOCLECON (numéro)  \DOCTITRE
\documentclass[11pt]{article}

\usepackage{fontspec}
\usepackage[french]{babel} % français = langue principale ; le chinois via \zh{..} / textebox (xeCJK)
\usepackage{xeCJK}
\usepackage{pifont} % \ding{51} \ding{55} utilisés par les modèles
\usepackage[a4paper,margin=2cm,top=3.4cm,bottom=2.8cm,headheight=1.4cm,footskip=1.3cm]{geometry}
\usepackage{amsmath,amssymb,amsfonts}
\usepackage{mathtools} % \xmapsto, \coloneqq... souvent utilisés par les modèles
\usepackage{cancel} % simplifications barrées (bilans de forces)
% \abs{z} (module, valeur absolue) et \norm{u} (norme), utilisés par les modèles
\providecommand{\abs}{}\let\abs\relax\DeclarePairedDelimiter{\abs}{\lvert}{\rvert}
\providecommand{\norm}{}\let\norm\relax\DeclarePairedDelimiter{\norm}{\lVert}{\rVert}
\usepackage[table]{xcolor} % [table] : fournit aussi \rowcolors (alternance de lignes)
\usepackage{pagecolor}
\usepackage{tikz}
\usetikzlibrary{arrows.meta,positioning,calc,shapes.geometric,shapes.misc,shapes.arrows,decorations.pathmorphing,decorations.pathreplacing,decorations.markings,patterns,shadows,mindmap,trees,fit,backgrounds,matrix,intersections,angles,quotes}
\usepackage{pgfplots}
\usepackage[version=4]{mhchem} % équations nucléaires \ce{^{235}_{92}U}
\usepackage[european,siunitx]{circuitikz} % schémas électriques
\pgfplotsset{compat=1.18}
\usepgfplotslibrary{fillbetween,groupplots} % aires entre courbes (calcul intégral)
\usepackage{enumitem}
\usepackage{fancyhdr}
% [most] charge toutes les bibliothèques tcolorbox (listings, minted, raster,
% documentation...) et gonfle inutilement la mémoire TeX sur un document long
% (~300 boîtes) : on ne charge que ce qui est réellement utilisé.
\usepackage[skins,breakable]{tcolorbox}
\usepackage{graphicx}
\usepackage{colortbl}
\usepackage{booktabs}
\usepackage{tabularx}
\usepackage{diagbox} % case d'en-tête diagonale des tableaux à double entrée
% Colonnes extensibles centrée / gauche / droite (C, L, R), souvent utilisées par les modèles
\newcolumntype{C}{>{\centering\arraybackslash}X}
\newcolumntype{L}{>{\raggedright\arraybackslash}X}
\newcolumntype{R}{>{\raggedleft\arraybackslash}X}
\usepackage{array}
\usepackage{multicol}
\usepackage{multirow}
\usepackage{siunitx}
% parse-numbers=false : accepte aussi \SI{2,5 \times 10^{-3}}{...}
\sisetup{locale=FR,output-decimal-marker={,},group-minimum-digits=4,parse-numbers=false}
\DeclareSIUnit\ohm{\ensuremath{\symup{\Omega}}} % U+2126 absent de la police
\DeclareSIUnit\division{div} % réglages d'oscilloscope (ms/div, V/div)
\DeclareSIUnit\tr{tr} % tours (vitesse de rotation, ex. \SI{1500}{\tr\per\minute})
\DeclareSIUnit\molar{M} % concentration molaire (ex. \SI{3}{\molar}, \SI{1}{\micro\molar})
\DeclareSIUnit\an{an} % année (le modèle confond parfois avec \year, un compteur TeX) — ex. \SI{5}{\kilo\meter\per\an}
\DeclareSIUnit\FCFA{FCFA} % franc CFA (ex. \SI{500}{\FCFA\per\kilo\gram})
\DeclareSIUnit\degreeBrix{\textdegree Brix} % teneur en sucre d'un sirop/jus (réfractométrie)
\DeclareSIUnit\month{mois} % le modèle utilise parfois \month, un compteur TeX, comme unité de durée
\usepackage{qrcode}
% \degree : commande fréquemment utilisée par les modèles pour un angle ou une
% température en dehors de \SI{}{} (non fournie par les paquets chargés).
\providecommand{\degree}{^{\circ}}
% \qed (fin de démonstration), utilisé par les modèles sans amsthm
\providecommand{\qed}{\hfill$\square$}
% \barre : double barre des valeurs interdites dans un tableau de variations (tkz-tab)
\providecommand{\barre}{\ensuremath{\|}}
% environnement proof (amsthm non chargé) utilisé par les modèles
\ifdefined\proof\else\newenvironment{proof}[1][Démonstration]{\par\noindent\textit{#1.}\ }{\qed\par}\fi
\usepackage{lastpage}
\usepackage{hyperref}
\hypersetup{hidelinks}

% ============================================================
% Palette « Pop » OnBuch+ — mêmes hex que la classe P (plus_ui.dart)
% ============================================================
\definecolor{poporange}{HTML}{F59321}
\definecolor{popdark}{HTML}{E07A0C}
\definecolor{popcream}{HTML}{FFF6E8}
\definecolor{popink}{HTML}{1C1714}
\definecolor{poppurple}{HTML}{7A5AE0}
\definecolor{popgreen}{HTML}{2E9E6A}
\definecolor{poppink}{HTML}{E0457B}
\definecolor{popblue}{HTML}{2F7DE1}
\definecolor{popgold}{HTML}{C98A12}
\definecolor{popmuted}{HTML}{8A7F76}
\definecolor{popline}{HTML}{EADFD2}
\definecolor{popgray}{HTML}{8A7F76}
\colorlet{popgrey}{popgray}
% Alias tolérants (le modèle nomme parfois les couleurs différemment)
\colorlet{popwhite}{white}
\colorlet{popblack}{popink}
\colorlet{popred}{poppink}
\colorlet{popyellow}{popgold}
% Teintes claires (fonds de tableaux/encadrés) — le modèle nomme parfois
% les couleurs avec un suffixe L (ex. popblueL) sans les avoir définies.
\colorlet{poporangeL}{poporange!20}
\colorlet{popdarkL}{popdark!20}
\colorlet{poppurpleL}{poppurple!20}
\colorlet{popgreenL}{popgreen!20}
\colorlet{poppinkL}{poppink!20}
\colorlet{popblueL}{popblue!20}
\colorlet{popgoldL}{popgold!20}
\colorlet{popredL}{popred!20}
\colorlet{popgrayL}{popgray!20}
% Teintes claires pour fonds de tableaux / zones
\colorlet{popblueL}{popblue!10!white}
\colorlet{poporangeL}{poporange!14!white}
\colorlet{popgreenL}{popgreen!12!white}
\colorlet{poppurpleL}{poppurple!12!white}
\colorlet{poppinkL}{poppink!12!white}
\colorlet{popgoldL}{popgold!14!white}

\pagecolor{popcream}

% ============================================================
% Typographie
% ============================================================
\defaultfontfeatures{Ligatures=TeX}
\setmainfont{PlusJakartaSans-Medium.ttf}[Path=../../../fonts/,BoldFont=PlusJakartaSans-Bold.ttf]
% Symboles, API et emojis absents de la police principale : repli sur DejaVu Sans (caractères actifs)
\newfontface\symfont{DejaVuSans.ttf}[Path=../../../fonts/]
\makeatletter
\newcommand\symactive[1]{\catcode#1=13 \begingroup\lccode`\~=#1 \lowercase{\endgroup\def~}{{\symfont\char#1}}}
\newcommand\symblank[1]{\catcode#1=13 \begingroup\lccode`\~=#1 \lowercase{\endgroup\def~}{}}
\newcommand\symhyph[1]{\catcode#1=13 \begingroup\lccode`\~=#1 \lowercase{\endgroup\def~}{\mbox{-}}}
\newcount\symc
\newcommand\symrange[2]{\symc=#1 \@whilenum\symc<#2 \do{\expandafter\symactive\expandafter{\the\symc}\advance\symc by 1 }}
\newcommand\symblankrange[2]{\symc=#1 \@whilenum\symc<#2 \do{\expandafter\symblank\expandafter{\the\symc}\advance\symc by 1 }}
\makeatother
\symrange{"0250}{"02FF}\symactive{"2713}\symactive{"2714}\symactive{"2717}\symactive{"2718}\symactive{"2610}\symactive{"2611}\symactive{"2612}
\symrange{"2600}{"27BF}
\catcode"2705=13 \begingroup\lccode`\~="2705 \lowercase{\endgroup\def~}{{\symfont\char"2713}}\symblank{"26BD}\symblank{"26A1}
\symrange{"2460}{"257F}\symactive{"223C}\symactive{"2248}\symactive{"2260}
\symhyph{"2011}\symhyph{"2E1A}\symblankrange{"1F300}{"1FAFF}\symblank{"FE0F}
% Caractères chinois : WenQuanYi Zen Hei (sous-ensemble GB2312 + pinyin), gras/italique simulés
\setCJKmainfont{WenQuanYiZenHei-subset.ttf}[Path=../../../fonts/,AutoFakeBold=3,AutoFakeSlant=0.2]
\xeCJKsetup{CJKspace=false}
% Pinyin : voyelles à caron (ǎ ǐ ǒ ǔ ǖ ǘ ǚ ǜ) absentes de la police principale -> caractères actifs
\newfontface\pyfont{WenQuanYiZenHei-subset.ttf}[Path=../../../fonts/]
\newcommand\pyactive[1]{\catcode#1=13 \begingroup\lccode`\~=#1 \lowercase{\endgroup\def~}{{\pyfont\char#1}}}
\pyactive{"01CD}\pyactive{"01CE}\pyactive{"01CF}\pyactive{"01D0}\pyactive{"01D1}\pyactive{"01D2}\pyactive{"01D3}\pyactive{"01D4}\pyactive{"01D5}\pyactive{"01D6}\pyactive{"01D7}\pyactive{"01D8}\pyactive{"01D9}\pyactive{"01DA}\pyactive{"01DB}\pyactive{"01DC}
% Maths en sans-serif (Fira Math) avec chiffres et lettres droites de Plus
% Jakarta, pour que les formules se fondent dans le texte.
\usepackage[math-style=ISO,bold-style=ISO]{unicode-math}
\setmathfont{FiraMath-Regular.otf}[Path=../../../fonts/]
\setmathfont{PlusJakartaSans-Medium.ttf}[Path=../../../fonts/,range={up/num,up/{Latin,latin},\mathup}]
\usepackage{icomma} % virgule décimale sans espace en mode math
% Caractères mathématiques Unicode absents des polices de texte (Plus Jakarta,
% Archivo Black) : rendus par leur équivalent mathématique, en texte comme en
% formule — sinon ils s'affichent en carré vide (ex. « Arithmétique dans ℤ »).
\usepackage{newunicodechar}
\newunicodechar{ℝ}{\ensuremath{\mathbb{R}}}
\newunicodechar{ℤ}{\ensuremath{\mathbb{Z}}}
\newunicodechar{ℂ}{\ensuremath{\mathbb{C}}}
\newunicodechar{ℕ}{\ensuremath{\mathbb{N}}}
\newunicodechar{ℚ}{\ensuremath{\mathbb{Q}}}
\newunicodechar{𝔻}{\ensuremath{\mathbb{D}}}
\newunicodechar{α}{\ensuremath{\alpha}}
\newunicodechar{β}{\ensuremath{\beta}}
\newunicodechar{γ}{\ensuremath{\gamma}}
\newunicodechar{δ}{\ensuremath{\delta}}
\newunicodechar{ε}{\ensuremath{\varepsilon}}
\newunicodechar{θ}{\ensuremath{\theta}}
\newunicodechar{λ}{\ensuremath{\lambda}}
\newunicodechar{μ}{\ensuremath{\mu}}
\newunicodechar{σ}{\ensuremath{\sigma}}
\newunicodechar{τ}{\ensuremath{\tau}}
\newunicodechar{φ}{\ensuremath{\varphi}}
\newunicodechar{ψ}{\ensuremath{\psi}}
\newunicodechar{ω}{\ensuremath{\omega}}
\newunicodechar{Σ}{\ensuremath{\Sigma}}
% majuscules produites par \MakeUppercase dans les titres (α-aminés -> Α)
\newunicodechar{Α}{\ensuremath{\alpha}}
\newunicodechar{Β}{\ensuremath{\beta}}
\newunicodechar{Γ}{\ensuremath{\Gamma}}
\newunicodechar{Δ}{\ensuremath{\Delta}}
\newunicodechar{Ε}{\ensuremath{\varepsilon}}
\newunicodechar{Θ}{\ensuremath{\Theta}}
\newunicodechar{Λ}{\ensuremath{\Lambda}}
\newunicodechar{Μ}{\ensuremath{\mu}}
\newunicodechar{Τ}{\ensuremath{\tau}}
\newunicodechar{Φ}{\ensuremath{\Phi}}
\newunicodechar{Ψ}{\ensuremath{\Psi}}
\newunicodechar{Ω}{\ensuremath{\Omega}}
\newunicodechar{↦}{\ensuremath{\mapsto}}
\newunicodechar{∈}{\ensuremath{\in}}
\newunicodechar{∉}{\ensuremath{\notin}}
\newunicodechar{∧}{\ensuremath{\wedge}}
\newunicodechar{⇔}{\ensuremath{\Leftrightarrow}}
\newunicodechar{⟺}{\ensuremath{\Longleftrightarrow}}
\newunicodechar{⇒}{\ensuremath{\Rightarrow}}
\newunicodechar{⟹}{\ensuremath{\Longrightarrow}}
\newunicodechar{∘}{\ensuremath{\circ}}
\newunicodechar{∪}{\ensuremath{\cup}}
\newunicodechar{∩}{\ensuremath{\cap}}
\newunicodechar{≡}{\ensuremath{\equiv}}
\newunicodechar{⊂}{\ensuremath{\subset}}
\newunicodechar{⊥}{\ensuremath{\perp}}
\newunicodechar{∃}{\ensuremath{\exists}}
\newunicodechar{∀}{\ensuremath{\forall}}
\newunicodechar{∛}{\ensuremath{\sqrt[3]{\,}}}
\newunicodechar{∜}{\ensuremath{\sqrt[4]{\,}}}
\newunicodechar{⃗}{\ensuremath{{}^{\to}}}
\newunicodechar{ⁿ}{\ifmmode^{n}\else\textsuperscript{\ensuremath{n}}\fi}
\newunicodechar{ˣ}{\ifmmode^{x}\else\textsuperscript{\ensuremath{x}}\fi}
\newunicodechar{ᵇ}{\ifmmode^{b}\else\textsuperscript{\ensuremath{b}}\fi}
\newunicodechar{ʳ}{\ifmmode^{r}\else\textsuperscript{\ensuremath{r}}\fi}
\newunicodechar{ᵘ}{\ifmmode^{u}\else\textsuperscript{\ensuremath{u}}\fi}
\newunicodechar{ʸ}{\ifmmode^{y}\else\textsuperscript{\ensuremath{y}}\fi}
\newunicodechar{ᵃ}{\ifmmode^{a}\else\textsuperscript{\ensuremath{a}}\fi}
\newunicodechar{ᵉ}{\ifmmode^{e}\else\textsuperscript{\ensuremath{e}}\fi}
\newunicodechar{ᵏ}{\ifmmode^{k}\else\textsuperscript{\ensuremath{k}}\fi}
\newunicodechar{ᵖ}{\ifmmode^{p}\else\textsuperscript{\ensuremath{p}}\fi}
\newunicodechar{ᴺ}{\ifmmode^{N}\else\textsuperscript{\ensuremath{N}}\fi}
\newunicodechar{ᶜ}{\ifmmode^{c}\else\textsuperscript{\ensuremath{c}}\fi}
\newunicodechar{ʰ}{\ifmmode^{h}\else\textsuperscript{\ensuremath{h}}\fi}
\newunicodechar{ᵠ}{\ifmmode^{q}\else\textsuperscript{\ensuremath{q}}\fi}
\newunicodechar{ₙ}{\ifmmode_{n}\else\textsubscript{\ensuremath{n}}\fi}
\newunicodechar{ₐ}{\ifmmode_{a}\else\textsubscript{\ensuremath{a}}\fi}
\newunicodechar{ᵢ}{\ifmmode_{i}\else\textsubscript{\ensuremath{i}}\fi}
\newunicodechar{ᵣ}{\ifmmode_{r}\else\textsubscript{\ensuremath{r}}\fi}
\newunicodechar{ₖ}{\ifmmode_{k}\else\textsubscript{\ensuremath{k}}\fi}
\newunicodechar{ᵦ}{\ifmmode_{\beta}\else\textsubscript{\ensuremath{\beta}}\fi}
\newunicodechar{ₓ}{\ifmmode_{x}\else\textsubscript{\ensuremath{x}}\fi}
\newfontfamily\popdisplay{ArchivoBlack-Regular.ttf}[Path=../../../fonts/]
\newfontfamily\poplogo{SpaceGrotesk-Bold.ttf}[Path=../../../fonts/]
\newfontfamily\popsemi{PlusJakartaSans-SemiBold.ttf}[Path=../../../fonts/]
\newfontfamily\popxbold{PlusJakartaSans-ExtraBold.ttf}[Path=../../../fonts/]
\newfontfamily\popxboldls{PlusJakartaSans-ExtraBold.ttf}[Path=../../../fonts/,LetterSpace=3.0]

\setlist[enumerate]{leftmargin=1.7em,itemsep=5pt,topsep=4pt}
\setlist[itemize]{leftmargin=1.7em,itemsep=4pt,topsep=4pt,label={\color{poporange}\textbullet}}
\setlength{\parindent}{0pt}
\setlength{\parskip}{5pt}
\renewcommand{\arraystretch}{1.25}

% pgfplots : style Pop par défaut
\pgfplotsset{
  every axis/.append style={axis line style={popink,line width=1pt},tick style={popink},
    grid style={popline,line width=0.5pt},label style={font=\small},tick label style={font=\footnotesize}},
  popcurve/.style={poporange,line width=1.8pt,no markers,smooth},
}
% Même style disponible dans un \draw TikZ simple (hors pgfplots)
\tikzset{popcurve/.style={poporange,line width=1.8pt}}
\tikzset{popgrid/.style={popline,very thin}}
\colorlet{popcurve}{poporange} % le nom du style est parfois utilisé comme couleur

% ============================================================
% Pill / squiggle
% ============================================================
\newcommand{\poppill}[3]{%
  \tikz[baseline=-0.55ex]\node[draw=popink,line width=1.2pt,rounded corners=2.6pt,%
    fill=#2,inner xsep=6.5pt,inner ysep=3pt]{%
    \popxboldls\fontsize{7}{7}\selectfont\color{#3}\MakeUppercase{#1}};%
}
\newcommand{\popsquiggle}[1]{%
  \tikz[baseline=-0.4ex]{
    \draw[#1,line width=2.6pt,line cap=round]
      plot [smooth] coordinates {(0,0.09) (0.34,0.01) (0.68,0.21) (1.02,0.01) (1.36,0.10)};
  }%
}
% Mot-clé surligné (marqueur orange sous le texte)
\newcommand{\cle}[1]{{\popxbold\color{popdark}#1}}

% ============================================================
% En-tête / pied
% ============================================================
\pagestyle{fancy}
\fancyhf{}
\renewcommand{\headrulewidth}{0pt}
\renewcommand{\footrulewidth}{0pt}
\lhead{\raisebox{-0.15ex}{\poplogo\fontsize{13}{13}\selectfont\color{popink}OnBuch\color{poporange}+}}
\rhead{\poppill{\DOCMATIERE\ \textbullet\ \DOCNIVEAU\ \textbullet\ Leçon \DOCLECON}{white}{popink}}
\lfoot{\popxbold\fontsize{7.6}{7.6}\selectfont\color{popmuted}\MakeUppercase{Cours signé OnBuch+}\ \textbullet\ {\popsemi\DOCTITRE}}
\rfoot{\tikz[baseline=-0.6ex]\node[rounded corners=8.5pt,fill=popink,minimum height=17pt,minimum width=17pt,inner xsep=5pt,inner sep=0]{\popxbold\fontsize{8}{8}\selectfont\color{popcream}\thepage\,/\,\pageref*{LastPage}};}

% ============================================================
% Titres
% ============================================================
\newcommand{\chaptitre}[2]{%
  \begin{tcolorbox}[enhanced,colback=poporange,colframe=popink,boxrule=2.6pt,arc=11pt,
      drop shadow=popink,left=15pt,right=15pt,top=12pt,bottom=11pt,before skip=3pt,after skip=18pt]
    \poppill{#2}{popink}{popcream}\par\vspace{9pt}
    {\raggedright\hyphenpenalty=10000\exhyphenpenalty=10000\popdisplay\fontsize{22}{24}\selectfont\color{popcream}\MakeUppercase{#1}\par}
  \end{tcolorbox}
}
% Section numérotée : pastille orange avec le numéro + titre Archivo
\newcounter{popsec}
\newcommand{\coursec}[1]{%
  \stepcounter{popsec}\setcounter{popsub}{0}%
  \par\vspace{16pt}\noindent
  \begin{minipage}{\linewidth}
  \tikz[baseline=-0.7ex]\node[draw=popink,line width=1.6pt,fill=poporange,rounded corners=4pt,minimum size=22pt,inner sep=0]{\popdisplay\fontsize{12}{12}\selectfont\color{popcream}\thepopsec};%
  \hspace{8pt}{\popdisplay\fontsize{15}{17}\selectfont\color{popink}\MakeUppercase{#1}}\par
  \vspace{4pt}\noindent{\color{poporange}\rule{36pt}{3.6pt}}
  \end{minipage}\par\nopagebreak\vspace{8pt}
}
\newcounter{popsub}
\newcommand{\courssub}[1]{%
  \stepcounter{popsub}%
  \par\vspace{9pt}\noindent{\popxbold\fontsize{11.5}{13}\selectfont\color{popink}{\color{poporange}\thepopsec.\thepopsub}\hspace{6pt}#1}\par\nopagebreak\vspace{4pt}
}

% ============================================================
% Boîtes pédagogiques — même recette : fond blanc, bordure épaisse colorée,
% ombre dure encre, badge plein. Titre optionnel [#1].
% ============================================================
\tcbset{popbase/.style={breakable,enhanced,colback=white,boxrule=2.3pt,arc=9pt,
  shadow={3pt}{-3pt}{0pt}{popink},left=14pt,right=14pt,top=11pt,bottom=11pt,before skip=11pt,after skip=15pt,
  fontupper=\popsemi\fontsize{10.6}{14.5}\selectfont\color{popink},
  fontlower=\popsemi\fontsize{10.6}{14.5}\selectfont\color{popink}}}
% Coche dessinée (le glyphe ✓ n'existe pas dans les polices du document)
\newcommand{\popcheck}[1]{\tikz[baseline=-0.5ex]\draw[#1,line width=1.6pt,line cap=round,line join=round](-0.13,0.0)--(-0.04,-0.09)--(0.13,0.1);}
\providecommand{\coche}{\popcheck{popgreen}}
\providecommand{\resultat}[1]{\par\smallskip\noindent\fcolorbox{popgreen}{popgreenL}{\parbox{\dimexpr\linewidth-2\fboxsep-2\fboxrule\relax}{\textbf{Résultat.} #1}}\par\smallskip}
\newcommand{\popboxhead}[3]{\poppill{#1}{#2}{white}\ifx\relax#3\relax\else\hspace{6pt}{\popxbold\fontsize{10.6}{12}\selectfont\color{popink}#3}\fi\par\vspace{7pt}}

\newtcolorbox{aretenir}[1][]{popbase,colframe=popgreen,before upper={\popboxhead{À retenir}{popgreen}{#1}}}
% Texte en chinois (document de lecture, dialogue, extrait) : cadre
% d'étude. Un titre optionnel (source, auteur) s'affiche dans l'en-tête.
\newtcolorbox{textebox}[1][]{popbase,colframe=popblue,before upper={\popboxhead{文本 · Texte}{popblue}{#1}},after upper={}}
% Marques ✓ / ✗ (forme correcte / fautive) souvent utilisées par les modèles
\providecommand{\cross}{\textcolor{poppink}{\ensuremath{\times}}}
\providecommand{\xmark}{\cross}\providecommand{\crossmark}{\cross}\providecommand{\cmark}{\ensuremath{\checkmark}}
% Ligne de réponse / justification à compléter (inventée par les modèles)
\providecommand{\justification}[1]{\par\noindent\textit{Justification~:} \dotfill\par}
% \textipa (package tipa non chargé) : la police ne couvre pas l'API -> simple passage
\providecommand{\textipa}[1]{#1}
% Soulignement (ulem non chargé) : \uline, \ul
\providecommand{\uline}[1]{\underline{#1}}\providecommand{\ul}[1]{\underline{#1}}
% \glossaire{\item ...} inventée par les modèles : liste de glossaire
\providecommand{\glossaire}[1]{\begin{itemize}#1\end{itemize}}
% \alert (beamer) inventée par les modèles : texte mis en évidence
\providecommand{\alert}[1]{\textbf{\textcolor{poppink}{#1}}}
% variantes ASCII d'ordinaux écrites en macro
\providecommand{\iere}{\textsuperscript{re}}\providecommand{\ieme}{\textsuperscript{e}}\providecommand{\ere}{\textsuperscript{re}}\providecommand{\eme}{\textsuperscript{e}}\providecommand{\er}{\textsuperscript{er}}\providecommand{\ie}{\textsuperscript{e}}
% Ordinaux français écrits en macro par les modèles : 1\ière, 3\ième
\newcommand{\ière}{\textsuperscript{re}}\newcommand{\ième}{\textsuperscript{e}}
% \exemple (inventée par les modèles) : étiquette « Exemple : »
\providecommand{\exemple}{\textbf{Exemple~:}\ }
% \square utilisable aussi hors mode math (cases à cocher)
\AtBeginDocument{\let\squareMath\square\renewcommand{\square}{\ensuremath{\squareMath}}}
% Mot ou phrase en chinois à l'intérieur d'un paragraphe français
\newcommand{\zh}[1]{\ifmmode\text{#1}\else{#1}\fi}
\let\eng\zh \let\esp\zh \let\all\zh % alias tolérés
% flèches d'intonation inventées par les modèles
\providecommand{\textsearrow}{}\renewcommand{\textsearrow}{\ensuremath{\searrow}}\providecommand{\textnearrow}{}\renewcommand{\textnearrow}{\ensuremath{\nearrow}}
% \fr{..} : mot français (inventé par les modèles)
\providecommand{\fr}[1]{\textit{#1}}
% zhbox : encadré de texte chinois inventé par les modèles
\newenvironment{zhbox}{\begin{quote}}{\end{quote}}
% \vs : « versus » inventé par les modèles
\providecommand{\vs}{\textit{vs}\ }
% \blank : case à compléter inventée par les modèles
\providecommand{\blank}[1][]{\underline{\hspace{1.6em}}}
\newtcolorbox{exemplebox}[1][]{popbase,colframe=poporange,before upper={\popboxhead{Exemple}{poporange}{#1}}}
\newtcolorbox{definition}[1][]{popbase,colframe=popblue,before upper={\popboxhead{Définition}{popblue}{#1}}}
\newtcolorbox{propriete}[1][]{popbase,colframe=poppurple,before upper={\popboxhead{Propriété}{poppurple}{#1}}}
\newtcolorbox{methode}[1][]{popbase,colframe=popgold,before upper={\popboxhead{Méthode}{popgold}{#1}}}
\newtcolorbox{attention}[1][]{popbase,colframe=poppink,before upper={\popboxhead{Attention}{poppink}{#1}}}
\newtcolorbox{experience}[1][]{popbase,colframe=popblue,colback=popblueL,before upper={\popboxhead{Expérience}{popblue}{#1}}}
% « Le savais-tu ? » : carte violette pleine (contexte, culture, Cameroun)
\newtcolorbox{savaistu}[1][]{popbase,colback=poppurple,colframe=popink,
  fontupper=\popsemi\fontsize{10.4}{14.2}\selectfont\color{white},
  before upper={\poppill{Le savais-tu\,?}{white}{poppurple}\ifx\relax#1\relax\else\hspace{6pt}{\popxbold\color{white}#1}\fi\par\vspace{7pt}}}
% Exercice résolu : énoncé en haut, \tcblower puis la solution sur fond crème
\newcounter{popexo}\newcounter{popexr}
\newtcolorbox{exoresolu}[1][]{popbase,colframe=poporange,colbacklower=popcream,
  segmentation style={popink,line width=1.2pt,dashed},
  before upper={\stepcounter{popexr}\popboxhead{Exercice résolu \thepopexr}{poporange}{#1}},
  before lower={\poppill{Solution}{popink}{popcream}\par\vspace{6pt}}}
% Exercice (non résolu en place) — la correction est dans la partie Corrigés
\newtcolorbox{exercice}[1][]{popbase,colframe=popink,
  before upper={\stepcounter{popexo}\popboxhead{Exercice \thepopexo}{popink}{#1}}}
% Corrigé
\newtcolorbox{corrige}[1][]{popbase,colframe=popgreen,colback=popgreenL,
  before upper={\popboxhead{Corrigé}{popgreen}{#1}}}
% Objectifs / prérequis / bilan
\newtcolorbox{objectifs}{popbase,colframe=popink,colback=poporangeL,
  before upper={\popboxhead{Objectifs de la leçon}{popink}{}}}
\newtcolorbox{prerequis}{popbase,colframe=popink,colback=popblueL,
  before upper={\popboxhead{Prérequis}{popblue}{}}}
\newtcolorbox{bilan}{popbase,colframe=popink,colback=white,boxrule=2.8pt,
  before upper={\popboxhead{Fiche bilan}{poporange}{L'essentiel à connaître pour l'examen}}}
% Figure encadrée avec légende
\newtcolorbox{popfigure}[1][]{enhanced,breakable=false,colback=white,colframe=popline,boxrule=1.4pt,arc=8pt,
  left=10pt,right=10pt,top=10pt,bottom=8pt,before skip=10pt,after skip=12pt,halign=center,
  lower separated=false,halign lower=center,
  fontlower=\popsemi\fontsize{9}{11.5}\selectfont\color{popmuted},
  #1}
\newcommand{\legende}[1]{\par\vspace{6pt}{\popsemi\fontsize{9}{11.5}\selectfont\color{popmuted}#1}}

% ============================================================
% Signature OnBuch+
% ============================================================
\newcommand{\poplogoinline}[1]{{\poplogo\fontsize{#1}{#1}\selectfont\color{popink}OnBuch\color{poporange}+}}
\newcommand{\signatureonbuch}{%
  \begin{tcolorbox}[enhanced,colback=popink,colframe=popink,boxrule=2.2pt,arc=10pt,
      drop shadow=poporange,left=14pt,right=14pt,top=10pt,bottom=10pt,before skip=0pt,after skip=0pt]
    \tikz[baseline=-0.7ex]\node[circle,fill=popgreen,draw=popcream,line width=1.3pt,minimum size=22pt,inner sep=0]{\popcheck{white}};%
    \hspace{9pt}\begin{minipage}[c]{0.62\linewidth}
      {\popxbold\fontsize{12}{13}\selectfont\color{popcream}Signé\ \textbullet\ {\poplogo OnBuch\color{poporange}+}}\par\vspace{2pt}
      {\raggedright\popsemi\fontsize{8}{10.5}\selectfont\color{popline}\MakeUppercase{Équipe pédagogique OnBuch+}\\\MakeUppercase{Conforme au programme officiel MINESEC}\par}
    \end{minipage}\hfill
    {\poplogo\fontsize{22}{22}\selectfont\color{popcream}OnBuch\color{poporange}+}
  \end{tcolorbox}%
}

% ============================================================
% Page de garde
% #1 titre · #2 matière · #3 niveau · #4 module · #5 n° leçon · #6 durée ·
% #7 description / objectifs (texte ou itemize)
% ============================================================
\newcommand{\pagegarde}[7]{%
  \thispagestyle{empty}
  \newgeometry{a4paper,margin=1.7cm,top=1.6cm,bottom=1.4cm}%
  \begin{tikzpicture}[remember picture,overlay]
    \fill[poporange!18] ([xshift=-3.2cm,yshift=-3.6cm]current page.north east) circle (4.6cm);
    \fill[poppurple!14] ([xshift=2.4cm,yshift=7.6cm]current page.south west) circle (3.8cm);
  \end{tikzpicture}%
  {\poplogo\fontsize{30}{30}\selectfont\color{popink}OnBuch\color{poporange}+}\hfill
  \raisebox{0.6ex}{\poppill{Cours signé \textbullet\ Premium}{popink}{popcream}}\par
  \vspace{1pt}\popsquiggle{poppurple}\par
  \vspace{8pt}
  \begin{tcolorbox}[enhanced,colback=poporange,colframe=popink,boxrule=2.8pt,arc=13pt,
      drop shadow=popink,left=18pt,right=18pt,top=12pt,bottom=13pt,after skip=10pt]
    \poppill{#2 \textbullet\ #3}{popink}{popcream}\hspace{5pt}\poppill{#4}{white}{popink}\par\vspace{8pt}
    {\popdisplay\fontsize{14}{14}\selectfont\color{popink}LEÇON #5}\par\vspace{4pt}
    {\raggedright\hyphenpenalty=10000\exhyphenpenalty=10000\popdisplay\fontsize{25}{27}\selectfont\color{popcream}\MakeUppercase{#1}\par}
    \vspace{8pt}
    {\popxbold\fontsize{9.5}{9.5}\selectfont\color{popink}\MakeUppercase{Durée conseillée : #6}}
  \end{tcolorbox}
  \begin{tcolorbox}[enhanced,colback=white,colframe=popink,boxrule=2.2pt,arc=10pt,
      drop shadow=popink,left=16pt,right=16pt,top=10pt,bottom=10pt,after skip=8pt]
    \poppill{Dans cette leçon}{poppurple}{white}\par\vspace{5pt}
    \popsemi\fontsize{9.3}{12.6}\selectfont\color{popink}#7
  \end{tcolorbox}
  \noindent\begin{minipage}[t]{0.487\linewidth}
    \begin{tcolorbox}[enhanced,colback=popgreen,colframe=popink,boxrule=2.2pt,arc=10pt,
        drop shadow=popink,left=11pt,right=11pt,top=10pt,bottom=10pt,height=3.1cm,valign=center]
      \tcbox[colback=white,colframe=popink,boxrule=1.3pt,arc=5pt,boxsep=0pt,nobeforeafter,
        left=3pt,right=3pt,top=3pt,bottom=3pt,valign=center]{\qrcode[height=1.9cm]{https://play.google.com/store/apps/details?id=com.luvvix.onbuch&hl=fr}}%
      \hspace{8pt}\begin{minipage}[c]{0.5\linewidth}
        {\popdisplay\fontsize{11}{12.5}\selectfont\color{white}\MakeUppercase{Télécharge OnBuch}}\par\vspace{4pt}
        {\raggedright\popsemi\fontsize{8.4}{11}\selectfont\color{white}Gratuit \textbullet\ Google Play\\Tuteur IA, annales, concours, résultats\par}
      \end{minipage}
    \end{tcolorbox}
  \end{minipage}\hfill
  \begin{minipage}[t]{0.487\linewidth}
    \begin{tcolorbox}[enhanced,colback=poppurple,colframe=popink,boxrule=2.2pt,arc=10pt,
        drop shadow=popink,left=11pt,right=11pt,top=10pt,bottom=10pt,height=3.1cm,valign=center]
      \tikz[baseline=-0.7ex]\node[circle,fill=white,draw=popink,line width=1.3pt,minimum size=1.6cm,inner sep=0]{\popdisplay\fontsize{24}{24}\selectfont\color{poppurple}+};%
      \hspace{8pt}\begin{minipage}[c]{0.6\linewidth}
        {\popdisplay\fontsize{11}{12.5}\selectfont\color{white}\MakeUppercase{Passe à OnBuch+}}\par\vspace{4pt}
        {\raggedright\popsemi\fontsize{8.4}{11}\selectfont\color{white}Tous les cours signés, Tuteur IA illimité, fiches PDF — onglet OnBuch+ de l'app\par}
      \end{minipage}
    \end{tcolorbox}
  \end{minipage}
  \vspace{8pt}
  \popcontenu
  \vfill
  \signatureonbuch
  \restoregeometry
  \newpage
}

% Rangée « Ce cours contient » (page de garde)
\newcommand{\poptile}[3]{% #1 couleur #2 gros texte #3 légende
  \begin{tcolorbox}[enhanced,colback=#1,colframe=popink,boxrule=2pt,arc=9pt,shadow={2.5pt}{-2.5pt}{0pt}{popink},
      left=6pt,right=6pt,top=9pt,bottom=9pt,halign=center,valign=center,height=1.75cm,width=\linewidth,nobeforeafter]
    {\popdisplay\fontsize{12.5}{13}\selectfont\color{white}\MakeUppercase{#2}}\par\vspace{3pt}
    {\centering\popsemi\fontsize{7.8}{9.5}\selectfont\color{white}#3\par}
  \end{tcolorbox}}
\newcommand{\popcontenu}{%
  \par\noindent{\popxboldls\fontsize{7.5}{7.5}\selectfont\color{popmuted}CE COURS SIGNÉ CONTIENT}\par\vspace{7pt}
  \noindent\begin{minipage}[t]{0.235\linewidth}\poptile{popblue}{Cours}{complet, illustré, pas à pas}\end{minipage}\hfill
  \begin{minipage}[t]{0.235\linewidth}\poptile{poporange}{Méthodes}{et exercices résolus}\end{minipage}\hfill
  \begin{minipage}[t]{0.235\linewidth}\poptile{poppink}{Exercices}{type examen + corrigés}\end{minipage}\hfill
  \begin{minipage}[t]{0.235\linewidth}\poptile{popgreen}{Fiche bilan}{l'essentiel à retenir}\end{minipage}\par}

% Sommaire
\newtcolorbox{sommaire}{enhanced,colback=white,colframe=popink,boxrule=2.2pt,arc=10pt,
  drop shadow=popink,left=18pt,right=18pt,top=13pt,bottom=13pt,before skip=6pt,after skip=20pt,
  before upper={\poppill{Sommaire}{poppurple}{white}\par\vspace{9pt}\popsemi\fontsize{10.6}{14.5}\selectfont\color{popink}}}

% Signature de fin de cours — poussée tout en bas de la dernière page
\newcommand{\findecours}{%
  \par\vfill\nopagebreak
  \begin{center}\popsquiggle{poporange}\end{center}\vspace{4pt}
  \signatureonbuch
}
\AtBeginDocument{\def\llbracket{\mathopen{[\mkern-3.5mu[}}\def\rrbracket{\mathclose{]\mkern-3.5mu]}}} % intervalles d'entiers (glyphe ⟦ absent de la police)
% Glyphes absents de Fira Math : redéfinis à partir de glyphes disponibles
\AtBeginDocument{%
  \def\setminus{\mathbin{\backslash}}\let\smallsetminus\setminus
  \def\vdots{\mathord{\vbox{\baselineskip3.2pt\lineskiplimit0pt\kern4pt\hbox{.}\hbox{.}\hbox{.}}}}%
  \def\ddots{\mathinner{\mkern1mu\raise6.5pt\hbox{.}\mkern2mu\raise3.6pt\hbox{.}\mkern2mu\raise0.7pt\hbox{.}\mkern1mu}}%
  \def\triangle{\mathord{\scalebox{0.9}{$\symup{\Delta}$}}}%
  \def\nabla{\mathord{\raisebox{\depth}{\scalebox{1}[-1]{$\symup{\Delta}$}}}}%
  \def\checkmark{\popcheck{popdark}}%
  \def\star{\mathbin{\ast}}\def\bigstar{\mathord{\ast}}%
  \def\diamond{\mathbin{\scalebox{0.6}{$\Diamond$}}}%
  \def\bigcup{\mathop{\mathchoice{\raisebox{-0.3em}{\scalebox{1.7}{$\cup$}}}{\scalebox{1.25}{$\cup$}}{\cup}{\cup}}\displaylimits}%
  \def\bigcap{\mathop{\mathchoice{\raisebox{-0.3em}{\scalebox{1.7}{$\cap$}}}{\scalebox{1.25}{$\cap$}}{\cap}{\cap}}\displaylimits}%
  \def\lesseqqgtr{\mathrel{\lessgtr}}\def\bigoplus{\mathop{\oplus}\displaylimits}%
}
\providecommand{\ssi}{si et seulement si} % abréviation fréquente
\AtBeginDocument{\providecommand{\wideparen}{\overparen}} % arc de cercle

\begin{document}

\pagegarde{Leçon 15 — Vocabulaire et compréhension de texte : protection de l’environnement}{Chinois}{1ère A}{Module 2 : Environnement et santé}{ZHA15}{2 h}{Cette leçon permet de lire et comprendre un court texte chinois sur la protection de l'environnement et d'acquérir le vocabulaire essentiel du thème (pollution, recyclage, économie d'eau et d'énergie). Elle développe la prononciation correcte des mots, la reconnaissance des radicaux et la capacité à répondre à des questions de compréhension.
\vspace{4pt}
\begin{multicols}{2}\small
\begin{itemize}
\item À la fin de cette leçon, je sais lire à voix haute un court texte chinois sur la protection de l'environnement avec une prononciation correcte (pinyin et tons).
\item À la fin de cette leçon, je sais reconnaître et traduire au moins 15 mots de vocabulaire liés à l'environnement.
\item À la fin de cette leçon, je sais identifier les radicaux (clés) des caractères du thème et écrire correctement 5 caractères en respectant l'ordre des traits.
\item À la fin de cette leçon, je sais répondre par écrit à des questions de compréhension portant sur un texte chinois.
\item À la fin de cette leçon, je sais employer les collocations courantes du thème (保护环境, 节约用水, 乱扔垃圾) dans des phrases complètes.
\item À la fin de cette leçon, je sais comparer brièvement les problèmes environnementaux du Cameroun et de la Chine avec le vocabulaire appris.
\end{itemize}\end{multicols}}

\begin{sommaire}
\begin{multicols}{2}
\begin{enumerate}
\item Mise en route et découverte du thème
\item Texte support : lecture et compréhension globale
\item Glossaire du thème
\item Radicaux et ordre des traits
\item Collocations et questions de compréhension
\item Production guidée et ancrage socioculturel
\item Activité d'intégration
\item Méthodes et exercices résolus
\item Exercices
\item Corrigés des exercices
\item Fiche bilan
\end{enumerate}
\end{multicols}
\end{sommaire}

\begin{objectifs}
\begin{itemize}
\item À la fin de cette leçon, je sais lire à voix haute un court texte chinois sur la protection de l'environnement avec une prononciation correcte (pinyin et tons).
\item À la fin de cette leçon, je sais reconnaître et traduire au moins 15 mots de vocabulaire liés à l'environnement.
\item À la fin de cette leçon, je sais identifier les radicaux (clés) des caractères du thème et écrire correctement 5 caractères en respectant l'ordre des traits.
\item À la fin de cette leçon, je sais répondre par écrit à des questions de compréhension portant sur un texte chinois.
\item À la fin de cette leçon, je sais employer les collocations courantes du thème (保护环境, 节约用水, 乱扔垃圾) dans des phrases complètes.
\item À la fin de cette leçon, je sais comparer brièvement les problèmes environnementaux du Cameroun et de la Chine avec le vocabulaire appris.
\item À la fin de cette leçon, je sais produire un court paragraphe guidé (4-5 phrases) sur un geste écologique.
\end{itemize}
\end{objectifs}

\begin{prerequis}
\begin{itemize}
\item Maîtrise du pinyin et des quatre tons (Seconde)
\item Structures de base : sujet + verbe + complément ; 是, 有, 很 (Seconde)
\item Vocabulaire de la nature et de la météo (leçons antérieures du Module 2)
\item Notion de radicaux (clés) et ordre général des traits (Seconde)
\item Modalité simple avec 应该 / 可以 (début d'année de Première)
\end{itemize}
\end{prerequis}

\newpage

\coursec{Mise en route et découverte du thème}

\courssub{Situation d'accroche et problématique}

Chaque matin à Yaoundé, des milliers d'élèves traversent à pied des quartiers où les ordures s'accumulent au bord des rues : sachets plastiques, bouteilles vides, restes de nourriture. Chaque saison sèche, dans le Nord-Cameroun, les feux de brousse détruisent des champs entiers et la fumée rend l'air difficile à respirer. À Douala, après une forte pluie, les caniveaux bouchés par les déchets provoquent des inondations. Tous ces problèmes ont un nom commun : ce sont des problèmes d'\cle{environnement}.

De l'autre côté du monde, la Chine a connu des difficultés semblables. Les grandes villes comme Pékin ont longtemps souffert de la pollution de l'air, au point que les habitants portaient souvent des masques. Puis le gouvernement chinois a lancé de vastes campagnes de protection de l'environnement : tri des déchets, énergies propres, plantations d'arbres. Aujourd'hui, les jeunes Chinois parlent de \zh{环保} (\emph{huánbǎo}, protection de l'environnement) comme les jeunes Camerounais parlent de recyclage et de reboisement.

\textbf{Problématique de la leçon :} comment dire en chinois « Nous devons protéger notre environnement » ? Quels mots utiliser pour parler de la pollution, des déchets, des arbres et de l'eau ? C'est ce que nous allons apprendre à partir d'un texte intitulé \zh{保护我们的环境} (\emph{Bǎohù wǒmen de huánjìng} — « Protégeons notre environnement »).

\courssub{Brainstorming : les problèmes environnementaux au Cameroun}

Avant de lire le moindre caractère chinois, parlons d'abord en français. Regardez autour de vous, dans votre quartier, votre ville, votre village : quels problèmes environnementaux observez-vous ?

\begin{experience}[Brainstorming en classe entière]
L'enseignant écrit au tableau les réponses des élèves et les regroupe en trois familles :
\begin{enumerate}
\item \textbf{Les déchets plastiques} : sachets, bouteilles, emballages jetés dans les rues de Douala ou de Yaoundé ;
\item \textbf{La déforestation} : coupe des arbres pour le bois de chauffe, feux de brousse, disparition des forêts du Sud ;
\item \textbf{La pollution de l'eau} : rivières sales, puits contaminés, plastiques dans le fleuve Wouri.
\end{enumerate}
Consigne : chaque élève cite un problème qu'il a vu de ses propres yeux cette semaine. L'enseignant note au moins huit réponses au tableau.
\end{experience}

Ces trois familles de problèmes ont leur équivalent exact en chinois, et ce sont précisément les mots du texte que nous allons étudier. Retenez déjà cette correspondance :

\begin{center}
\begin{tabularx}{\linewidth}{|X|X|X|X|}
\hline
\rowcolor{popblueL} Caractères & Pinyin & Nature & Traduction \\
\hline
环境 & huánjìng & nom & environnement \\
\hline
污染 & wūrǎn & nom / verbe & pollution ; polluer \\
\hline
垃圾 & lājī & nom & déchets, ordures \\
\hline
森林 & sēnlín & nom & forêt \\
\hline
水 & shuǐ & nom & eau \\
\hline
空气 & kōngqì & nom & air \\
\hline
保护 & bǎohù & verbe & protéger \\
\hline
节约 & jiéyuē & verbe & économiser (eau, énergie) \\
\hline
\end{tabularx}
\end{center}

\begin{attention}[Prononciation : les pièges pour les francophones]
\begin{itemize}
\item \zh{环境} \emph{huánjìng} : la lettre \textbf{j} se prononce comme un « dj » doux, jamais comme le « j » français de « jour ».
\item \zh{水} \emph{shuǐ} : \textbf{sh} se prononce comme « ch » français, mais la langue est recourbée vers le palais.
\item \zh{垃圾} \emph{lājī} : deux premiers tons, bien monter la voix sur les deux syllabes : \emph{lājī}, et non « ladji » plat.
\item \zh{保护} \emph{bǎohù} : deux troisième et quatrième tons ; \emph{bǎo} descend puis remonte, \emph{hù} tombe franchement.
\end{itemize}
\end{attention}

\courssub{Observation d'images et élicitation du vocabulaire}

L'enseignant présente maintenant trois images (ou les dessine schématiquement au tableau) : une ville polluée avec de la fumée, des élèves qui trient des déchets dans trois poubelles de couleurs, et une grande forêt verte. Pour chaque image, il pose la question : « Qu'est-ce que vous voyez ? » puis il donne le mot chinois.

\begin{exemplebox}[De l'image au mot chinois]
Image 1 (ville grise, fumée) : \zh{空气污染。} \emph{Kōngqì wūrǎn.} « L'air est pollué. »\\
Image 2 (trois poubelles) : \zh{垃圾分类。} \emph{Lājī fēnlèi.} « Le tri des déchets. »\\
Image 3 (forêt verte) : \zh{保护森林。} \emph{Bǎohù sēnlín.} « Protéger la forêt. »
\end{exemplebox}

Remarquez la logique de construction des mots chinois : \zh{垃圾} (déchets) + \zh{分类} (classer) = \zh{垃圾分类} (tri des déchets) ; \zh{空气} (air) + \zh{污染} (pollution) = \zh{空气污染} (pollution de l'air). Le chinois assemble des mots simples pour créer des mots composés, comme des briques. C'est une excellente nouvelle pour vous : chaque nouveau mot appris multiplie votre vocabulaire.

\begin{popfigure}
\begin{center}
\begin{tikzpicture}[mindmap, grow cyclic, every node/.style={concept, align=center}, root concept/.append style={concept color=popblue, font=\large\bfseries}, level 1/.append style={level distance=3.2cm, sibling angle=90}, level 1 concept/.append style={concept color=poporange!70, font=\small}]
\node[root concept, align=center]{环境\\ huánjìng\\ environnement}
  child { node[align=center]{污染\\ wūrǎn\\ pollution} }
  child { node[align=center]{垃圾\\ lājī\\ déchets} }
  child { node[align=center]{节约\\ jiéyuē\\ économiser} }
  child { node[align=center]{保护\\ bǎohù\\ protéger} };
\end{tikzpicture}
\end{center}
\legende{Carte mentale du champ lexical de l'environnement : un mot central, quatre familles de sens.}
\end{popfigure}

\begin{methode}[Comment mémoriser ce champ lexical]
\begin{enumerate}
\item Partez du mot central \zh{环境} (environnement) et prononcez-le à voix haute trois fois.
\item Parcourez chaque branche de la carte mentale en associant le mot à une image mentale : la fumée pour \zh{污染}, une poubelle pour \zh{垃圾}, un robinet fermé pour \zh{节约}, un bouclier pour \zh{保护}.
\item Construisez une phrase simple avec chaque branche : \zh{我们要保护环境。} \emph{Wǒmen yào bǎohù huánjìng.} « Nous devons protéger l'environnement. »
\item Récitez la carte mentale les yeux fermés, puis redessinez-la de mémoire dans votre cahier.
\end{enumerate}
\end{methode}

\courssub{Annonce du titre et prédiction}

Voici maintenant le titre du texte que nous allons étudier dans cette leçon :

\begin{textebox}[Titre du texte]
保护我们的环境\\
\emph{Bǎohù wǒmen de huánjìng}\\
« Protégeons notre environnement »
\end{textebox}

Décortiquons ce titre mot à mot, car il contient déjà une grammaire précieuse :

\begin{center}
\begin{tabularx}{\linewidth}{|X|X|X|X|}
\hline
\rowcolor{popblueL} Caractères & Pinyin & Nature & Traduction \\
\hline
保护 & bǎohù & verbe & protéger \\
\hline
我们 & wǒmen & pronom & nous \\
\hline
的 & de & particule & (marque de possession / de relation) \\
\hline
环境 & huánjìng & nom & environnement \\
\hline
\end{tabularx}
\end{center}

\begin{propriete}[La particule 的 \emph{de} entre un pronom et un nom]
Structure : \textbf{pronom / nom + 的 + nom}. La particule \zh{的} relie deux noms (ou un pronom et un nom) et exprime la possession ou l'appartenance, comme « de » ou l'apostrophe « s » en français.
\begin{itemize}
\item \zh{我们的环境} \emph{wǒmen de huánjìng} : « notre environnement » (litt. « l'environnement de nous ») ;
\item \zh{我的书} \emph{wǒ de shū} : « mon livre » ;
\item \zh{喀麦隆的森林} \emph{Kāmàilóng de sēnlín} : « les forêts du Cameroun ».
\end{itemize}
Attention à l'ordre : en chinois, le possesseur vient \textbf{toujours avant} le possédé, alors qu'en français avec « de » c'est l'inverse : \zh{我们的学校} = « l'école \emph{de} nous » = « notre école ».
\end{propriete}

\begin{attention}[Erreur fréquente des francophones]
Ne dites jamais \zh{环境的我们} en calquant « l'environnement de nous ». En chinois, on inverse : le possesseur d'abord, \zh{的} au milieu, le possédé ensuite. Mémorisez le schéma : \textbf{A 的 B = B de A}.
\end{attention}

\begin{experience}[Prédiction : que va dire le texte ?]
À partir du titre \zh{保护我们的环境} et des trois images observées, les élèves formulent des prédictions en français, puis l'enseignant les aide à les transformer en chinois simple :
\begin{itemize}
\item « Le texte va parler de la pollution » : \zh{空气和水有污染。} \emph{Kōngqì hé shuǐ yǒu wūrǎn.}
\item « Il va parler des déchets » : \zh{垃圾很多。} \emph{Lājī hěn duō.} « Il y a beaucoup de déchets. »
\item « Il va donner des conseils » : \zh{我们要节约水。} \emph{Wǒmen yào jiéyuē shuǐ.} « Nous devons économiser l'eau. »
\end{itemize}
Chaque prédiction est notée au tableau ; à la lecture du texte (section suivante), la classe vérifiera quelles prédictions étaient justes.
\end{experience}

\courssub{Première écoute : saisir le sens global}

L'enseignant lit maintenant le texte \textbf{une première fois, à voix haute et à vitesse naturelle, sans que les élèves aient le texte sous les yeux}. Cahiers fermés, oreilles ouvertes. L'objectif n'est pas de tout comprendre, mais de répondre à trois questions globales :

\begin{enumerate}
\item De quoi parle le texte ? (le sujet général)
\item Quels mots de la carte mentale avez-vous reconnus ?
\item Le ton du texte est-il plutôt inquiet ou plutôt optimiste ?
\end{enumerate}

\begin{methode}[Aborder un texte chinois en quatre étapes]
\begin{enumerate}
\item \textbf{Lire le titre} : il annonce le sujet et contient souvent le mot-clé du texte (ici \zh{保护}, protéger).
\item \textbf{Observer les images} : elles donnent le contexte et activent le vocabulaire déjà connu.
\item \textbf{Écouter sans lire} : une première écoute globale pour saisir le sens général, sans s'arrêter sur les mots inconnus.
\item \textbf{Lire en repérant les mots connus} : souligner les mots déjà appris, deviner les autres grâce au contexte, puis seulement consulter le glossaire.
\end{enumerate}
Règle d'or : ne jamais commencer par traduire mot à mot. En chinois comme en français, on comprend d'abord le sens global, puis on précise.
\end{methode}

\begin{savaistu}[Le tri des déchets en Chine]
En Chine, le tri des déchets, \zh{垃圾分类} (\emph{lājī fēnlèi}), est obligatoire dans les grandes villes comme Shanghai depuis 2019. Les habitants doivent séparer leurs ordures en quatre catégories : déchets recyclables, déchets dangereux, déchets de cuisine et autres déchets. Des caméras et des agents vérifient les poubelles, et les contrevenants paient une amende ! Au Cameroun, des associations de jeunes à Yaoundé et à Douala organisent aussi des collectes de plastique et des campagnes de nettoyage des quartiers. Protéger l'environnement est devenu un combat commun aux jeunes des deux pays.
\end{savaistu}

\begin{exoresolu}[Première écoute : questionnaire global]
Énoncé : après la première écoute du texte lue par l'enseignant, répondez en français. 1) Quel est le sujet général du texte ? 2) Citez trois mots chinois reconnus à l'écoute. 3) Le texte se termine-t-il sur une idée positive ou négative ?
\tcblower
Solution détaillée : 1) Le texte parle de la protection de l'environnement : la pollution de l'air et de l'eau, les déchets, et les actions pour protéger la nature. 2) Mots reconnaissables grâce à la carte mentale : \zh{环境} (environnement), \zh{污染} (pollution), \zh{垃圾} (déchets), \zh{保护} (protéger), \zh{节约} (économiser) — trois de ces mots suffisent. 3) Le texte se termine sur une idée positive : si chacun fait un petit effort, notre environnement sera plus beau. Astuce : à la première écoute, repérez les mots déjà connus et la mélodie générale des phrases ; ne cherchez pas à comprendre chaque caractère.
\end{exoresolu}

\begin{exercice}[À vous de jouer]
Sans regarder le tableau de vocabulaire, complétez les phrases suivantes avec le mot chinois qui convient : \zh{污染}, \zh{垃圾}, \zh{森林}, \zh{保护}.
\begin{enumerate}
\item Les feux de brousse détruisent la \_\_\_\_\_\_\_. (forêt)
\item Il y a trop de \_\_\_\_\_\_\_ dans les rues de Douala. (déchets)
\item La fumée des voitures, c'est de la \_\_\_\_\_\_\_. (pollution)
\item Nous devons \_\_\_\_\_\_\_ notre environnement. (protéger)
\end{enumerate}
\end{exercice}

\begin{corrige}[Exercice « À vous de jouer »]
\begin{enumerate}
\item \zh{森林} (\emph{sēnlín}) : les feux de brousse détruisent la forêt.
\item \zh{垃圾} (\emph{lājī}) : il y a trop de déchets dans les rues.
\item \zh{污染} (\emph{wūrǎn}) : la fumée des voitures, c'est de la pollution.
\item \zh{保护} (\emph{bǎohù}) : nous devons protéger notre environnement — phrase complète : \zh{我们要保护我们的环境。} \emph{Wǒmen yào bǎohù wǒmen de huánjìng.}
\end{enumerate}
\end{corrige}

\begin{aretenir}[Ce qu'il faut retenir de cette mise en route]
\begin{itemize}
\item Le mot central du thème est \zh{环境} (\emph{huánjìng}, environnement), entouré de quatre familles : \zh{污染} (pollution), \zh{垃圾} (déchets), \zh{节约} (économiser), \zh{保护} (protéger).
\item Le titre du texte est \zh{保护我们的环境} : « Protégeons notre environnement ».
\item La particule \zh{的} suit le schéma \textbf{A 的 B = B de A} : le possesseur avant, le possédé après.
\item Pour aborder un texte chinois : titre, images, écoute globale, puis lecture en repérant les mots connus.
\end{itemize}
\end{aretenir}

\coursec{Texte support : lecture et compréhension globale}

Dans la section précédente, nous avons découvert le thème du jour à travers des images et des mots déjà familiers : \zh{水} (\emph{shuǐ}, l'eau), \zh{树} (\emph{shù}, l'arbre), \zh{人} (\emph{rén}, la personne). Nous avons aussi entendu pour la première fois le mot-clé de la leçon : \zh{环境} (\emph{huánjìng}, « environnement »). Il est maintenant temps de rencontrer le texte support de la leçon. Notre objectif n'est pas encore de tout analyser mot à mot — ce sera le travail des sections suivantes — mais de \cle{lire}, de \cle{comprendre globalement} et de \cle{repérer} ce que nous connaissons déjà.

\courssub{Première écoute : le texte sans le regarder}

Avant même de lire les caractères, fermez votre manuel et écoutez simplement le professeur lire le texte à voix haute, deux fois. À la première écoute, essayez seulement de répondre à cette question : de quoi parle ce texte ? À la deuxième écoute, comptez mentalement combien de phrases vous entendez.

\begin{textebox}[Texte support — 保护环境]
我们的环境越来越差了。\\
Wǒmen de huánjìng yuèláiyuè chà le.\\
« Notre environnement devient de plus en plus mauvais. »\\[0.4em]
空气和水都被污染了，很多人乱扔垃圾。\\
Kōngqì hé shuǐ dōu bèi wūrǎn le, hěn duō rén luàn rēng lājī.\\
« L'air et l'eau sont pollués, et beaucoup de gens jettent leurs déchets n'importe où. »\\[0.4em]
我们应该保护地球：节约用水、节约用电，还要多种树。\\
Wǒmen yīnggāi bǎohù dìqiú: jiéyuē yòng shuǐ, jiéyuē yòng diàn, hái yào duō zhòng shù.\\
« Nous devons protéger la Terre : économiser l'eau, économiser l'électricité, et aussi planter plus d'arbres. »\\[0.4em]
保护环境，人人有责！\\
Bǎohù huánjìng, rénrén yǒuzé!\\
« Protéger l'environnement est l'affaire de tous ! »
\end{textebox}

\begin{methode}[Comment aborder un texte chinois inconnu ?]
Voici la démarche en quatre étapes que nous appliquerons à chaque texte de l'année :
\begin{enumerate}
\item \textbf{Lire une première fois sans s'arrêter}, pour saisir le sens général. Ne bloquez jamais sur un caractère inconnu.
\item \textbf{Souligner les mots inconnus} au crayon, sans encore chercher leur sens.
\item \textbf{Deviner le sens grâce au contexte} : les mots qui entourent le mot inconnu donnent souvent la réponse. Par exemple, dans \zh{空气和水都被污染了}, même sans connaître \zh{污染}, on comprend qu'il arrive quelque chose de négatif à l'air et à l'eau.
\item \textbf{Consulter le glossaire} seulement après avoir essayé de deviner, pour vérifier son hypothèse.
\end{enumerate}
Cette méthode entraîne votre autonomie : à l'examen, vous rencontrerez toujours des mots nouveaux, et le contexte sera votre meilleur allié.
\end{methode}

\courssub{Lecture en écho : travailler la prononciation et les tons}

Nous allons maintenant lire le texte à voix haute selon la technique de la \cle{lecture en écho}, en quatre temps :

\begin{enumerate}
\item \textbf{Enseignant → classe entière} : le professeur lit une phrase, la classe répète en chœur, en imitant exactement la mélodie des tons.
\item \textbf{Classe → groupes} : chaque rangée lit une phrase à tour de rôle ; les autres groupes écoutent et lèvent la main s'ils entendent une erreur de ton.
\item \textbf{Groupes → binômes} : à deux, chacun lit le texte entier pendant que l'autre suit du doigt et corrige.
\item \textbf{Lecture individuelle} : deux ou trois volontaires lisent seuls devant la classe.
\end{enumerate}

\begin{attention}[Les tons difficiles de ce texte]
Surveillez particulièrement :
\begin{itemize}
\item \zh{越来越} \emph{yuèláiyuè} : deux quatrièmes tons entourent un deuxième ton. Descendez bien franchement sur les deux \zh{越}.
\item \zh{差} \emph{chà} : quatrième ton, bref et sec. Attention : \zh{差} (\emph{chà}, 4\textsuperscript{e} ton) signifie « mauvais, médiocre ». Ne le confondez pas avec \zh{差不多} (\emph{chàbuduō}, « presque, à peu près »), que vous rencontrerez plus tard et qui a un sens très différent !
\item \zh{污染} \emph{wūrǎn} : premier ton puis troisième ton — le troisième ton doit plonger puis remonter légèrement.
\item \zh{节约} \emph{jiéyuē} : deuxième ton montant puis premier ton haut et plat.
\item \zh{种} \emph{zhòng} : quatrième ton (verbe « planter »). Ne pas confondre avec \zh{种} \emph{zhǒng} (3\textsuperscript{e} ton, nom « espèce, graine »).
\end{itemize}
\end{attention}

\courssub{Repérage : mots connus et mots nouveaux}

Relisez maintenant le texte en silence et complétez mentalement le travail suivant : entourez les mots que vous connaissez déjà, soulignez les mots nouveaux. Voici le bilan attendu sous forme de tableau de vocabulaire (la colonne « Nature » indique la classe grammaticale).

\begin{center}
\begin{tabularx}{\linewidth}{|X|X|X|X|}
\hline
\rowcolor{popblueL} Caractères & Pinyin & Nature & Traduction \\
\hline
我们 & wǒmen & pronom (代词) & nous \\
\hline
的 & de & particule (助词) & (marque de possession / attribut) \\
\hline
环境 & huánjìng & nom (名词) & environnement \\
\hline
越来越 & yuèláiyuè & adverbe (副词) & de plus en plus \\
\hline
差 & chà & adjectif verbe (形容词) & mauvais, médiocre \\
\hline
了 & le & particule (助词) & (marque de changement d'état / aspect accompli) \\
\hline
空气 & kōngqì & nom (名词) & air \\
\hline
和 & hé & conjonction (连词) & et \\
\hline
水 & shuǐ & nom (名词) & eau \\
\hline
都 & dōu & adverbe (副词) & tous, toutes \\
\hline
被 & bèi & préposition (介词) & (marque du passif) \\
\hline
污染 & wūrǎn & verbe (动词) & polluer ; pollution \\
\hline
很 & hěn & adverbe (副词) & très \\
\hline
多 & duō & adjectif verbe (形容词) & beaucoup, nombreux \\
\hline
人 & rén & nom (名词) & personne, gens \\
\hline
乱 & luàn & adverbe (副词) & en désordre, n'importe comment \\
\hline
扔 & rēng & verbe (动词) & jeter \\
\hline
垃圾 & lājī & nom (名词) & déchets, ordures \\
\hline
应该 & yīnggāi & verbe modal (能愿动词) & devoir, il faut \\
\hline
保护 & bǎohù & verbe (动词) & protéger \\
\hline
地球 & dìqiú & nom (名词) & la Terre (planète) \\
\hline
节约 & jiéyuē & verbe (动词) & économiser \\
\hline
用 & yòng & verbe (动词) & utiliser \\
\hline
电 & diàn & nom (名词) & électricité \\
\hline
还 & hái & adverbe (副词) & aussi, en plus \\
\hline
要 & yào & verbe modal (能愿动词) & vouloir, devoir \\
\hline
多 & duō & adverbe (副词) & plus (en quantité) \\
\hline
种 & zhòng & verbe (动词) & planter \\
\hline
树 & shù & nom (名词) & arbre \\
\hline
人人 & rénrén & pronom (代词) & tout le monde, chacun \\
\hline
有 & yǒu & verbe (动词) & avoir \\
\hline
责 & zé & nom (名词) & responsabilité (abrégé de \zh{责任}) \\
\hline
\end{tabularx}
\end{center}

Bonne nouvelle : vous constatez que le texte est bâti sur une charpente de mots déjà connus (\zh{我们}, \zh{的}, \zh{水}, \zh{和}, \zh{人}, \zh{了}, \zh{很}, \zh{多}, \zh{用}, \zh{树}, \zh{有}). Les mots nouveaux seront étudiés en détail dans la section 3 (Glossaire du thème) ; pour l'instant, il suffit de les reconnaître et de comprendre leur rôle dans la phrase.

\courssub{Observation grammaticale : la structure passive avec 被}

Regardons de près cette phrase du texte :

\zh{空气和水都被污染了。}

\emph{Kōngqì hé shuǐ dōu bèi wūrǎn le.}

« L'air et l'eau sont pollués. »

En français, nous utilisons le passif (« sont pollués »). Le chinois fait de même avec le mot-outil \zh{被} (\emph{bèi}), qui introduit la voix passive.

\begin{propriete}[Le passif avec 被 (bèi)]
\textbf{Structure :} Sujet (celui qui subit) + \zh{被} + (Agent) + Verbe + Complément / \zh{了} / \zh{完} etc.

\textbf{Emploi :} on l'utilise quand le sujet \emph{subit} l'action, surtout pour des événements négatifs ou subis (polluer, casser, voler, critiquer, frapper). L'agent peut être omis s'il est inconnu, indéfini ou sans importance (passif « sans agent »).

\textbf{Comparaison avec le français :} le français conjugue « être + participe passé » ; le chinois, lui, ne conjugue rien : \zh{被} reste invariable, et ce sont les particules d'aspect (\zh{了}, \zh{完}, \zh{过}…) qui situent l'action dans le temps.
\end{propriete}

\begin{exemplebox}[Exemples de passif avec 被]
\zh{空气被污染了。} \emph{Kōngqì bèi wūrǎn le.} « L'air est pollué. » (agent omis)\\[0.3em]
\zh{水被污染了。} \emph{Shuǐ bèi wūrǎn le.} « L'eau est polluée. »\\[0.3em]
\zh{我的书被弟弟拿走了。} \emph{Wǒ de shū bèi dìdi ná zǒu le.} « Mon livre a été emporté par mon petit frère. » (agent \zh{弟弟} exprimé)\\[0.3em]
\zh{他被批评了。} \emph{Tā bèi pīpíng le.} « Il a été critiqué. »
\end{exemplebox}

\begin{attention}[Piège fréquent : 被 vs 把]
Ne confondez pas \zh{被} (passif : le sujet \textbf{subit}) et \zh{把} (actif : le sujet \textbf{agit} sur l'objet mis en avant).
\begin{itemize}
\item \zh{我把书看了。} (\emph{Wǒ bǎ shū kàn le.}) « J'ai lu le livre. » (Actif, focus sur l'objet).
\item \zh{书被我看了。} (\emph{Shū bèi wǒ kàn le.}) « Le livre a été lu par moi. » (Passif).
\end{itemize}
Pour l'instant, retenez seulement à \textbf{reconnaître} \zh{被} + Verbe + \zh{了} = passif accompli.
\end{attention}

\courssub{Compréhension globale : problèmes et solutions}

Un bon lecteur sait résumer un texte en une phrase et en dégager la structure logique. Ce texte est construit de façon très claire : il présente d'abord des \cle{problèmes} (constat alarmant), puis des \cle{solutions} (actions concrètes), et se termine par un \cle{slogan} mobilisateur.

\begin{popfigure}
\begin{tikzpicture}[
  box/.style={draw=popblue, fill=popblueL, rounded corners, align=center, text width=4.6cm, minimum height=1.1cm, font=\small},
  title/.style={font=\bfseries\small, text=popdark},
  prob/.style={box, fill=poppinkL, draw=poppink},
  sol/.style={box, fill=popgreenL, draw=popgreen}
]
\node[title] at (-3.2,1.6) {问题 Problèmes};
\node[title] at (3.2,1.6) {办法 Solutions};
\node[prob, align=center] at (-3.2,0.6){空气污染\\ kōngqì wūrǎn\\ pollution de l'air};
\node[prob, align=center] at (-3.2,-0.8){水污染\\ shuǐ wūrǎn\\ pollution de l'eau};
\node[prob, align=center] at (-3.2,-2.2){乱扔垃圾\\ luàn rēng lājī\\ jeter les déchets n'importe où};
\node[sol, align=center] at (3.2,0.6){节约用水\\ jiéyuē yòng shuǐ\\ économiser l'eau};
\node[sol, align=center] at (3.2,-0.8){节约用电\\ jiéyuē yòng diàn\\ économiser l'électricité};
\node[sol, align=center] at (3.2,-2.2){多种树\\ duō zhòng shù\\ planter plus d'arbres};
\draw[-{Stealth}, very thick, poporange] (-0.6,-0.8) -- (0.6,-0.8) node[midway, above, font=\tiny, text=poporange] {action};
\end{tikzpicture}
\legende{Structure du texte : à chaque problème, une solution correspondante.}
\end{popfigure}

\begin{exoresolu}[Questions de compréhension globale]
Énoncé — Répondez en français, puis essayez de formuler la réponse en chinois avec les mots du texte.
\begin{enumerate}
\item De quoi parle ce texte ?
\item Quels sont les problèmes environnementaux cités ?
\item Quelles sont les trois solutions proposées par le texte ?
\item Que signifie la dernière phrase \zh{保护环境，人人有责！} ?
\end{enumerate}
\tcblower
Solution détaillée :
\begin{enumerate}
\item Le texte parle de la \textbf{protection de l'environnement} : la situation actuelle se dégrade (pollution, incivilités), mais chacun a le pouvoir et le devoir d'agir. En chinois : \zh{这个课文讲保护环境。} (\emph{Zhège kèwén jiǎng bǎohù huánjìng.})
\item Trois problèmes sont cités : la pollution de l'air (\zh{空气污染}), la pollution de l'eau (\zh{水污染}) et un problème de comportement civique : beaucoup de gens jettent leurs déchets n'importe où (\zh{很多人乱扔垃圾}).
\item Les trois solutions concrètes sont : économiser l'eau (\zh{节约用水}), économiser l'électricité (\zh{节约用电}) et planter plus d'arbres (\zh{多种树}).
\item \zh{保护环境，人人有责！} signifie « Protéger l'environnement, c'est la responsabilité de chacun. » \zh{人人} (\emph{rénrén}, « tout le monde », redoublement de \zh{人} pour exprimer la généralité) + \zh{有责} (\emph{yǒuzé}, « avoir une responsabilité », abréviation de \zh{有责任}).
\end{enumerate}
\end{exoresolu}

\begin{textebox}[Un slogan célèbre à retenir]
保护环境，人人有责！\\
Bǎohù huánjìng, rénrén yǒuzé!\\
« Protéger l'environnement, c'est la responsabilité de chacun ! »\\[0.4em]
Ce slogan est omniprésent en Chine : affiches dans les écoles, les parcs, les rues, les médias. Retenez-le par cœur : il est parfait pour conclure une production écrite ou un exposé sur l'environnement, et il montre à l'examinateur que vous maîtrisez un élément de culture chinoise partagé.
\end{textebox}

\begin{attention}[Pièges de traduction pour les francophones]
\begin{itemize}
\item \zh{越来越差了} se traduit par « devient de plus en plus mauvais » : \zh{越来越} + adjectif exprime une évolution progressive, et le \zh{了} final marque le changement d'état (nouvelle situation). Ne traduisez pas mot à mot par « de plus en plus mauvais est ».
\item \zh{乱扔垃圾} : \zh{乱} (\emph{luàn}) signifie « en désordre, n'importe comment ». Placé \emph{avant} le verbe, il signifie « faire quelque chose n'importe comment / sans ordre » : \zh{乱扔} = « jeter n'importe où », \zh{乱写} = « griffonner », \zh{乱跑} = « courir dans tous les sens ».
\item \zh{节约用水} se construit \zh{节约} (économiser) + \zh{用} (utiliser) + \zh{水} (eau). On retrouve la même structure pour \zh{节约用电} (« économiser l'électricité »). Ne dites pas \zh{节约水用} : l'ordre est strictement \cle{Verbe d'action + 用 + Nom de la ressource}.
\item \zh{多种树} : \zh{多} (plus) + \zh{种} (planter, 4\textsuperscript{e} ton) + \zh{树} (arbre). Attention à ne pas lire \zh{种} \emph{zhǒng} (3\textsuperscript{e} ton, « espèce »).
\end{itemize}
\end{attention}

\begin{exercice}[À vous de jouer]
\begin{enumerate}
\item Lisez le texte à voix haute trois fois, en veillant particulièrement aux tons de \zh{越来越} (4-2-4), \zh{污染} (1-3), \zh{节约} (2-1) et \zh{种} (4\textsuperscript{e} ton).
\item Recopiez \textbf{toutes les phrases} du texte qui contiennent au moins un mot déjà connu (\zh{我们}, \zh{水}, \zh{树}, \zh{人}, \zh{和}, \zh{了}, \zh{很}, \zh{多}, \zh{用}, \zh{有}) et traduisez-les en français.
\item Sans regarder le texte, citez de mémoire les trois problèmes et les trois solutions, en français puis en chinois (caractères + pinyin).
\item Traduisez en chinois :
      \begin{enumerate}
      \item L'eau est polluée.
      \item Nous devons économiser l'eau.
      \item L'air et l'eau sont pollués.
      \item Tout le monde doit protéger la Terre.
      \end{enumerate}
\end{enumerate}
\end{exercice}

\begin{corrige}[Exercice « À vous de jouer »]
\begin{enumerate}
\item Lecture orale vérifiée en classe par le professeur (insister sur l'enchaînement \zh{yuè-lái-yuè} et le ton 3 de \zh{rǎn}).
\item Phrases du texte contenant des mots connus (les 4 phrases en contiennent toutes) :
      \begin{itemize}
      \item \zh{我们的环境越来越差了。} — « Notre environnement devient de plus en plus mauvais. » (\zh{我们}, \zh{了})
      \item \zh{空气和水都被污染了，很多人乱扔垃圾。} — « L'air et l'eau sont pollués, et beaucoup de gens jettent leurs déchets n'importe où. » (\zh{和}, \zh{水}, \zh{了}, \zh{很}, \zh{多}, \zh{人})
      \item \zh{我们应该保护地球：节约用水、节约用电，还要多种树。} — « Nous devons protéger la Terre : économiser l'eau, économiser l'électricité, et aussi planter plus d'arbres. » (\zh{我们}, \zh{用}, \zh{水}, \zh{用}, \zh{电}, \zh{树})
      \item \zh{保护环境，人人有责！} — « Protéger l'environnement, c'est la responsabilité de chacun ! » (\zh{人} dans \zh{人人}, \zh{有})
      \end{itemize}
\item Problèmes : \zh{空气污染} (kōngqì wūrǎn), \zh{水污染} (shuǐ wūrǎn), \zh{乱扔垃圾} (luàn rēng lājī). Solutions : \zh{节约用水} (jiéyuē yòng shuǐ), \zh{节约用电} (jiéyuē yòng diàn), \zh{多种树} (duō zhòng shù).
\item \begin{enumerate}
      \item \zh{水被污染了。} (\emph{Shuǐ bèi wūrǎn le.})
      \item \zh{我们应该节约用水。} (\emph{Wǒmen yīnggāi jiéyuē yòng shuǐ.})
      \item \zh{空气和水都被污染了。} (\emph{Kōngqì hé shuǐ dōu bèi wūrǎn le.})
      \item \zh{人人都应该保护地球。} (\emph{Rénrén dōu yīnggāi bǎohù dìqiú.})
      \end{enumerate}
\end{enumerate}
\end{corrige}

Nous avons maintenant une vision d'ensemble du texte : son sujet, sa structure problèmes/solutions et son slogan final. Dans la prochaine section, nous entrerons dans le détail du vocabulaire grâce au glossaire du thème, pour pouvoir réemployer chaque mot dans vos propres phrases.

\coursec{Glossaire du thème}

\courssub{Transition depuis la lecture globale}

Après avoir découvert le texte \emph{« Protégeons notre Terre »} dans la section précédente et en avoir saisi le sens global, nous allons maintenant \textbf{fixer le lexique} indispensable pour parler de l'environnement en chinois. Ce glossaire structuré vous permettra de reconnaître, prononcer, écrire et employer les mots-clés du thème, tout en évitant les pièges classiques des apprenants francophones (tons, polyphonie, collocations).

\courssub{Tableau récapitulatif du vocabulaire de base}

Voici les \textbf{16 mots fondamentaux} du thème, classés par familles sémantiques pour faciliter la mémorisation. Le tableau utilise les quatre colonnes standard : Caractères, Pinyin (tons marqués), Nature grammaticale, Traduction française.

\begin{center}
\begin{tabularx}{\linewidth}{|X|X|X|X|}
\hline
\rowcolor{popblueL} \textbf{Caractères} & \textbf{Pinyin} & \textbf{Nature} & \textbf{Traduction} \\ \hline
\zh{环境} & huánjìng & n. & environnement \\ \hline
\zh{保护} & bǎohù & v. & protéger \\ \hline
\zh{污染} & wūrǎn & v./n. & polluer / pollution \\ \hline
\zh{空气} & kōngqì & n. & air \\ \hline
\zh{垃圾} & lājī & n. & déchets, ordures \\ \hline
\zh{扔} & rēng & v. & jeter \\ \hline
\zh{乱} & luàn & adv. & n'importe comment, en désordre \\ \hline
\zh{节约} & jiéyuē & v. & économiser \\ \hline
\zh{地球} & dìqiú & n. & la Terre \\ \hline
\zh{树} & shù & n. & arbre \\ \hline
\zh{种} & zhòng & v. & planter \\ \hline
\zh{越来越} & yuèláiyuè & expr. & de plus en plus \\ \hline
\zh{应该} & yīnggāi & v. mod. & devoir \\ \hline
\zh{责任} & zérèn & n. & responsabilité \\ \hline
\zh{干净} & gānjìng & adj. & propre \\ \hline
\zh{回收} & huíshōu & v. & recycler (complément examen) \\ \hline
\end{tabularx}
\end{center}

\courssub{Analyse détaillée par familles sémantiques}

Pour aller au-delà de la liste brute, observons comment ces mots s'articulent logiquement : \textbf{problèmes} (pollution, déchets), \textbf{solutions} (protéger, recycler, planter, économiser), \textbf{acteurs et lieux} (Terre, air, arbres, responsabilité).

\begin{definition}[Famille « Problèmes : Pollution et Déchets »]
    \textbf{污染 wūrǎn} (v./n.) est le mot-clé central. Comme verbe : « polluer » ; comme nom : « pollution ».
    \begin{itemize}
        \item \zh{空气污染} kōngqì wūrǎn — pollution de l'air (litt. « air pollution »).
        \item \zh{水污染} shuǐ wūrǎn — pollution de l'eau.
    \end{itemize}
    \textbf{垃圾 lājī} (n.) « déchets, ordures ». \textbf{Attention :} c'est un mot « inséparable » (voir \emph{Le savais-tu ?}).
    \textbf{扔 rēng} (v.) « jeter ». Souvent suivi de \textbf{乱 luàn} (adv.) pour dire « jeter n'importe où / en désordre » : \zh{乱扔垃圾} luàn rēng lājī — jeter les déchets n'importe comment.
\end{definition}

\begin{definition}[Famille « Solutions : Actions positives »]
    \textbf{保护 bǎohù} (v.) « protéger ». Structure fréquente : \zh{保护环境} bǎohù huánjìng (protéger l'environnement).
    \textbf{节约 jiéyuē} (v.) « économiser ». Souvent suivi d'un objet : \zh{节约用水} jiéyuē yòng shuǐ (économiser l'eau), \zh{节约用电} jiéyuē yòng diàn (économiser l'électricité).
    \textbf{种 zhòng} (v.) « planter (un arbre, des fleurs) ». \textbf{Attention à la polyphonie} (voir \emph{Attention} ci-dessous).
    \textbf{回收 huíshōu} (v.) « recycler, récupérer ». Complément examen fréquent : \zh{回收塑料} huíshōu sùliào (recycler le plastique).
\end{definition}

\begin{definition}[Famille « Acteurs, Lieux, Devoirs »]
    \textbf{环境 huánjìng} (n.) « environnement » (litt. « entourage + situation »).
    \textbf{空气 kōngqì} (n.) « air » (litt. « air/vide + gaz/souffle »).
    \textbf{地球 dìqiú} (n.) « la Terre » (litt. « terre + boule »).
    \textbf{树 shù} (n.) « arbre ». Classificateur : \zh{棵} kē (\zh{一棵树} yī kē shù).
    \textbf{应该 yīnggāi} (v. mod.) « devoir, doit ». Place : \textbf{Sujet + 应该 + Verbe}. Ex: \zh{我们应该保护环境。}
    \textbf{责任 zérèn} (n.) « responsabilité ». Collocation : \zh{尽责任} jìn zérèn (assumer sa responsabilité), \zh{有责任} yǒu zérèn (avoir la responsabilité).
    \textbf{干净 gānjìng} (adj.) « propre, net ». Antonyme : \zh{脏} zāng (sale).
    \textbf{越来越 yuèláiyuè} (adv.) « de plus en plus ». Structure : \zh{越来越 + Adjectif/Verbe}. Ex: \zh{天气越来越热。} Tiānqì yuèláiyuè rè. (Il fait de plus en plus chaud.)
\end{definition}

\courssub{Entraînement à la prononciation : points de vigilance}

Les tons sont la mélodie du chinois. Trois mots de cette liste piègent systématiquement les francophones. Exercez-vous à les dire à voix haute en exagérant la courbe de ton.

\begin{propriete}[Points de prononciation critiques]
    \begin{enumerate}
        \item \textbf{保护 bǎohù} : 3e ton (bas-descendant-remontant) + 4e ton (chutant fort). \emph{Piège :} ne pas dire « bao-hu » (2+1) ni « bao-hu » (3+2). Le 3e ton isolé est bas ; enchaîné devant un 4e ton, il reste bas (demi-3e ton).
        \item \textbf{垃圾 lājī} : 1er ton (haut plat) + 1er ton (haut plat). \emph{Piège :} le « ji » ne se prononce pas « dji » (français) mais \cle{[tɕi]} (consonne alvéolo-palatale sourde, comme un « tchi » très doux non aspiré). Les deux syllabes restent hautes.
        \item \textbf{节约 jiéyuē} : 2e ton (montant) + 1er ton (haut plat). \emph{Piège :} « yuē » se lit \cle{[ɥɛ]} (u + è fermé), pas « you-è ». Le 2e ton « jié » monte nettement depuis le milieu.
    \end{enumerate}
\end{propriete}

\begin{exemplebox}[Modèles de prononciation en contexte]
    \begin{itemize}
        \item \zh{我们要保护环境。} Wǒmen yào bǎohù huánjìng. (Nous devons protéger l'environnement.)
        \item \zh{请不要乱扔垃圾。} Qǐng bùyào luàn rēng lājī. (S'il vous plaît, ne jetez pas les déchets n'importe où.)
        \item \zh{大家都要节约用水。} Dàjiā dōu yào jiéyuē yòng shuǐ. (Tout le monde doit économiser l'eau.)
    \end{itemize}
\end{exemplebox}

\courssub{Focus sur les caractères : radicaux et ordre des traits}

Comprendre la structure des caractères (clé phonétique / clé sémantique) aide à la mémorisation et à la reconnaissance. Voici 4 caractères stratégiques de la leçon.

\begin{definition}[Analyse graphémique de 4 caractères clés]
    \begin{itemize}
        \item \textbf{环 huán} (environnement) : Clé \zh{王} wáng (jade/roi) en bas → lien avec quelque chose de précieux/entouré. Haut : \zh{瞢} méng (phonétique). Ordre : haut (艹, 目, 皿) → bas (王). \emph{17 traits.}
        \item \textbf{污 wū} (pollution/sale) : Clé \zh{氵} shuǐ (eau, 3 traits) à gauche → sémantique (liquide sale). Droite : \zh{于} yú (phonétique). Ordre : clé eau (点, 点, 提) → droite (一, 丨, 乚). \emph{6 traits.}
        \item \textbf{种 zhòng/zhǒng} (planter / graine/espèce) : Clé \zh{禾} hé (riz/plant) en haut → sémantique (agriculture). Bas : \zh{中} zhōng (phonétique + sens « au milieu » = graine dans la terre). Ordre : haut (禾 : 撇, 横, 竖, 撇, 捺) → bas (中 : 竖, 横折, 横). \emph{9 traits.}
        \item \textbf{责 zé} (responsabilité) : Clé \zh{贝} bèi (coquillage/monnaie) en bas → sémantique (valeur, dette morale). Haut : \zh{朁} (phonétique). Ordre : haut → bas (贝 : 横, 竖, 横, 竖, 横折, 横, 横). \emph{11 traits.}
    \end{itemize}
\end{definition}

\begin{methode}[Comment mémoriser l'ordre des traits ?]
    \begin{enumerate}
        \item \textbf{Identifiez la clé (radical)} : elle donne l'indice de sens (eau, plante, jade, coquillage).
        \item \textbf{Appliquez les règles universelles} : Haut → Bas (环, 种) ; Gauche → Droite (污) ; Extérieur → Intérieur → Fermeture.
        \item \textbf{Écrivez 5 fois} en dictant l'ordre à voix haute (ex: « pie, heng, shu, pie, na... »).
        \item \textbf{Associez à un mot connu} : \zh{污} dans \zh{污染} / \zh{脏污} ; \zh{种} dans \zh{种树} / \zh{种子}.
    \end{enumerate}
\end{methode}

\courssub{Collocations usuelles et structures types}

Le vocabulaire ne vit pas isolé. Voici les \textbf{collocations (associations mots-mots) incontournables} pour réussir l'expression écrite et orale.

\begin{center}
\begin{tabularx}{\linewidth}{|X|X|X|}
\hline
\rowcolor{popblue} \textbf{Structure / Collocation} & \textbf{Exemple (Caractères + Pinyin)} & \textbf{Traduction} \\ \hline
\zh{保护 + 环境/地球/树/水} & \zh{保护地球} bǎohù dìqiú & Protéger la Terre \\ \hline
\zh{污染 + 空气/水/环境} (v. ou n.) & \zh{严重污染了空气} yánzhòng wūrǎn le kōngqì & A gravement pollué l'air \\ \hline
\zh{乱扔 + 垃圾} & \zh{不乱扔垃圾} bù luàn rēng lājī & Ne pas jeter les ordures n'importe où \\ \hline
\zh{节约 + 用水/用电/纸/粮食} & \zh{节约用电} jiéyuē yòng diàn & Économiser l'électricité \\ \hline
\zh{种 + 树/花/草} & \zh{种树} zhòng shù & Planter des arbres \\ \hline
\zh{回收 + 垃圾/塑料/纸/电池} & \zh{回收塑料瓶} huíshōu sùliào píng & Recycler les bouteilles en plastique \\ \hline
\zh{越来越 + Adj./Verbe d'état} & \zh{环境越来越好} huánjìng yuèláiyuè hǎo & L'environnement est de mieux en mieux \\ \hline
\zh{应该 + V} (obligation morale) & \zh{我们应该负责} wǒmen yīnggāi fùzé & Nous devons être responsables \\ \hline
\zh{是...的责任} & \zh{保护环境是大家的责任} & Protéger l'env. est la responsabilité de tous \\ \hline
\end{tabularx}
\end{center}

\courssub{Boîtes pédagogiques : Pièges, Astuces et Culture}

\begin{attention}[Le piège de la polyphonie : \zh{种} zhòng vs zhǒng]
    C'est \textbf{l'erreur n°1} sur ce caractère.
    \begin{itemize}
        \item \textbf{zhòng (4e ton)} = \textbf{Verbe} « planter, semer, cultiver ». Ex: \zh{种树} zhòng shù, \zh{种菜} zhòng cài.
        \item \textbf{zhǒng (3e ton)} = \textbf{Nom} « graine, semence, espèce, type, sorte ». Ex: \zh{种子} zhǒngzi (graine), \zh{树种} shùzhǒng (essence d'arbre), \zh{这种} zhè zhǒng (ce genre de).
    \end{itemize}
    \textbf{Astuce mnémotechnique} : Pour \textbf{planter} (action dynamique, energy), le ton \textbf{chute} (4e). Pour la \textbf{graine} (petite, posée, statique), le ton \textbf{descend puis remonte} (3e, comme une graine qu'on enterre et qui va germer).
\end{attention}

\begin{savaistu}[\zh{垃圾 lājī} : un mot « inséparable »]
    Contrairement à la plupart des mots chinois bisyllabiques dont les caractères existent seuls (\zh{环}境, \zh{保}护, \zh{污}染), \zh{垃} lā et \zh{圾} jī \textbf{ne s'emploient quasi jamais isolément} en chinois moderne.
    \begin{itemize}
        \item \zh{垃} : n'existe plus seul (anciennement « déchets de céréales »).
        \item \zh{圾} : n'existe pas seul.
    \end{itemize}
    C'est un \emph{mot lié} (合成词 héchéngcí) : les deux morphèmes ont fusionné pour ne former qu'une seule unité lexicale signifiant « ordures ». C'est pourquoi on ne dit pas « un \zh{垃} » ni « un \zh{圾} ».
\end{savaistu}

\begin{methode}[Stratégie de mémorisation : Regroupement par familles de sens]
    Ne mémorisez pas la liste dans l'ordre alphabétique du pinyin. Construisez une \textbf{carte mentale} à 3 branches :
    \begin{enumerate}
        \item \textbf{LES PROBLÈMES} (mots « négatifs ») : \zh{污染, 垃圾, 乱扔, 脏}.
        \item \textbf{LES SOLUTIONS} (verbes d'action positive) : \zh{保护, 节约, 种, 回收, 扔} (dans le sens « jeter à la poubelle » → \zh{扔进垃圾桶}).
        \item \textbf{LE CADRE ET LE DEVOIR} (noms + modaux) : \zh{环境, 空气, 地球, 树, 塑料, 应该, 责任, 干净, 越来越}.
    \end{enumerate}
    \textbf{Exercice} : Dessinez cette carte mentale sur une feuille A4 paysage. Utilisez des flèches : Problèmes $\xrightarrow{\text{nécessitent}}$ Solutions ; Cadre $\xrightarrow{\text{impose}}$ Devoir.
\end{methode}

\courssub{Exercice résolu : Association mot-image et traduction}

\begin{exoresolu}[Entraînement rapide : Vocabulaire $\leftrightarrow$ Sens]
\textbf{Consigne :} Associez chaque mot chinois (1--8) à sa définition française (A--H) et donnez le pinyin avec tons.

\begin{enumerate}
    \item \zh{回收} \hfill A. Devoir, obligation morale
    \item \zh{越来越} \hfill B. Planter (verbe)
    \item \zh{责任} \hfill C. Recycler
    \item \zh{种} \hfill D. Déchets, ordures
    \item \zh{应该} \hfill E. De plus en plus
    \item \zh{垃圾} \hfill F. Responsabilité
    \item \zh{干净} \hfill G. Propre
    \item \zh{节约} \hfill H. Économiser
\end{enumerate}

\tcblower
\textbf{Corrigé détaillé :}
\begin{enumerate}
    \item \zh{回收} \textbf{huíshōu} — \textbf{C} (Recycler). Composé de \zh{回} « retour » + \zh{收} « collecter ».
    \item \zh{越来越} \textbf{yuèláiyuè} — \textbf{E} (De plus en plus). Structure figée.
    \item \zh{责任} \textbf{zérèn} — \textbf{F} (Responsabilité). \zh{责} (devoir) + \zh{任} (charge/mission).
    \item \zh{种} \textbf{zhòng} — \textbf{B} (Planter). \textbf{Verbe = 4e ton}. Ne pas confondre avec \zh{zhǒng} (graine).
    \item \zh{应该} \textbf{yīnggāi} — \textbf{A} (Devoir). Verbe modal, pas de « à » après (pas de *应该 \textbf{到} 去).
    \item \zh{垃圾} \textbf{lājī} — \textbf{D} (Déchets). Mot inséparable, 1+1 ton.
    \item \zh{干净} \textbf{gānjìng} — \textbf{G} (Propre). \zh{干} sec/net + \zh{净} pur. Antonyme : \zh{脏} zāng.
    \item \zh{节约} \textbf{jiéyuē} — \textbf{H} (Économiser). 2+1 ton. \zh{节} festival/économiser + \zh{约} 约束.
\end{enumerate}
\end{exoresolu}

\courssub{Complément « Spécial Examen » : \zh{回收} et \zh{塑料}}

Les sujets de compréhension (écrite/orale) du Baccalauréat et du HSK 3/4 intègrent très fréquemment le tri des déchets et le plastique. Maîtrisez ce duo :

\begin{definition}[Duo « Examen » : \zh{回收 huíshōu} + \zh{塑料 sùliào}]
    \begin{itemize}
        \item \textbf{回收 huíshōu} (v.) : Recycler, récupérer pour réutiliser. \zh{回} « retour » + \zh{收} « collecter ».
            \begin{itemize}
                \item \zh{可回收垃圾} kě huíshōu lājī — déchets recyclables (poubelle bleue en Chine).
                \item \zh{回收站} huíshōu zhàn — centre de collecte / déchetterie.
            \end{itemize}
        \item \textbf{塑料 sùliào} (n.) : Plastique. \zh{塑} « modeler » + \zh{料} « matière ».
            \begin{itemize}
                \item \zh{塑料袋} sùliào dài — sac en plastique.
                \item \zh{塑料瓶} sùliào píng — bouteille en plastique.
                \item \zh{减少使用塑料} jiǎnshǎo shǐyòng sùliào — réduire l'utilisation du plastique.
            \end{itemize}
    \end{itemize}
    \textbf{Phrase type examen} : \zh{我们应该把塑料瓶扔进可回收垃圾桶。} Wǒmen yīnggāi bǎ sùliào píng rēng jìn kě huíshōu lājītǒng. (Nous devons jeter les bouteilles en plastique dans la poubelle recyclable.)
    \begin{itemize}
        \item Structure \zh{把 bǎ} (objet déplacé avant le verbe) : \zh{把 + Objet + 动词 + 结果/方向}.
        \item \zh{扔进} rēng jìn — jeter \textbf{dans} (verbe + complément de direction).
    \end{itemize}
\end{definition}

\courssub{Exercice d'application : Complétion de phrases}

\begin{exercice}[Mise en œuvre du glossaire]
\textbf{Complétez les phrases avec le mot approprié du glossaire (caractères + pinyin).}
\begin{enumerate}
    \item 我们大家都 \_\_\_\_\_\_\_\_ \zh{保护环境}。 (obligation morale)
    \item 请把废电池放进 \_\_\_\_\_\_\_\_ 箱。 (recyclage)
    \item 这条河很 \_\_\_\_\_\_\_\_，鱼都死了。 (pollution)
    \item 春天来了，我们去公园 \_\_\_\_\_\_\_\_ 树吧。 (action positive, verbe 4e ton)
    \item 现在的空气 \_\_\_\_\_\_\_\_ \_\_\_\_\_\_\_\_ 好多了。 (évolution positive)
\end{enumerate}
\end{exercice}

\begin{corrige}[Exercice d'application]
\begin{enumerate}
    \item \zh{应该} \textbf{yīnggāi} (Nous devons tous protéger l'environnement.)
    \item \zh{回收} \textbf{huíshōu} (Mettez les piles usagées dans la boîte de \textbf{recyclage}.)
    \item \zh{污染} \textbf{wūrǎn} (Cette rivière est très \textbf{polluée}, les poissons sont morts.) — \emph{Note :} 很 \zh{污染} = très pollué (verbe d'état).
    \item \zh{种} \textbf{zhòng} (Le printemps est venu, allons \textbf{planter} des arbres au parc.) — \emph{Rappel :} Verbe = 4e ton.
    \item \zh{越来越} \textbf{yuèláiyuè} (L'air est \textbf{de plus en plus} bon.) — \zh{好多了} hǎo duō le (bien mieux).
\end{enumerate}
\end{corrige}

\courssub{Texte d'entraînement intégré : « Une journée verte à Yaoundé »}

Pour consolider le glossaire dans un contexte camerounais authentique, lisez ce texte original (niveau HSK 2/3, $\approx$ 220 caractères). Le glossaire des mots nouveaux (hors liste principale) est donné en bas.

\begin{textebox}[Texte original : 一天绿色的雅温得]
    今天是六月五日，世界环境日。\\
    雅温得的天空很蓝，空气很干净。\\
    早上，我和同学一起去公园种树。\\
    我们种了五棵小树，很开心。\\
    \\
    下午，我们在学校做环保宣传。\\
    我们告诉同学们：不乱扔垃圾，要分类。\\
    塑料瓶、纸、电池要放进回收箱。\\
    我们还要节约用水，节约用电。\\
    \\
    保护地球是我们的责任。\\
    如果人人都行动，环境会越来越好。\\
    让我们一起行动吧！
\end{textebox}

\begin{center}
\begin{tabularx}{\linewidth}{|X|X|X|X|}
\hline
\rowcolor{popblueL} \textbf{Caractères} & \textbf{Pinyin} & \textbf{Nature} & \textbf{Traduction} \\ \hline
\zh{世界环境日} & shìjiè huánjìng rì & n. & Journée mondiale de l'environnement (5 juin) \\ \hline
\zh{雅温得} & Yǎwēndé & n. & Yaoundé (capitale du Cameroun) \\ \hline
\zh{分类} & fēnlèi & v. & Trier (les déchets) / classer \\ \hline
\zh{电池} & diànchí & n. & Pile, batterie \\ \hline
\zh{回收箱} & huíshōu xiāng & n. & Poubelle de recyclage / bac de récupération \\ \hline
\zh{如果} & rúguǒ & conj. & Si / dans le cas où \\ \hline
\zh{人人} & rénrén & pron. & Tout le monde / chacun \\ \hline
\zh{行动} & xíngdòng & v./n. & Agir / action \\ \hline
\zh{一起} & yīqǐ & adv. & Ensemble \\ \hline
\end{tabularx}
\end{center}

\courssub{Questions de compréhension (entraînement examen)}

Répondez en chinois (phrase complète si possible) ou en français.

\begin{exercice}[Compréhension du texte « Une journée verte »]
\begin{enumerate}
    \item 今天是什么日子？(Quelle date / jour spécial est-ce aujourd'hui ?)
    \item 早上他们做了什么？(Qu'ont-ils fait le matin ?)
    \item 下午他们在学校做什么？(Qu'ont-ils fait l'après-midi à l'école ?)
    \item 塑料瓶和电池应该放在哪里？(Où doit-on mettre les bouteilles en plastique et les piles ?)
    \item 根据课文，保护地球是谁的责任？(D'après le texte, de qui est la responsabilité de protéger la Terre ?)
\end{enumerate}
\end{exercice}

\begin{corrige}[Compréhension du texte]
\begin{enumerate}
    \item 今天是世界环境日。(Aujourd'hui c'est la Journée mondiale de l'environnement.)
    \item 他们去公园种树，种了五棵小树。(Ils sont allés au parc planter des arbres, ils en ont planté cinq.)
    \item 他们做环保宣传，告诉同学不乱扔垃圾，要分类，回收塑料瓶、纸、电池，节约用水用电。(Ils ont fait de la sensibilisation écolo : ne pas jeter n'importe où, trier, recycler bouteilles/papier/piles, économiser eau/électricité.)
    \item 应该放进回收箱。(Il faut les mettre dans le bac de recyclage.)
    \item 是我们大家的责任 / 是人人的责任。(C'est la responsabilité de nous tous / de chacun.)
\end{enumerate}
\end{corrige}

\courssub{Synthèse et passage à la section 4}

Vous maîtrisez maintenant le \textbf{lexique de base}, sa \textbf{prononciation}, ses \textbf{collocations}, l'\textbf{analyse graphémique} de 4 caractères clés et vous avez testé votre compréhension sur un texte contextualisé au Cameroun. La section 4 (« Radicaux et ordre des traits ») approfondira l'écriture pas à pas de 6 caractères nouveaux du thème (\zh{环, 污, 种, 责, 回, 塑}) avec grilles de tracé et exercices de calligraphie. Préparez votre cahier d'écriture (田字格) !

\coursec{Radicaux et ordre des traits}

Dans la section précédente, nous avons découvert le glossaire du thème « protection de l'environnement » : \zh{环境} (huánjìng, environnement), \zh{保护} (bǎohù, protéger), \zh{污染} (wūrǎn, pollution). Vous savez maintenant les lire, les prononcer et les traduire. Mais savez-vous les \emph{écrire} ? En chinois, écrire un caractère n'est pas un dessin libre : c'est une construction réglée, avec des pièces (les clés) et un ordre de montage (l'ordre des traits). C'est ce que nous allons apprendre ici, en prenant pour modèles cinq caractères de notre thème : \zh{环}, \zh{境}, \zh{保}, \zh{护}, \zh{污}.

\courssub{Observer avant d'écrire : les caractères se décomposent}

Regardons d'abord ces caractères de près. Chacun est formé de deux blocs : un bloc à \cle{gauche}, un bloc à \cle{droite}.

\begin{textebox}[Observation : cinq caractères, cinq constructions]
环 = 王 + 不 \quad (huán, anneau, entourer)\\
境 = 土 + 竟 \quad (jìng, territoire, milieu)\\
保 = 亻 + 呆 \quad (bǎo, protéger)\\
护 = 扌 + 户 \quad (hù, défendre, garder)\\
污 = 氵 + 亏 \quad (wū, sale, polluer)
\end{textebox}

On remarque tout de suite une régularité : la partie de gauche est souvent une petite forme simple (\zh{王}, \zh{土}, \zh{亻}, \zh{扌}, \zh{氵}) qui donne une indication de \emph{sens}, tandis que la partie de droite donne souvent une indication de \emph{prononciation}. Cette petite forme de classement s'appelle la \cle{clé} (ou radical).

\begin{definition}[La clé (radical)]
Une \cle{clé} (en chinois \zh{部首}, bùshǒu) est un élément graphique qui sert à \textbf{classer le caractère dans le dictionnaire} et qui \textbf{suggère souvent son sens}. Ainsi, un caractère contenant la clé \zh{氵} (eau) désigne presque toujours quelque chose en rapport avec l'eau ; un caractère contenant la clé \zh{木} (arbre) renvoie au monde végétal ou au bois. La clé est donc à la fois un \emph{outil de recherche} (pour trouver le caractère) et un \emph{indice de sens} (pour le deviner).
\end{definition}

\begin{exemplebox}[Le radical annonce la famille de sens]
\zh{树} \emph{shù} « arbre » contient la clé \zh{木} (arbre, bois) : le sens est directement visible.\\
\zh{河} \emph{hé} « rivière » contient la clé \zh{氵} (eau) : une rivière, c'est de l'eau.\\
\zh{污} \emph{wū} « sale, polluer » contient aussi \zh{氵} : la pollution, dans l'esprit chinois, commence par l'eau souillée.\\
\zh{境} \emph{jìng} « territoire, environnement » contient \zh{土} (terre) : un environnement, c'est d'abord une terre, un sol.
\end{exemplebox}

\begin{savaistu}[Terre, eau, arbre : les trois clés de l'environnement]
Ce n'est pas un hasard si le vocabulaire de l'écologie chinoise porte souvent les clés \zh{土} (terre), \zh{氵} (eau) et \zh{木} (arbre) : \zh{垃圾} (lājī, ordures) porte la clé de la terre, \zh{河} (hé, rivière) et \zh{海} (hǎi, mer) portent la clé de l'eau, \zh{树} (shù, arbre) et \zh{林} (lín, forêt) portent la clé du bois. Au Cameroun comme en Chine, protéger l'environnement, c'est protéger ces trois éléments : le sol, l'eau et la forêt — pensez à la forêt du bassin du Congo, deuxième poumon vert de la planète.
\end{savaistu}

\courssub{La règle générale de l'ordre des traits}

\begin{propriete}[Les trois lois d'écriture]
Tout caractère chinois s'écrit selon trois principes, appliqués dans cet ordre de priorité :
\begin{enumerate}
\item \cle{De haut en bas} : on écrit d'abord les éléments du haut, puis ceux du bas. Exemple : \zh{三} (sān, trois) se trace du trait supérieur vers le trait inférieur.
\item \cle{De gauche à droite} : dans un caractère composé de deux blocs latéraux, on écrit d'abord la partie gauche, puis la partie droite. C'est le cas de nos cinq caractères du jour.
\item \cle{L'extérieur avant l'intérieur} : dans un caractère à enclos (comme \zh{国} guó, pays), on trace d'abord le cadre extérieur (sans le fermer), puis le contenu, puis on ferme.
\end{enumerate}
À l'intérieur de chaque bloc, on ajoute : l'horizontal avant le vertical (\zh{十}), et les traits se terminent parfois par un point ou un crochet.
\end{propriete}

Comparaison avec le français : en français, nous écrivons des lettres de gauche à droite sur une ligne, et l'ordre des gestes importe peu (on peut dotter le « i » à la fin). En chinois, l'ordre des traits est \textbf{obligatoire} : il garantit des caractères équilibrés, lisibles, et il est indispensable pour utiliser les dictionnaires et les claviers à traits. Un caractère écrit « dans le désordre » est immédiatement repéré par un lecteur chinois — c'est comme écrire un mot français en commençant par la dernière lettre.

\courssub{Étude détaillée des cinq caractères du thème}

\textbf{1. 环 (huán) — anneau, entourer — 8 traits}\par
Clé : \zh{王} à gauche. Attention : cette clé est en réalité la forme latérale de \zh{玉} (yù, jade) ; on l'appelle « clé du jade ». Les caractères qui la portent désignent souvent des objets précieux en jade : un anneau (\zh{环}) était autrefois un bijou de jade. Ordre des traits : on écrit d'abord la clé \zh{王} en 4 traits (horizontal, horizontal, vertical, horizontal — le dernier trait remonte légèrement), puis la partie droite \zh{不} en 4 traits (horizontal, puis le trait descendant à gauche, le vertical, le point).

\textbf{2. 境 (jìng) — territoire, milieu — 14 traits}\par
Clé : \zh{土} (terre) à gauche, en 3 traits (horizontal court, vertical, horizontal qui remonte). Puis la partie droite \zh{竟} en 11 traits, de haut en bas : d'abord \zh{音} (le son), puis \zh{儿} en bas. Règle appliquée : gauche avant droite, puis haut avant bas à l'intérieur du bloc droit. C'est le caractère le plus complexe de notre liste : prenez le temps de le découper mentalement en \zh{土} + \zh{音} + \zh{儿}.

\textbf{3. 保 (bǎo) — protéger — 9 traits}\par
Clé : \zh{亻} à gauche, la forme debout de \zh{人} (rén, homme, personne) : 2 traits (le trait oblique descendant à gauche, puis le vertical). Puis la partie droite \zh{呆} en 7 traits : d'abord \zh{口} (la bouche, 3 traits, enclos : on trace le vertical gauche, le coin supérieur, on écrit rien à l'intérieur, on ferme en bas), puis \zh{木} (4 traits). Image à retenir : une \emph{personne} (\zh{亻}) qui veille — protéger, c'est l'action d'un homme.

\textbf{4. 护 (hù) — défendre, garder — 7 traits}\par
Clé : \zh{扌} à gauche, la forme latérale de \zh{手} (shǒu, main) : 3 traits (horizontal, vertical avec crochet, trait remontant). Puis \zh{户} (porte, maison) à droite en 4 traits. Image : la \emph{main} (\zh{扌}) devant la \emph{porte} (\zh{户}) — garder la maison, c'est protéger.

\textbf{5. 污 (wū) — sale, polluer — 6 traits}\par
Clé : \zh{氵} à gauche, les « trois gouttes d'eau », forme abrégée de \zh{水} (shuǐ, eau) : 3 traits (deux points obliques en haut, puis un trait remontant en bas). Puis \zh{亏} à droite en 3 traits. C'est le caractère le plus simple de la liste : idéal pour commencer l'entraînement.

\begin{methode}[Comment écrire correctement un caractère composé]
\begin{enumerate}
\item \textbf{Identifier la clé} : repérer la partie qui classe le caractère (souvent à gauche) et se demander quel sens elle annonce.
\item \textbf{Découper le caractère} en blocs (gauche/droite, ou haut/bas, ou extérieur/intérieur).
\item \textbf{Écrire la partie gauche} (ou le haut, ou l'extérieur) en entier, en respectant l'ordre interne de ses traits.
\item \textbf{Écrire la partie droite} (ou le bas, ou l'intérieur), trait par trait, de haut en bas.
\item \textbf{Vérifier l'équilibre} : la clé à gauche est étroite (environ un tiers de la largeur), la partie droite occupe le reste. Un caractère doit tenir dans un carré imaginaire.
\item \textbf{Réécrire} le caractère cinq fois en prononçant le pinyin à voix haute : la mémoire du geste et la mémoire du son s'ancrent ensemble.
\end{enumerate}
\end{methode}

\begin{center}
\begin{tabularx}{\linewidth}{|X|X|X|X|}
\hline
\rowcolor{popblueL} Caractère & Clé (sens) & Nombre de traits & Décomposition et ordre \\
\hline
\zh{环} huán & \zh{王} (jade) & 8 & \zh{王} (4 traits) puis \zh{不} (4 traits) \\
\hline
\zh{境} jìng & \zh{土} (terre) & 14 & \zh{土} (3 traits) puis \zh{竟} (11 traits, haut $\rightarrow$ bas) \\
\hline
\zh{保} bǎo & \zh{亻} (homme) & 9 & \zh{亻} (2 traits) puis \zh{呆} (7 traits) \\
\hline
\zh{护} hù & \zh{扌} (main) & 7 & \zh{扌} (3 traits) puis \zh{户} (4 traits) \\
\hline
\zh{污} wū & \zh{氵} (eau) & 6 & \zh{氵} (3 traits) puis \zh{亏} (3 traits) \\
\hline
\end{tabularx}
\end{center}

\begin{popfigure}
\begin{tikzpicture}[node distance=0.5cm, every node/.style={font=\normalsize}]
\node[draw=popblue, fill=popblueL, rounded corners, minimum width=1.1cm, minimum height=1.1cm, font=\Large] (bao) {\zh{保}};
\node[right=0.4cm of bao, font=\Large] (eq1) {=};
\node[draw=popgreen, fill=popgreenL, rounded corners, minimum width=1.1cm, minimum height=1.1cm, right=0.4cm of eq1, font=\Large] (ren) {\zh{亻}};
\node[right=0.3cm of ren, font=\Large] (plus1) {+};
\node[draw=poporange, fill=poporangeL, rounded corners, minimum width=1.1cm, minimum height=1.1cm, right=0.3cm of plus1, font=\Large] (dai) {\zh{呆}};
\node[below=0.25cm of ren, text=popdark, font=\small] {homme (clé)};
\node[below=0.25cm of dai, text=popdark, font=\small] {partie phonétique};

\node[draw=popblue, fill=popblueL, rounded corners, minimum width=1.1cm, minimum height=1.1cm, font=\Large, below=1.5cm of bao] (wu) {\zh{污}};
\node[right=0.4cm of wu, font=\Large] (eq2) {=};
\node[draw=popgreen, fill=popgreenL, rounded corners, minimum width=1.1cm, minimum height=1.1cm, right=0.4cm of eq2, font=\Large] (shui) {\zh{氵}};
\node[right=0.3cm of shui, font=\Large] (plus2) {+};
\node[draw=poporange, fill=poporangeL, rounded corners, minimum width=1.1cm, minimum height=1.1cm, right=0.3cm of plus2, font=\Large] (kui) {\zh{亏}};
\node[below=0.25cm of shui, text=popdark, font=\small] {eau (clé)};
\node[below=0.25cm of kui, text=popdark, font=\small] {partie phonétique};

\node[draw=popblue, fill=popblueL, rounded corners, minimum width=1.1cm, minimum height=1.1cm, font=\Large, below=1.5cm of wu] (jing) {\zh{境}};
\node[right=0.4cm of jing, font=\Large] (eq3) {=};
\node[draw=popgreen, fill=popgreenL, rounded corners, minimum width=1.1cm, minimum height=1.1cm, right=0.4cm of eq3, font=\Large] (tu) {\zh{土}};
\node[right=0.3cm of tu, font=\Large] (plus3) {+};
\node[draw=poporange, fill=poporangeL, rounded corners, minimum width=1.1cm, minimum height=1.1cm, right=0.3cm of plus3, font=\Large] (jingr) {\zh{竟}};
\node[below=0.25cm of tu, text=popdark, font=\small] {terre (clé)};
\node[below=0.25cm of jingr, text=popdark, font=\small] {partie phonétique};
\end{tikzpicture}
\legende{Décomposition des caractères : à gauche la clé (sens), à droite la partie phonétique. On écrit toujours la clé d'abord.}
\end{popfigure}

\courssub{Pièges fréquents des francophones}

\begin{attention}[La clé 扌 (main) ne se confond pas avec 木 (arbre)]
Erreur très fréquente : \zh{扌} (main, \textbf{3 traits}) et \zh{木} (arbre, \textbf{4 traits}) se ressemblent, mais ils diffèrent par deux points. D'abord, \zh{木} possède deux traits obliques en bas (les « racines »), que \zh{扌} n'a pas ; ensuite, le dernier trait de \zh{扌} est un trait fin qui \emph{remonte} vers la droite, tandis que le vertical de \zh{木} est traversé par les obliques. Conséquence : \zh{护} (hù, protéger, avec la main) et un hypothétique caractère avec \zh{木} n'ont rien à voir. Astuce : la main \emph{cueille} vers le haut (trait remontant), l'arbre \emph{s'enracine} vers le bas (deux obliques). De même, ne confondez pas \zh{王} (jade, dans \zh{环}) avec \zh{土} (terre, dans \zh{境}) : dans \zh{土}, le trait vertical \emph{dépasse} au-dessus du premier horizontal ; dans \zh{王}, il ne dépasse pas.
\end{attention}

\begin{attention}[Ne jamais écrire la partie droite avant la clé]
Beaucoup de débutants francophones, habitués à écrire les lettres dans l'ordre qui les arrange, tracent d'abord la partie droite « parce qu'elle est plus grande ». C'est une faute : la clé à gauche se trace \textbf{toujours en premier}. Elle fixe la hauteur et la position du caractère ; la partie droite vient ensuite s'aligner sur elle.
\end{attention}

\begin{exoresolu}[Décomposer et décrire l'ordre d'écriture]
Énoncé : pour chaque caractère, indiquez la clé, le nombre total de traits et l'ordre d'écriture des blocs : a) \zh{护} ; b) \zh{河} ; c) \zh{境}.
\tcblower
Solution détaillée :
\begin{itemize}
\item a) \zh{护} (hù, protéger) : clé \zh{扌} (main) à gauche ; 7 traits au total. On écrit d'abord \zh{扌} (3 traits : horizontal, vertical-crochet, trait remontant), puis \zh{户} (4 traits) à droite. Règle : gauche avant droite.
\item b) \zh{河} (hé, rivière) : clé \zh{氵} (eau) à gauche ; 8 traits au total. On trace d'abord les trois gouttes d'eau (3 traits), puis \zh{可} (5 traits) à droite. Le sens est annoncé par la clé : une rivière est faite d'eau.
\item c) \zh{境} (jìng, environnement) : clé \zh{土} (terre) à gauche ; 14 traits au total. On écrit \zh{土} (3 traits), puis \zh{竟} (11 traits) en procédant de haut en bas : \zh{音} d'abord, \zh{儿} ensuite. Règles combinées : gauche avant droite, puis haut avant bas.
\end{itemize}
\end{exoresolu}

\begin{exercice}[À vous d'écrire]
1. Décomposez les caractères suivants en clé + partie droite, et donnez le sens suggéré par la clé : \zh{环}, \zh{污}, \zh{树}, \zh{保}.\\
2. Classez ces caractères selon leur clé (\zh{土}, \zh{氵}, \zh{木}, \zh{亻}, \zh{扌}) : \zh{境}, \zh{河}, \zh{树}, \zh{保}, \zh{护}, \zh{污}, \zh{海}, \zh{林}.\\
3. Recopiez cinq fois \zh{环境} et \zh{保护} en respectant l'ordre des traits, puis écrivez la phrase \zh{我们保护环境。} (Wǒmen bǎohù huánjìng. — Nous protégeons l'environnement.)
\end{exercice}

\begin{corrige}[Exercice : décomposition et classement]
1. \zh{环} = \zh{王} (jade) + \zh{不} ; \zh{污} = \zh{氵} (eau) + \zh{亏} ; \zh{树} = \zh{木} (arbre) + \zh{对} ; \zh{保} = \zh{亻} (homme) + \zh{呆}.\\
2. Clé \zh{土} : \zh{境} ; clé \zh{氵} : \zh{河}, \zh{污}, \zh{海} ; clé \zh{木} : \zh{树}, \zh{林} ; clé \zh{亻} : \zh{保} ; clé \zh{扌} : \zh{护}.\\
3. Vérification : dans \zh{环} et \zh{境}, la clé de gauche est écrite en premier ; dans \zh{保} et \zh{护}, idem. La phrase complète compte 6 caractères, chacun écrit de gauche à droite dans ses blocs internes.
\end{corrige}

\begin{aretenir}[L'essentiel de la section]
\begin{itemize}
\item La \cle{clé} (radical) classe le caractère dans le dictionnaire et suggère son sens : \zh{土} terre, \zh{氵} eau, \zh{木} arbre, \zh{亻} homme, \zh{扌} main, \zh{王} jade.
\item Ordre des traits : \textbf{haut $\rightarrow$ bas}, \textbf{gauche $\rightarrow$ droite}, \textbf{extérieur $\rightarrow$ intérieur} ; dans un caractère latéral, la clé de gauche s'écrit toujours en premier.
\item Nos cinq modèles : \zh{环} (8 traits), \zh{境} (14), \zh{保} (9), \zh{护} (7), \zh{污} (6).
\item Pièges : \zh{扌} (3 traits, trait remontant) contre \zh{木} (4 traits, racines) ; \zh{王} contre \zh{土} (le vertical de \zh{土} dépasse).
\end{itemize}
\end{aretenir}

\coursec{Collocations et questions de compréhension}

Dans la section précédente, nous avons appris à écrire les caractères du thème de l'environnement en respectant leurs clés (radicaux) et l'ordre des traits. Savoir écrire un caractère, c'est bien ; savoir l'employer dans une phrase, c'est encore mieux. Nous allons maintenant voir comment les mots du thème s'assemblent naturellement entre eux — ce sont les \cle{collocations} — puis nous apprendrons à répondre correctement, en chinois, aux questions de compréhension d'un texte, une compétence essentielle à l'examen.

\courssub{Les collocations essentielles du thème}

\begin{definition}[Qu'est-ce qu'une collocation ?]
Une \cle{collocation} est une combinaison habituelle de mots qui vont naturellement ensemble. En français, on dit « \emph{faire} attention » et non « \emph{prendre} attention » ; en chinois, on dit \zh{保护环境} (protéger l'environnement) et non \zh{保护情况}. Apprendre les mots par collocations, c'est apprendre le chinois tel que les Chinois le parlent réellement.
\end{definition}

Observons d'abord ces six collocations tirées de notre texte support :

\begin{center}
\begin{tabularx}{\linewidth}{|X|X|X|X|}
\hline
\rowcolor{popblueL} Caractères & Pinyin & Structure & Traduction \\
\hline
保护环境 & bǎohù huánjìng & verbe + nom & protéger l'environnement \\
\hline
节约用水 & jiéyuē yòng shuǐ & verbe + verbe + nom & économiser l'eau (litt. « économiser l'usage de l'eau ») \\
\hline
节约用电 & jiéyuē yòng diàn & verbe + verbe + nom & économiser l'électricité \\
\hline
乱扔垃圾 & luàn rēng lājī & adverbe + verbe + nom & jeter les déchets n'importe où \\
\hline
空气污染 & kōngqì wūrǎn & nom + nom/verbe & pollution de l'air \\
\hline
多种树 & duō zhòng shù & adverbe + verbe + nom & planter plus d'arbres \\
\hline
\end{tabularx}
\end{center}

\begin{propriete}[Le rôle de 多 (duō) et 乱 (luàn) devant le verbe]
\zh{多} (duō, « beaucoup, plus ») et \zh{乱} (luàn, « en désordre, n'importe comment ») se placent \textbf{avant le verbe} qu'ils modifient, comme des adverbes :
\begin{itemize}
\item \zh{多种树} : duō zhòng shù — « planter \emph{plus} d'arbres » (多 + verbe 种)
\item \zh{乱扔垃圾} : luàn rēng lājī — « jeter les déchets \emph{n'importe où} » (乱 + verbe 扔)
\end{itemize}
Attention à l'ordre : on dit \zh{多种树} et jamais \zh{种多树}. De même, on dit \zh{乱扔} et jamais \zh{扔乱}.
\end{propriete}

\begin{exemplebox}[Chaque collocation dans une phrase complète]
\zh{我们应该保护环境。}\\
\emph{Wǒmen yīnggāi bǎohù huánjìng.}\\
« Nous devons protéger l'environnement. »

\zh{洗手以后，请节约用水。}\\
\emph{Xǐshǒu yǐhòu, qǐng jiéyuē yòng shuǐ.}\\
« Après vous être lavé les mains, économisez l'eau, s'il vous plaît. »

\zh{离开教室的时候要节约用电。}\\
\emph{Líkāi jiàoshì de shíhou yào jiéyuē yòng diàn.}\\
« En quittant la salle de classe, il faut économiser l'électricité. »

\zh{请不要乱扔垃圾！}\\
\emph{Qǐng búyào luàn rēng lājī!}\\
« Ne jetez pas vos déchets n'importe où, s'il vous plaît ! »

\zh{雅温得的空气污染越来越严重。}\\
\emph{Yǎwēndé de kōngqì wūrǎn yuèláiyuè yánzhòng.}\\
« La pollution de l'air à Yaoundé devient de plus en plus grave. »

\zh{每年春天，我们学校多种树。}\\
\emph{Měinián chūntiān, wǒmen xuéxiào duō zhòng shù.}\\
« Chaque printemps, notre école plante plus d'arbres. »
\end{exemplebox}

\begin{attention}[La structure figée 越来越 + adjectif]
\zh{越来越} (yuèláiyuè) + adjectif exprime une \textbf{progression} : « de plus en plus… ».
\begin{itemize}
\item \zh{天气越来越热。} \emph{Tiānqì yuèláiyuè rè.} « Il fait de plus en plus chaud. »
\item \zh{我们的环境越来越差了。} \emph{Wǒmen de huánjìng yuèláiyuè chà le.} « Notre environnement devient de plus en plus mauvais. »
\end{itemize}
C'est une \textbf{structure figée} : on ne dit jamais \zh{很越来越}. En effet, \zh{很} (hěn, « très ») décrit un état fixe, alors que \zh{越来越} décrit un changement : les deux sont incompatibles. On ne dit pas non plus \zh{越来越很热}. La particule finale \zh{了} marque le constat du changement : \zh{越来越差了。}
\end{attention}

\courssub{Le texte support : relecture ciblée}

Relisons le texte de la leçon pour y repérer les collocations et préparer les réponses aux questions.

\begin{textebox}[保护我们的环境 — Bǎohù wǒmen de huánjìng]
\zh{现在，我们的环境越来越差了。空气被污染了，水也被污染了。在很多大城市里，人们常常看不到蓝蓝的天。}\\
\emph{Xiànzài, wǒmen de huánjìng yuèláiyuè chà le. Kōngqì bèi wūrǎn le, shuǐ yě bèi wūrǎn le. Zài hěn duō dà chéngshì lǐ, rénmen chángcháng kànbujiào lánlán de tiān.}\\
« Aujourd'hui, notre environnement devient de plus en plus mauvais. L'air est pollué, l'eau aussi est polluée. Dans beaucoup de grandes villes, on ne voit souvent plus le ciel bleu. »

\zh{为什么会这样？因为车太多了，工厂太多了，还有人乱扔垃圾。}\\
\emph{Wèishénme huì zhèyàng? Yīnwèi chē tài duō le, gōngchǎng tài duō le, hái yǒu rén luàn rēng lājī.}\\
« Pourquoi en est-il ainsi ? Parce qu'il y a trop de voitures, trop d'usines, et que des gens jettent leurs déchets n'importe où. »

\zh{我们应该怎么做？第一，要节约用水。第二，要节约用电。第三，不要乱扔垃圾。我们还要多种树，因为树可以让空气更干净。}\\
\emph{Wǒmen yīnggāi zěnme zuò? Dì-yī, yào jiéyuē yòng shuǐ. Dì-èr, yào jiéyuē yòng diàn. Dì-sān, búyào luàn rēng lājī. Wǒmen hái yào duō zhòng shù, yīnwèi shù kěyǐ ràng kōngqì gèng gānjìng.}\\
« Que devons-nous faire ? Premièrement, économiser l'eau. Deuxièmement, économiser l'électricité. Troisièmement, ne pas jeter les déchets n'importe où. Nous devons aussi planter plus d'arbres, car les arbres rendent l'air plus propre. »

\zh{保护地球，就是保护我们的家。让我们一起努力吧！}\\
\emph{Bǎohù dìqiú, jiùshì bǎohù wǒmen de jiā. Ràng wǒmen yìqǐ nǔlì ba!}\\
« Protéger la Terre, c'est protéger notre maison. Faisons des efforts ensemble ! »
\end{textebox}

\begin{definition}[Trois points de langue du texte]
\begin{itemize}
\item \zh{被} (bèi) introduit le \textbf{passif} : Sujet (qui subit) + 被 + verbe. \zh{空气被污染了。} « L'air est pollué. »
\item \zh{看不到} (kànbujiào) : verbe + \zh{不到} exprime l'\textbf{incapacité de résultat} : « ne pas arriver à voir ». À l'affirmatif : \zh{看得到} « arriver à voir ».
\item \zh{就是} (jiùshì) : « c'est précisément, c'est-à-dire ». \zh{保护地球，就是保护我们的家。} « Protéger la Terre, c'est (précisément) protéger notre maison. »
\end{itemize}
\end{definition}

\begin{exemplebox}[Exemple modèle à retenir]
\zh{我们应该保护地球。}\\
\emph{Wǒmen yīnggāi bǎohù dìqiú.}\\
« Nous devons protéger la Terre. »\\
Structure : Sujet (\zh{我们}) + \zh{应该} (devoir) + verbe (\zh{保护}) + complément (\zh{地球}).
\end{exemplebox}

\courssub{Méthode : répondre à une question de compréhension}

\begin{methode}[Répondre en phrase complète, comme à l'examen]
\begin{enumerate}
\item \textbf{Lire la question} et souligner le mot-clé interrogatif : \zh{怎么样} (comment), \zh{什么} (quoi), \zh{怎么} (comment faire), \zh{为什么} (pourquoi).
\item \textbf{Repérer dans le texte} la phrase qui contient le même mot-clé ou le même thème.
\item \textbf{Reprendre les mots de la question} dans la réponse : si la question commence par \zh{我们的环境怎么样？}, la réponse commence par \zh{我们的环境…}.
\item \textbf{Compléter avec l'information du texte}, sans inventer.
\item \textbf{Vérifier} : phrase complète (sujet + verbe), tons corrects, ponctuation chinoise 。
\end{enumerate}
\end{methode}

\begin{popfigure}
\begin{tikzpicture}[
  node distance=0.55cm,
  box/.style={draw=popblue, fill=popblueL, rounded corners=3pt, align=center, text width=4.6cm, inner sep=4pt, font=\small},
  arr/.style={-{Stealth[length=2.5mm]}, thick, poporange}
]
\node[box, align=center] (q){\textbf{1. Question}\\ 我们的环境怎么样？};
\node[box, below=of q, align=center] (kw){\textbf{2. Mot-clé}\\ 怎么样 « comment »};
\node[box, below=of kw, align=center] (txt){\textbf{3. Phrase du texte}\\ 环境越来越差了};
\node[box, below=of txt, align=center] (rep){\textbf{4. Réponse complète}\\ 我们的环境越来越差了。};
\draw[arr] (q) -- (kw);
\draw[arr] (kw) -- (txt);
\draw[arr] (txt) -- (rep);
\end{tikzpicture}
\legende{Organigramme de la méthode : de la question à la réponse complète.}
\end{popfigure}

\courssub{Questions de compréhension : analyse et réponses modèles}

\textbf{Question 1 :} \zh{我们的环境怎么样？} \emph{Wǒmen de huánjìng zěnmeyàng?} « Comment est notre environnement ? »

Réponse modèle : \zh{我们的环境越来越差了。} \emph{Wǒmen de huánjìng yuèláiyuè chà le.} « Notre environnement devient de plus en plus mauvais. »

On reprend \zh{我们的环境} (de la question) et on ajoute l'information du texte \zh{越来越差了}.

\textbf{Question 2 :} \zh{什么被污染了？} \emph{Shénme bèi wūrǎn le?} « Qu'est-ce qui est pollué ? »

Réponse modèle : \zh{空气和水都被污染了。} \emph{Kōngqì hé shuǐ dōu bèi wūrǎn le.} « L'air et l'eau sont tous les deux pollués. »

\begin{propriete}[La place de 都 (dōu, « tous, tout »)]
\zh{都} se place \textbf{avant le verbe} (et après le sujet pluriel ou composé) :
\begin{center}
Sujet (A + 和 + B) + 都 + verbe
\end{center}
\zh{空气和水都被污染了。} « L'air et l'eau sont \emph{tous les deux} pollués. »\\
\zh{我们都喜欢汉语。} \emph{Wǒmen dōu xǐhuan Hànyǔ.} « Nous aimons tous le chinois. »\\
Erreur fréquente des francophones : placer \zh{都} avant le sujet, par calque du français « tous… ». On ne dit jamais \zh{都空气和水被污染了}.
\end{propriete}

\textbf{Question 3 :} \zh{我们应该怎么做？} \emph{Wǒmen yīnggāi zěnme zuò?} « Que devons-nous faire ? »

Réponse modèle : \zh{我们应该节约用水、节约用电，还要多种树。} \emph{Wǒmen yīnggāi jiéyuē yòng shuǐ, jiéyuē yòng diàn, hái yào duō zhòng shù.} « Nous devons économiser l'eau, économiser l'électricité, et aussi planter plus d'arbres. »

Remarquez la virgule d'énumération chinoise \zh{、} (dunhao), réservée aux listes de mots ou de groupes parallèles. Remarquez aussi \zh{还要} (hái yào) : « il faut en outre », qui ajoute une action supplémentaire.

\textbf{Question 4 (ouverte) :} \zh{在你的城市，有什么环境问题？} \emph{Zài nǐ de chéngshì, yǒu shénme huánjìng wèntí?} « Quels problèmes environnementaux y a-t-il dans ta ville ? »

Réponse possible (exemple pour Douala) : \zh{在我的城市，空气污染很严重，也有人乱扔垃圾。} \emph{Zài wǒ de chéngshì, kōngqì wūrǎn hěn yánzhòng, yě yǒu rén luàn rēng lājī.} « Dans ma ville, la pollution de l'air est grave, et il y a aussi des gens qui jettent leurs déchets n'importe où. »

\begin{center}
\begin{tabularx}{\linewidth}{|X|X|X|}
\hline
\rowcolor{popblueL} Question & Réponse modèle (caractères + pinyin) & Traduction \\
\hline
我们的环境怎么样？ & 我们的环境越来越差了。Wǒmen de huánjìng yuèláiyuè chà le. & Notre environnement devient de plus en plus mauvais. \\
\hline
什么被污染了？ & 空气和水都被污染了。Kōngqì hé shuǐ dōu bèi wūrǎn le. & L'air et l'eau sont tous les deux pollués. \\
\hline
我们应该怎么做？ & 我们应该节约用水、节约用电，还要多种树。Wǒmen yīnggāi jiéyuē yòng shuǐ, jiéyuē yòng diàn, hái yào duō zhòng shù. & Nous devons économiser l'eau et l'électricité, et planter plus d'arbres. \\
\hline
在你的城市，有什么环境问题？ & 在我的城市，空气污染很严重。Zài wǒ de chéngshì, kōngqì wūrǎn hěn yánzhòng. & Dans ma ville, la pollution de l'air est grave. \\
\hline
\end{tabularx}
\end{center}

\courssub{Entraînement au format examen}

\begin{exoresolu}[Répondre par une phrase complète]
Énoncé : répondez en chinois, par une phrase complète, à la question suivante : \zh{我们应该怎么做？} \emph{Wǒmen yīnggāi zěnme zuò?}
\tcblower
Solution détaillée :
\begin{enumerate}
\item Mot-clé de la question : \zh{怎么做} (« comment faire ») ; sujet : \zh{我们} + modal \zh{应该}.
\item Dans le texte, on lit : \zh{要节约用水。要节约用电。不要乱扔垃圾。我们还要多种树。}
\item On reprend les mots de la question et on complète :\\
\zh{我们应该节约用水、节约用电，还要多种树。}\\
\emph{Wǒmen yīnggāi jiéyuē yòng shuǐ, jiéyuē yòng diàn, hái yào duō zhòng shù.}\\
« Nous devons économiser l'eau, économiser l'électricité, et aussi planter plus d'arbres. »
\item Vérification : sujet + modal + verbe présents, énumération avec \zh{、}, point final 。 — réponse complète et correcte.
\end{enumerate}
Réponse incomplète à éviter : \zh{节约用水。} seul (sans sujet ni \zh{应该}) perd des points à l'examen, car elle ne reprend pas les mots de la question.
\end{exoresolu}

\begin{exercice}[À vous de répondre]
Répondez en chinois par des phrases complètes :
\begin{enumerate}
\item \zh{我们的环境怎么样？}
\item \zh{什么被污染了？}
\item \zh{在你的城市，有什么环境问题？} (réponse personnelle, avec \zh{有} ou \zh{很})
\end{enumerate}
\end{exercice}

\begin{corrige}[Exercice : répondre par des phrases complètes]
\begin{enumerate}
\item \zh{我们的环境越来越差了。} \emph{Wǒmen de huánjìng yuèláiyuè chà le.} « Notre environnement devient de plus en plus mauvais. »
\item \zh{空气和水都被污染了。} \emph{Kōngqì hé shuǐ dōu bèi wūrǎn le.} « L'air et l'eau sont tous les deux pollués. » (Bien placer \zh{都} avant le verbe \zh{被污染}.)
\item Réponse personnelle, par exemple : \zh{在我的城市，有很多人乱扔垃圾，水也被污染了。} \emph{Zài wǒ de chéngshì, yǒu hěn duō rén luàn rēng lājī, shuǐ yě bèi wūrǎn le.} « Dans ma ville, beaucoup de gens jettent leurs déchets n'importe où, et l'eau aussi est polluée. »
\end{enumerate}
\end{corrige}

\begin{attention}[Pièges fréquents des francophones]
\begin{itemize}
\item Ne traduisez pas mot à mot « de plus en plus » : utilisez la structure figée \zh{越来越 + adjectif}, jamais \zh{很越来越}.
\item \zh{都} se place \textbf{avant le verbe}, jamais en tête de phrase : \zh{空气和水都被污染了。}
\item À l'examen, une réponse sans sujet (\zh{越来越差了。} seul) est incomplète : reprenez toujours les mots de la question.
\item N'oubliez pas la particule \zh{了} finale avec \zh{越来越}, qui marque le constat du changement : \zh{越来越差了。}
\item Ne confondez pas \zh{差} (chà, « mauvais, médiocre ») avec \zh{坏} (huài, « mauvais, cassé ») : pour l'environnement, on dit \zh{环境差} ou \zh{环境不好}, pas \zh{环境坏}.
\end{itemize}
\end{attention}

\begin{aretenir}[L'essentiel de la section]
\begin{itemize}
\item Six collocations à connaître par cœur : \zh{保护环境}, \zh{节约用水}, \zh{节约用电}, \zh{乱扔垃圾}, \zh{空气污染}, \zh{多种树}.
\item \zh{越来越 + adjectif} = « de plus en plus » ; structure figée, sans \zh{很}, avec \zh{了} final.
\item \zh{都} (« tous ») se place avant le verbe : \zh{空气和水都被污染了。}
\item Le passif avec \zh{被} : Sujet + 被 + verbe (\zh{空气被污染了。}).
\item Méthode de réponse : reprendre les mots de la question + compléter avec l'information du texte = phrase complète.
\end{itemize}
\end{aretenir}

\coursec{Production guidée et ancrage socioculturel}

\courssub{Transition et objectif de la séance}

Après avoir travaillé le vocabulaire, les collocations et la compréhension fine du texte sur la protection de l'environnement, nous passons maintenant à l'\textbf{appropriation active}. L'objectif est double : \textbf{produire} un court paragraphe cohérent en réinvestissant le lexique et les structures vus (notamment \zh{应该} / \zh{不应该}, \zh{节约}, \zh{保护}, \zh{垃圾分类}) et \textbf{situer} les enjeux écologiques dans une perspective comparée Cameroun -- Chine, conformément à l'approche socioculturelle du programme MINESEC.

\courssub{Modèle de paragraphe guidé : observation et analyse}

Avant de vous lancer dans l'écriture personnelle, observons un modèle type. Ce texte sert de \textbf{trame} (squelette) que vous compléterez avec vos propres idées.

\begin{textebox}[Modèle de paragraphe guidé : « Protéger notre environnement »]
在我们的城市，环境问题很多。很多人乱扔垃圾，水也被污染了。我们应该节约用水，不乱扔塑料袋，还要参加植树活动。保护环境，人人有责。\\
Zài wǒmen de chéngshì, huánjìng wèntí hěn duō. Hěn duō rén luàn rēng lājī, shuǐ yě bèi wūrǎn le. Wǒmen yīnggāi jiéyuē yòng shuǐ, bù luàn rēng sùliào dài, hái yào cānjiā zhíshù huódòng. Bǎohù huánjìng, rénrén yǒu zé.\\
« Dans notre ville, les problèmes d'environnement sont nombreux. Beaucoup de gens jettent des déchets n'importe où, l'eau est aussi polluée. Nous devons économiser l'eau, ne pas jeter de sacs en plastique, et participer à des activités de plantation d'arbres. Protéger l'environnement, c'est la responsabilité de tous. »
\end{textebox}

\begin{aretenir}[Analyse structurelle du modèle]
\begin{itemize}
    \item \textbf{Présentation du contexte :} \zh{在我们的城市，环境问题很多} (Sujet + Locatif + Sujet + Adjectif). Structure simple pour poser le décor.
    \item \textbf{Constat (faits) :} \zh{很多人乱扔垃圾，水也被污染了}. Utilisation de \zh{乱扔} (jeter n'importe comment / désordonné) et de la \textbf{passive \zh{被} (bèi)} : \zh{水被污染了} (L'eau a été polluée). \textit{Rappel : la passive \zh{被} marque l'accent sur la victime de l'action.}
    \item \textbf{Propositions / Devoir (modalité) :} \zh{我们应该...，不乱扔...，还要...}. Enchaînement de trois verbes modaux/auxiliaires :
    \begin{enumerate}
        \item \zh{应该} (devoir / faut) + Verbe d'action positive (\zh{节约用水}).
        \item \zh{不} + Verbe (négation de l'action négative \zh{乱扔塑料袋}).
        \item \zh{还要} (encore faut / de plus il faut) + Verbe d'engagement collectif (\zh{参加植树活动}).
    \end{enumerate}
    \item \textbf{Conclusion proverbialesque :} \zh{保护环境，人人有责} (Sujet-Verbe-Objet, Sujet-Prédicat). Structure en 4+4 caractères, très rhétorique.
\end{itemize}
\end{aretenir}

\courssub{Méthodologie de la production écrite}

Pour réussir votre paragraphe personnel (4 à 5 phrases), suivez cette méthode en trois étapes.

\begin{methode}[Rédiger un paragraphe écologique simple et correct]
\textbf{Étape 1 : Sélectionner 3 à 4 mots-clés du glossaire (Section 3).}
Piochez dans les listes : \zh{节约用水}, \zh{垃圾分类}, \zh{少用塑料袋}, \zh{植树 / 种树}, \zh{保护环境}, \zh{污染}, \zh{回收}, \zh{环保袋}.

\textbf{Étape 2 : Construire des phrases simples selon le schéma SVO (Sujet + Verbe + Objet/Complément).}
Utilisez les structures maîtrisées :
\begin{itemize}
    \item \zh{我们应该 + [Verbe Positif]} (Nous devons...)
    \item \zh{我们不应该 + [Verbe Négatif]} / \zh{不要 + [Verbe Négatif]} (Nous ne devons pas... / Ne pas...)
    \item \zh{我可以 + [Verbe]} (Je peux... -- engagement personnel).
    \item \zh{[Lieu] + 要 + [Action]} (Au/À [Lieu], il faut [Action]).
\end{itemize}

\textbf{Étape 3 : Vérifier les points de vigilance (Check-list).}
\begin{enumerate}
    \item \textbf{Tons :} Chaque mot a-t-il son ton correct ? (ex: \zh{节约} jiéyuē 2-1, pas 1-2).
    \item \textbf{Ordre des mots :} Le complément de lieu/temps est-il \textbf{avant} le verbe ? (\zh{在学校我们...} ou \zh{我们在学校...} selon le focus). L'adverbe \zh{不} / \zh{也} / \zh{还} est-il \textbf{avant} le verbe/auxiliaire ?
    \item \textbf{Caractères :} Pas de confusion \zh{污} (sale) vs \zh{乌} (noir/corbeau) ; \zh{染} (teindre/polluer) vs \zh{然} (ainsi/feu).
    \item \textbf{Ponctuation :} Virgule chinoise \zh{，} pour énumérer, point final \zh{。}.
\end{enumerate}
\end{methode}

\courssub{Exercice guidé : Complétion de la trame (Entraînement collectif)}

Avant l'écrit individuel, complétons ensemble le texte ci-dessous au tableau. Choisissez parmi les options entre parenthèses.

\begin{exoresolu}[Complétion guidée]
\textbf{Énoncé :} Complétez le paragraphe suivant en choisissant la forme correcte.
\par
在我们的学校，环境\rule{2cm}{0.5pt}（问题 / 题目）不少。同学们经常\rule{2cm}{0.5pt}（乱扔 / 扔乱）纸屑，水龙头\rule{2cm}{0.5pt}（也 / 还）滴水。我们\rule{2cm}{0.5pt}（应该 / 可以）\rule{2cm}{0.5pt}（节约 / 节省）用水，\rule{2cm}{0.5pt}（不 / 没）应该\rule{2cm}{0.5pt}（乱扔 / 扔）垃圾，要\rule{2cm}{0.5pt}（做 / 做到）垃圾分类。\rule{2cm}{0.5pt}（保护 / 爱护）环境，从我做起。
\par
\textbf{Solution détaillée :}
\begin{enumerate}
    \item \zh{问题} (wèntí) = problème. \zh{题目} = sujet d'examen/titre.
    \item \zh{乱扔} (luàn rēng) = jeter en désordre (V+V de manière). \zh{扔乱} n'existe pas.
    \item \zh{也} (yě) = aussi. \zh{还} (hái) = encore/en plus (souvent pour actions futures ou additionnelles). Ici, constat simultané : \zh{也}.
    \item \zh{应该} (yīnggāi) = devoir moral/conseil. \zh{可以} = pouvoir/permission.
    \item \zh{节约} (jiéyuē) = économiser (ressources). \zh{节省} = économiser (souvent argent/temps), moins collé à \zh{用水}.
    \item \zh{不应该} (bù yīnggāi) = ne pas devoir (interdiction/conseil négatif). \zh{没应该} \textbf{incorrect} (\zh{没} nie le passé/accompli, pas le modal).
    \item \zh{乱扔} (voir 2).
    \item \zh{做到} (zuòdào) = réaliser/appliquer (un principe, une règle). Collocation fixe : \zh{做到垃圾分类}. \zh{做} seul est trop vague.
    \item \zh{保护} (bǎohù) = protéger (environnement, nature). \zh{爱护} (àihù) = prendre soin de / chérir (objets, livres, animaux). Les deux possibles, \zh{保护} plus formel pour « environnement ».
\end{enumerate}
\textbf{Texte final :}
\zh{在我们的学校，环境问题不少。同学们经常乱扔纸屑，水龙头也滴水。我们应该节约用水，不应该乱扔垃圾，要做到垃圾分类。保护环境，从我做起。}\\
Zài wǒmen de xuéxiào, huánjìng wèntí bù shǎo. Tóngxuémen jīngcháng luàn rēng zhǐxiè, shuǐlóngtóu yě dī shuǐ. Wǒmen yīnggāi jiéyuē yòng shuǐ, bù yīnggāi luàn rēng lājī, yào zuòdào lājī fēnlèi. Bǎohù huánjìng, cóng wǒ zuòqǐ.
\end{exoresolu}

\begin{attention}[Piège majeur : La négation de \zh{应该}]
En français, « Je ne dois pas » = obligation de ne pas faire (interdiction).
En chinois, \zh{不应该} (bù yīnggāi) exprime exactement cela : \textbf{« Il ne faut pas / Il est mal de... »}.
\par
\textbf{Erreur fréquente :} \zh{*应该不} (ordre impossible) ou \zh{*不用} (pas besoin de -- absence d'obligation, sens différent).
\par
\textbf{Exemples :}
\begin{itemize}
    \item \zh{我们不应该乱扔垃圾。} (Wǒmen bù yīnggāi luàn rēng lājī.) -- \textbf{Interdiction / Devoir de ne pas faire.}
    \item \zh{我们不用买瓶装水。} (Wǒmen bù yòng mǎi píngzhuāng shuǐ.) -- \textbf{Absence de nécessité} (on a une gourde, par ex.).
\end{itemize}
\cle{Ne confondez pas \zh{不应该} (interdiction morale) et \zh{不用} (absence de besoin).}
\end{attention}

\courssub{Production individuelle : « Mon geste éco au lycée / au quartier »}

\emph{Consigne :} Rédigez un court paragraphe de \textbf{4 à 5 phrases} en caractères chinois (avec pinyin si besoin pour vous relire) sur le thème : \textbf{« Un geste concret que je fais / peux faire pour l'environnement dans mon école ou mon quartier »}.

Utilisez au moins \textbf{3 mots du glossaire} de la Section 3. Inspirez-vous de la trame du modèle mais personnalisez le lieu (\zh{在我们的学校}, \zh{在我的小区}, \zh{在雅温得}, \zh{在杜阿拉}) et les actions.

\begin{exemplebox}[Exemples de productions attendues (Niveau HSK 2-3)]
\textbf{Exemple A (Focus école - Yaoundé) :}
\zh{在我们的学校，垃圾分类还不够好。很多同学乱扔塑料瓶。我决定每天用环保袋，不买瓶装水。我们应该保护校园环境，人人做环保小卫士。}\\
Zài wǒmen de xuéxiào, lājī fēnlèi hái bù gòu hǎo. Hěnduō tóngxué luàn rēng sùliào píng. Wǒ juédìng měitiān yòng huánbào dài, bù mǎi píngzhuāng shuǐ. Wǒmen yīnggāi bǎohù xiàoyuán huánjìng, rénrén zuò huánbào xiǎo wèishì.\\
« Dans notre école, le tri des déchets n'est pas encore bon. Beaucoup de camarades jettent des bouteilles en plastique. Je décide d'utiliser un sac éco tous les jours, de ne plus acheter d'eau en bouteille. Nous devons protéger l'environnement du campus, chacun être un petit garde écolo. »

\textbf{Exemple B (Focus quartier - Douala) :}
\zh{在我的小区，河水被污染了，很脏。邻居们应该少用塑料袋，多用篮子买菜。我和爸爸周末去植树，种了两棵树。保护环境，从现在做起。}\\
Zài wǒ de xiǎoqū, heshuǐ bèi wūrǎn le, hěn zāng. Línjūmen yīnggāi shǎo yòng sùliào dài, duō yòng lánzi mǎicài. Wǒ hé bàba zhōumò qù zhíshù, zhòng le liǎng kē shù. Bǎohù huánjìng, cóng xiànzài zuòqǐ.\\
« Dans mon quartier, l'eau de la rivière est polluée, très sale. Les voisins devraient moins utiliser de sacs plastique, plus utiliser des paniers pour acheter la nourriture. Mon père et moi, le week-end, on va planter des arbres, on a planté deux arbres. Protéger l'environnement, commencer dès maintenant. »
\end{exemplebox}

\courssub{Ancrage socioculturel : Comparaison Cameroun -- Chine}

La protection de l'environnement prend des formes différentes selon les contextes urbains, les politiques publiques et les habitudes. Comparons deux réalités proches de vous.

\begin{center}
\begin{tabularx}{\linewidth}{|X|X|X|}
\hline
\rowcolor{popblueL}
\textbf{Thème} & \textbf{Cameroun (ex: Douala, Yaoundé)} & \textbf{Chine (ex: Shanghai, Pékin)} \\
\hline
\textbf{Gestion des déchets plastiques} &
\zh{塑料垃圾} (sùliào lājī) omniprésents.
Caniveaux bouchés, brûlage à l'air libre.
Interdiction officielle sacs < 60µm (2012) mais application difficile.
&
\zh{垃圾分类} (lājī fēnlèi) \textbf{obligatoire} (Shanghai 2019).
4 catégories : \zh{干垃圾} (sec), \zh{湿垃圾} (humide), \zh{可回收物} (recyclable), \zh{有害垃圾} (dangereux).
Amendes pour non-respect. Apps de suivi. \\
\hline
\textbf{Reboisement / Lutte désertification} &
\zh{砍树} (kǎn shù) / \zh{砍伐森林} pour bois de chauffe, agriculture.
Projet « Grande Muraille Verte » sahélienne (transnational, lent).
&
\zh{种树} (zhòng shù) / \zh{植树节} (12 mars) fête nationale.
\zh{三北防护林} (Sānběi Fánghùlín) -- « Grande Muraille Verte » chinoise (1978--2050) : plus grand reboisement artificiel monde. \\
\hline
\textbf{Implication citoyenne} &
Associations locales, \zh{清洁运动} (nettoyage) ponctuels (Journée Mondiale Env.).
&
\zh{志愿者} (zhìyuànzhě) systémiques. Crédits carbone personnels (Ant Forest / \zh{蚂蚁森林} sur Alipay : marcher, payer sans papier $\rightarrow$ arbre réel planté). \\
\hline
\end{tabularx}
\end{center}

\begin{popfigure}
\begin{tikzpicture}[
    node distance=1.2cm and 1.5cm,
    box/.style={rectangle, rounded corners, draw=popdark, thick, fill=popcream, text width=4.5cm, align=center, minimum height=1.8cm},
    arrow/.style={-Stealth, thick, popblue},
    title/.style={font=\bfseries\large, text=popblue, fill=popblueL, rounded corners, draw=popblue, inner sep=4pt}
]
% Cameroon column
\node[title] (camtitle) {Cameroun};
\node[box, below=of camtitle, align=center] (cam1){\zh{塑料垃圾}\\(sùliào lājī)\\\emph{Déchets plastiques}};
\node[box, below=of cam1, align=center] (cam2){\zh{砍树 / 森林砍伐}\\(kǎn shù / sēnlín kǎnfá)\\\emph{Coupe d'arbres / Déforestation}};
\node[box, below=of cam2, align=center] (cam3){\zh{垃圾分类起步}\\(lājī fēnlèi qǐbù)\\\emph{Tri des déchets : débuts}};

% China column
\node[title, right=3.5cm of camtitle] (chntitle) {Chine};
\node[box, right=3.5cm of cam1, align=center] (chn1){\zh{垃圾分类强制}\\(lājī fēnlèi qiángzhì)\\\emph{Tri obligatoire}};
\node[box, right=3.5cm of cam2, align=center] (chn2){\zh{植树节 / 种树}\\(zhíshùjié / zhòng shù)\\\emph{Fête de l'arbre / Planter}};
\node[box, right=3.5cm of cam3, align=center] (chn3){\zh{蚂蚁森林 / 碳普惠}\\(mǎyǐ sēnlín / tàn pǔhuì)\\\emph{Fourmi Forest / Crédits carbone}};

% Arrows / Links
\draw[arrow, dashed] (cam1.east) -- node[above, font=\small, text=poporange] {Évolution visée} (chn1.west);
\draw[arrow, dashed] (cam2.east) -- node[above, font=\small, text=poporange] {Modèle massif} (chn2.west);
\draw[arrow, dashed] (cam3.east) -- node[above, font=\small, text=poporange] {Innovation numérique} (chn3.west);
\end{tikzpicture}
\legende{Comparatif visuel : Défis camerounais (gauche) vs Politiques/Pratiques chinoises (droite). Les flèches pointent vers des leviers d'action ou modèles inspirants.}
\end{popfigure}

\begin{savaistu}[La « Grande Muraille Verte » : Deux continents, un même combat]
Saviez-vous que la \textbf{Chine} et l'\textbf{Afrique (Sahel)} mènent un projet homonyme ?
\begin{itemize}
    \item \textbf{Chine :} \zh{三北防护林工程} (Sānběi Fánghùlín Gōngchéng), lancé en 1978. Objectif : planter une ceinture forestière de 4 500 km à travers le Nord, le Nord-Est et le Nord-Ouest pour stopper l'avancée du désert de Gobi. C'est le plus grand projet de reboisement artificiel de l'histoire. La Chine est devenue le pays qui a le plus augmenté sa surface forestière ces 20 dernières années.
    \item \textbf{Afrique :} \textbf{Grande Muraille Verte pour le Sahara et le Sahel} (initiative UA, 2007). 8 000 km de Dakar à Djibouti. Le Cameroun (région Nord/Extrême-Nord) y participe activement : lutte contre l'érosion, restauration des terres, agroforesterie.
\end{itemize}
\cle{Point commun :} La plantation d'arbres (\zh{植树} / \zh{造林}) n'est pas seulement écologique, c'est une question de \textbf{sécurité alimentaire} et de \textbf{moyens de subsistance} pour les populations locales des deux côtés.
\end{savaistu}

\courssub{Expression orale : Lecture et correction collective}

\begin{experience}[Activité : « Ma meilleure phrase »]
\textbf{Déroulement (10 min) :}
\begin{enumerate}
    \item \textbf{Préparation (2 min) :} Chaque élève relit sa production écrite, choisit \textbf{sa phrase préférée} (la plus correcte, la plus originale ou celle dont il est le plus fier).
    \item \textbf{Lecture à voix haute (5 min) :} Tour de classe rapide. L'élève lit sa phrase en chinois (caractères ou pinyin). La classe écoute.
    \item \textbf{Feedback immédiat (3 min) :} Le professeur note au tableau 2 à 3 phrases entendues (anonymisées) pour une micro-correction collective.
\end{enumerate}
\textbf{Critères d'écoute pour les camarades :} Tons justes ? Ordre des mots (Lieu avant Verbe) ? Vocabulaire du thème utilisé ?
\end{experience}

\begin{exoresolu}[Correction collective type : Analyse de deux productions anonymes]
\textbf{Production 1 (Élève X) :}
\zh{在我们学校，我们应该节约水，不乱扔垃圾，要做垃圾分类。保护环境很重要。}\\
Zài wǒmen xuéxiào, wǒmen yīnggāi jiéyuē shuǐ, bù luàn rēng lājī, yào zuò lājī fēnlèi. Bǎohù huánjìng hěn zhòngyào.
\par
\textbf{Analyse :}
\begin{itemize}
    \item \textbf{Points forts :} Structure \zh{在我们学校} (Locatif initial) correcte. Enchaînement \zh{应该...，不...，要...} maîtrisé. Vocabulaire précis (\zh{节约水}, \zh{垃圾分类}). Conclusion claire.
    \item \textbf{Point perfectible :} \zh{节约水} $\rightarrow$ mieux : \zh{节约用水} (collocation standard \zh{节约} + \zh{用水/电/纸}). \zh{做垃圾分类} $\rightarrow$ mieux : \zh{做到垃圾分类} ou \zh{实行垃圾分类}.
\end{itemize}

\textbf{Production 2 (Élève Y) :}
\zh{我喜欢环境保护。在雅温得，有很多垃圾。我们不应该用塑料袋。我种一棵树。}\\
Wǒ xǐhuan huánjìng bǎohù. Zài Yǎwēndé, yǒu hěnduō lājī. Wǒmen bù yīnggāi yòng sùliào dài. Wǒ zhòng yī kē shù.
\par
\textbf{Analyse :}
\begin{itemize}
    \item \textbf{Points forts :} Phrases courtes, simples, communication réussie. \zh{不应该用} correct. \zh{种一棵树} (classificateur \zh{棵} pour arbre) excellent !
    \item \textbf{Points à corriger :}
    \begin{enumerate}
        \item \zh{我喜欢环境保护} : \zh{喜欢} (aimer) + nom abstrait d'action $\rightarrow$ préférer \zh{我关心环境保护} (je me soucie de...) ou \zh{我支持环境保护}.
        \item \zh{有很多垃圾} : Correct, mais \zh{垃圾很严重} (les déchets sont un problème grave) ou \zh{到处是垃圾} (il y a des déchets partout) plus naturel.
        \item \zh{我种一棵树} : Aspect accompli \zh{了} manquant si c'est fait (\zh{我种了一棵树}) ; ou futur \zh{我要种一棵树} / \zh{我打算种一棵树} si c'est un projet.
    \end{enumerate}
\end{itemize}
\textbf{Bilan au tableau :} Écrire les versions corrigées à côté des originaux. Insister sur \zh{节约用水} vs \zh{节约水} et l'aspect \zh{了} / \zh{要} pour \zh{种树}.
\end{exoresolu}

\courssub{Bilan de la leçon : Auto-évaluation des compétences}

Remplissez cette grille honnêtement. Elle vous servira de base pour la révision du Module 2.

\begin{center}
\begin{tabularx}{\linewidth}{|X|c|c|c|}
\hline
\rowcolor{popblueL}
\textbf{Compétence visée (Programme MINESEC 1ère A)} & \zh{ Maîtrisé} & \zh{ En cours} & \zh{☹ Difficulté} \\
\hline
Je reconnais et prononce (tons) le vocabulaire de l'environnement (\zh{污染, 保护, 分类, 节约, 植树}...). & & & \\
\hline
Je lis et comprends un texte court (200-300 mots) sur l'écologie sans dictionnaire. & & & \\
\hline
J'utilise correctement \zh{应该 / 不应该 / 还要} pour exprimer le devoir/conseil. & & & \\
\hline
J'écris un paragraphe cohérent (4-5 phrases) sur un geste éco (SVO, ordre mots, tons). & & & \\
\hline
Je compare une situation camerounaise et chinoise (déchets, arbres) en chinois simple. & & & \\
\hline
Je connais l'ordre des traits et le radical de \zh{污, 染, 护, 种, 棵}. & & & \\
\hline
\end{tabularx}
\end{center}

\begin{methode}[Comment progresser sur les items « En cours / Difficulté » ?]
\begin{enumerate}
    \item \textbf{Vocabulaire / Tons :} Refaites les cartes flash (Anki/Quizlet) en \textbf{disant le mot à haute voix} avant de retourner la carte. Groupez par radical (\zh{氵} pour eau/liquide : \zh{污, 河, 洗, 浇}).
    \item \textbf{Grammaire \zh{应该} :} Écrivez 3 phrases \zh{应该} + 3 phrases \zh{不应该} sur votre cahier ce soir. Vérifiez l'ordre : Suj + \zh{不} + \zh{应该} + V.
    \item \textbf{Écriture caractères :} Utilisez la méthode « \textbf{Doigt sur la paume} » : tracez le caractère en grand sur la paume de la main gauche avec l'index droit en disant l'ordre des traits (ex: \zh{污} : \zh{氵} 3 traits, puis \zh{亏} -- haut/bas/gauche/droite).
    \item \textbf{Comparaison Culturelle :} Regardez une vidéo courte (3 min) sur \zh{垃圾分类上海} ou \zh{蚂蚁森林} sur YouTube/Bilibili (sous-titres chinois). Notez 3 mots entendus.
\end{enumerate}
\end{methode}

\begin{exercice}[Devoirs / Approfondissement personnel]
\begin{enumerate}
    \item \textbf{Rédaction soignée :} Recopiez votre paragraphe personnel (Section 6.3) sur votre cahier de \textbf{texte chinois} en soignant l'écriture des caractères (grille 10x10 ou 12x12). Ajoutez le pinyin au-dessus des mots nouveaux uniquement.
    \item \textbf{Traduction (Version) :} Traduisez en français le texte suivant (vocabulaire connu + 2 mots nouveaux contextuels) :
    \par
    \zh{上海市民现在垃圾分类做得很好。干垃圾扔黑色桶，湿垃圾扔棕色桶。如果不分类，要罚款。大家都觉得保护环境是自己的事。}
    \item \textbf{Recherche culturelle (Oral prochain cours) :} Préparez 2 phrases en chinois pour présenter \textbf{une initiative écologique locale au Cameroun} (ex: \zh{雅温得清洁日}, \zh{芒贝绿色城市项目}, \zh{回收塑料瓶协会}). Vocabulaire clé à préparer : \zh{项目} (projet), \zh{志愿者} (bénévole), \zh{社区} (communauté/quartier), \zh{改善} (améliorer).
\end{enumerate}
\end{exercice}

\begin{corrige}[Exercice 2 -- Traduction]
\textbf{Texte source :}
\zh{上海市民现在垃圾分类做得很好。干垃圾扔黑色桶，湿垃圾扔棕色桶。如果不分类，要罚款。大家都觉得保护环境是自己的事。}\\
Shànghǎi shìmín xiànzài lājī fēnlèi zuò de hěn hǎo. Gān lājī rēng hēisè tǒng, shī lājī rēng zōngsè tǒng. Rúguǒ bù fēnlèi, yào fákuǎn. Dàjiā dōu juéde bǎohù huánjìng shì zìjǐ de shì.

\textbf{Traduction :}
« Les habitants de Shanghai font maintenant le tri des déchets très bien. Les déchets secs vont dans la poubelle noire, les déchets humides dans la poubelle brune. Si on ne trie pas, il faut payer une amende. Tout le monde trouve que protéger l'environnement, c'est sa propre affaire / responsabilité. »

\textbf{Notes lexicales :}
\begin{itemize}
    \item \zh{市民} (shìmín) : citadins / habitants de la ville.
    \item \zh{做得很好} (zuò de hěn hǎo) : structure \zh{得} (de) complément de degré -- « faire \textbf{bien} ».
    \item \zh{干垃圾} (gān lājī) / \zh{湿垃圾} (shī lājī) : déchets secs / humides (terminologie officielle Shanghai).
    \item \zh{桶} (tǒng) : seau, poubelle, conteneur. Classificateur : \zh{个} 或 \zh{只}.
    \item \zh{罚款} (fákuǎn) : amende (verbe/nom). \zh{要罚款} = risquer une amende / devoir payer une amende.
    \item \zh{自己的事} (zìjǐ de shì) : sa propre affaire / responsabilité personnelle.
\end{itemize}
\end{corrige}

\newpage

\coursec{Activité d'intégration}

\courssub{Objectifs de l'activité}

À la fin de cette activité d'intégration, tu seras capable de :

\begin{itemize}
\item mobiliser le \cle{vocabulaire} de la protection de l'environnement (\zh{环境}， \zh{保护}， \zh{节约}， \zh{垃圾}， \zh{树}…) dans une production authentique ;
\item employer correctement les \cle{collocations} étudiées : \zh{保护环境}， \zh{节约用水}， \zh{节约用电}， \zh{乱扔垃圾}， \zh{多种树} ;
\item rédiger un \cle{slogan} et des conseils courts en chinois simplifié, avec une ponctuation chinoise correcte ;
\item écrire soigneusement au moins quatre caractères du thème en respectant l'\cle{ordre des traits} ;
\item lire à voix haute un texte chinois avec une prononciation correcte (pinyin et tons) ;
\item travailler en équipe et évaluer la production des autres groupes selon des critères précis.
\end{itemize}

\courssub{Situation}

Le club écologique de ton lycée à Yaoundé organise la « Semaine verte ». Les élèves de la section chinois ont été choisis pour décorer les couloirs avec des affiches de sensibilisation \textbf{bilingues chinois-français}, car le lycée accueille cette année une délégation d'élèves d'une école partenaire de Chine. Le proviseur a annoncé que la meilleure affiche serait reproduite et affichée à l'entrée de l'établissement.

Voici le message que la présidente du club écologique a envoyé à tous les groupes :

\begin{textebox}[Message du club écologique]
同学们好！\\
Tóngxuémen hǎo!\\
« Bonjour à tous les élèves ! »\\[4pt]
下个星期是我们学校的“绿色周”。\\
Xià ge xīngqī shì wǒmen xuéxiào de “lǜsè zhōu”.\\
« La semaine prochaine, c'est la Semaine verte de notre école. »\\[4pt]
请每个小组做一张海报，题目是“保护我们的环境”。\\
Qǐng měi ge xiǎozǔ zuò yì zhāng hǎibào, tímù shì “bǎohù wǒmen de huánjìng”.\\
« Chaque groupe doit faire une affiche intitulée "Protégeons notre environnement". »\\[4pt]
海报上要有一个口号和三个建议。\\
Hǎibào shàng yào yǒu yí ge kǒuhào hé sān ge jiànyì.\\
« L'affiche doit contenir un slogan et trois conseils. »\\[4pt]
我们要选出最好的海报，放在学校门口。\\
Wǒmen yào xuǎn chū zuì hǎo de hǎibào, fàng zài xuéxiào ménkǒu.\\
« Nous choisirons la meilleure affiche pour la placer à l'entrée de l'école. »
\end{textebox}

\courssub{Consignes}

Par groupes de 4, vous créez une affiche intitulée \zh{保护我们的环境} (\emph{bǎohù wǒmen de huánjìng}, « Protégeons notre environnement »). Travaillez étape par étape :

\begin{enumerate}
\item \textbf{Compréhension du message.} Relisez le message du club écologique et répondez oralement : Quel est le thème de la semaine ? Que doit contenir l'affiche ? Où sera exposée la meilleure affiche ?
\item \textbf{Choix du slogan.} Écrivez un slogan court en chinois. Vous pouvez reprendre le modèle \zh{保护环境，人人有责！} (\emph{bǎohù huánjìng, rén rén yǒu zé!}, « Protéger l'environnement, c'est l'affaire de tous ! ») ou en créer un nouveau avec le vocabulaire du thème.
\item \textbf{Rédaction des trois conseils.} Rédigez trois conseils en utilisant les collocations apprises : \zh{节约用水}， \zh{节约用电}， \zh{不乱扔垃圾}， \zh{多种树}. Chaque conseil doit être une phrase complète et correcte (ex. \zh{我们要节约用水。}).
\item \textbf{Écriture des caractères.} Choisissez au moins quatre caractères du thème (par exemple \zh{环}， \zh{境}， \zh{保}， \zh{护}， \zh{树}， \zh{水}) et écrivez-les en grand sur l'affiche, en respectant l'ordre des traits : de haut en bas, de gauche à droite, l'horizontal avant le vertical.
\item \textbf{Version bilingue.} Ajoutez sous chaque phrase chinoise sa traduction française en plus petit.
\item \textbf{Présentation orale.} Chaque groupe présente son affiche en une minute : un élève lit le slogan, les autres lisent chacun un conseil. Attention aux tons !
\item \textbf{Vote de la classe.} La classe vote pour la meilleure affiche selon trois critères : exactitude du vocabulaire, qualité de la prononciation, soin de l'écriture des caractères. L'affiche gagnante sera exposée dans la salle.
\end{enumerate}

\courssub{Ressources à mobiliser}

\begin{aretenir}[Vocabulaire et collocations utiles]
\begin{itemize}
\item \zh{保护环境} (\emph{bǎohù huánjìng}) : protéger l'environnement ; \zh{保护我们的环境} : protéger notre environnement.
\item \zh{节约用水} (\emph{jiéyuē yòng shuǐ}) : économiser l'eau ; \zh{节约用电} (\emph{jiéyuē yòng diàn}) : économiser l'électricité.
\item \zh{不乱扔垃圾} (\emph{bù luàn rēng lājī}) : ne pas jeter les déchets n'importe où.
\item \zh{多种树} (\emph{duō zhòng shù}) : planter plus d'arbres ; \zh{多} + verbe = « faire davantage ».
\item Structure des conseils : \textbf{Sujet + 要 / 应该 + verbe} → \zh{我们要节约用水。} (\emph{Wǒmen yào jiéyuē yòng shuǐ.}, « Nous devons économiser l'eau. »)
\item Structure du slogan : \zh{人人有责} (\emph{rén rén yǒu zé}) : « chacun est responsable ».
\end{itemize}
\end{aretenir}

\begin{attention}[Pièges à éviter]
\begin{itemize}
\item N'écris pas \zh{乱扔垃圾} sans la négation dans un conseil : il faut \zh{不要乱扔垃圾} ou \zh{不乱扔垃圾} (« ne pas jeter »).
\item \zh{水} s'écrit en 4 traits : le trait vertical central d'abord, puis les deux côtés. \zh{树} : la clé du bois \zh{木} à gauche d'abord.
\item Ponctuation chinoise pleine chasse : \zh{，} \zh{。} \zh{！} et non la virgule ou le point français.
\item Prononce bien les tons : \zh{shuǐ} (3\up{e} ton), \zh{lājī} (1\up{er} ton), \zh{diàn} (4\up{e} ton).
\end{itemize}
\end{attention}

\courssub{Proposition de production et corrigé commenté}

\begin{exoresolu}[Modèle d'affiche du groupe « Les Verts de Yaoundé »]
Rédige l'affiche complète de ton groupe : titre, slogan, trois conseils avec traduction française, et quatre caractères du thème écrits soigneusement.
\tcblower
\textbf{Titre de l'affiche :}\\
\zh{保护我们的环境}\\
\emph{bǎohù wǒmen de huánjìng} — « Protégeons notre environnement »\\[4pt]
\textbf{Slogan :}\\
\zh{保护环境，人人有责！}\\
\emph{bǎohù huánjìng, rén rén yǒu zé!} — « Protéger l'environnement, c'est l'affaire de tous ! »\\[4pt]
\textbf{Trois conseils :}\\
1. \zh{我们要节约用水。} \emph{Wǒmen yào jiéyuē yòng shuǐ.} « Nous devons économiser l'eau. »\\
2. \zh{我们要节约用电。} \emph{Wǒmen yào jiéyuē yòng diàn.} « Nous devons économiser l'électricité. »\\
3. \zh{我们不应该乱扔垃圾，要多种树。} \emph{Wǒmen bù yīnggāi luàn rēng lājī, yào duō zhòng shù.} « Nous ne devons pas jeter les déchets n'importe où, nous devons planter plus d'arbres. »\\[4pt]
\textbf{Caractères écrits en grand :} \zh{环}（clé du jade 王， 8 traits）、\zh{境}（clé de la terre 土， 14 traits）、\zh{树}（clé du bois 木， 9 traits）、\zh{水}（4 traits）。\\[4pt]
\textbf{Commentaire des choix de langue :} le slogan reprend la structure figée \zh{人人有责}, très fréquente dans les campagnes de sensibilisation en Chine. Les conseils utilisent \zh{要} (« devoir, il faut ») et \zh{不应该} (« ne devrait pas »), ce qui permet d'alterner conseil positif et interdiction. La troisième phrase combine deux collocations avec la virgule chinoise \zh{，}, ce qui enrichit la production sans dépasser le niveau HSK 2-3. Les caractères choisis couvrent trois clés différentes du thème (jade, terre, bois), ce qui montre la maîtrise des radicaux étudiés.
\end{exoresolu}

\courssub{Bilan de l'activité}

Cette activité t'a permis de réinvestir l'ensemble des savoir-faire de la leçon dans une tâche concrète et motivante :

\begin{itemize}
\item \textbf{Compréhension de l'écrit :} tu as lu et exploité un message authentique en chinois pour en dégager les consignes de production ;
\item \textbf{Lexique :} tu as employé le vocabulaire de la protection de l'environnement et ses collocations dans des phrases correctes ;
\item \textbf{Écriture :} tu as tracé des caractères du thème en respectant les clés et l'ordre des traits ;
\item \textbf{Expression orale :} tu as lu ton affiche à voix haute en soignant les tons ;
\item \textbf{Compétence sociale :} tu as coopéré en groupe, présenté un travail collectif et évalué celui des autres selon des critères objectifs.
\end{itemize}

\begin{savaistu}[Le saviez-vous ?]
En Chine comme au Cameroun, les campagnes de sensibilisation à l'environnement utilisent des slogans courts et rythmés, faciles à retenir. Le slogan \zh{保护环境，人人有责！} est l'un des plus connus en Chine : on le voit dans les écoles, les parcs et les transports. À Yaoundé et à Douala, les journées de nettoyage des quartiers rappellent la tradition chinoise des actions collectives pour l'embellissement des villes : dans les deux pays, protéger l'environnement est considéré comme l'affaire de tous, exactement ce que dit ton affiche !
\end{savaistu}

\newpage

\coursec{Méthodes et exercices résolus}

\courssub{Mise en route et découverte du thème}

\begin{experience}[Brainstorming : l'environnement autour de nous]
En binômes, répondez d'abord en français, puis essayez en chinois avec l'aide du professeur :
\begin{enumerate}
\item Quels problèmes d'environnement vois-tu dans ta ville (Yaoundé, Douala, Buea) ? Essaie de dire : \zh{垃圾很多。} \emph{Lājī hěn duō.} « Il y a beaucoup de déchets. »
\item Que font les élèves de ton lycée pour garder la cour propre ? Modèle : \zh{我们种树。} \emph{Wǒmen zhòng shù.} « Nous plantons des arbres. »
\item Connais-tu le slogan chinois \zh{保护环境，人人有责} ? Devine son sens avant de regarder la traduction : \emph{Bǎohù huánjìng, rén rén yǒu zé.} « Protéger l'environnement, c'est l'affaire de tous. »
\end{enumerate}
\end{experience}

\courssub{Texte support : lecture et compréhension globale}

\begin{textebox}[保护我们的环境 (Protégeons notre environnement)]
\zh{现在，环境污染是一个大问题。城市的空气不太好，河水也很脏。为了保护环境，我们应该少开车，多骑自行车，还应该节约用水。去年，我们学校的学生在校园里种了一百多棵树。大家都说：保护环境，人人有责。}\\
\emph{Xiànzài, huánjìng wūrǎn shì yí ge dà wèntí. Chéngshì de kōngqì bù tài hǎo, hé shuǐ yě hěn zāng. Wèile bǎohù huánjìng, wǒmen yīnggāi shǎo kāi chē, duō qí zìxíngchē, hái yīnggāi jiéyuē yòng shuǐ. Qùnián, wǒmen xuéxiào de xuésheng zài xiàoyuán lǐ zhòng le yì bǎi duō kē shù. Dàjiā dōu shuō : bǎohù huánjìng, rén rén yǒu zé.}\\
« Aujourd'hui, la pollution de l'environnement est un grand problème. L'air des villes n'est pas très bon, et l'eau des rivières est aussi très sale. Pour protéger l'environnement, nous devrions moins conduire, plus faire de vélo, et aussi économiser l'eau. L'année dernière, les élèves de notre école ont planté plus de cent arbres dans la cour de l'école. Tout le monde dit : protéger l'environnement, c'est l'affaire de tous. »
\end{textebox}

\begin{aretenir}[Ce qu'il faut retenir de la première lecture]
Le texte suit un plan en trois temps typique des textes de sensibilisation : (1) le constat du problème (\zh{污染} pollution, \zh{脏} sale) ; (2) les actions proposées (\zh{少开车}, \zh{多骑自行车}, \zh{节约用水}, \zh{种树}) ; (3) le slogan final (\zh{保护环境，人人有责}). Ce plan te servira de modèle pour tes propres productions écrites.
\end{aretenir}

\courssub{Glossaire du thème}

\begin{center}
\begin{tabularx}{\linewidth}{|X|X|X|X|}
\hline
\rowcolor{popblueL} Caractères & Pinyin & Nature & Traduction \\
\hline
环境 & huánjìng & nom & environnement \\
\hline
保护 & bǎohù & verbe & protéger \\
\hline
污染 & wūrǎn & verbe / nom & polluer ; pollution \\
\hline
垃圾 & lājī & nom & déchets, ordures \\
\hline
空气 & kōngqì & nom & air \\
\hline
河水 & hé shuǐ & nom & eau de rivière \\
\hline
干净 & gānjìng & adjectif & propre \\
\hline
脏 & zāng & adjectif & sale \\
\hline
节约 & jiéyuē & verbe & économiser \\
\hline
回收 & huíshōu & verbe & recycler \\
\hline
扔 & rēng & verbe & jeter \\
\hline
捡 & jiǎn & verbe & ramasser \\
\hline
种 & zhòng & verbe & planter \\
\hline
树 & shù & nom & arbre \\
\hline
棵 & kē & classificateur & (classificateur des arbres et plantes) \\
\hline
条 & tiáo & classificateur & (classificateur des rivières, poissons, routes) \\
\hline
应该 & yīnggāi & verbe modal & devoir (conseil) \\
\hline
为了 & wèile & préposition & pour, afin de \\
\hline
问题 & wèntí & nom & problème, question \\
\hline
校园 & xiàoyuán & nom & cour de l'école, campus \\
\hline
\end{tabularx}
\end{center}

\begin{attention}[Prononciation : les tons qui piègent les francophones]
\begin{itemize}
\item \zh{水} shuǐ (3\up{e} ton, « eau ») et \zh{树} shù (4\up{e} ton, « arbre ») : même initiale \emph{sh}, mais des tons différents. Les confondre change le sens.
\item \zh{河} hé (2\up{e} ton, « rivière ») : ne prononce jamais hě.
\item \zh{垃圾} lājī : les deux syllabes au 1\up{er} ton, haut et plat.
\item Sandhi de \zh{不} : devant une syllabe au 4\up{e} ton, \zh{不} passe au 2\up{e} ton : \zh{不要} se prononce \emph{búyào}, \zh{不太} se prononce \emph{bú tài}.
\end{itemize}
\end{attention}

\courssub{Clés (radicaux) et ordre des traits}

\begin{definition}[Les clés du thème environnement]
\begin{itemize}
\item \zh{河} (hé, rivière) : clé de l'\cle{eau} \zh{氵} (trois gouttes, à gauche) + \zh{可} (kě). Ordre : d'abord les trois gouttes de haut en bas, puis la partie droite.
\item \zh{树} (shù, arbre) : clé du \cle{bois} \zh{木} (à gauche) + \zh{对}. Ordre : \zh{木} (horizontal, vertical, puis les deux obliques), puis la partie droite, de gauche à droite.
\item \zh{垃} et \zh{圾} (lājī, déchets) : clé de la \cle{terre} \zh{土} (à gauche). Attention : \zh{垃} se termine par \zh{立}, \zh{圾} par \zh{及}.
\item \zh{境} (jìng, dans \zh{环境}) : clé de la \cle{terre} \zh{土} + \zh{竟}. Ne pas confondre avec \zh{镜} (jìng, miroir), qui a la clé du métal \zh{钅}.
\item \zh{保} (bǎo, protéger) : clé de l'\cle{homme} \zh{亻} (à gauche) + \zh{呆}. Ordre : d'abord le radical homme (deux traits), puis la partie droite de haut en bas.
\end{itemize}
Règle générale d'écriture : de haut en bas, de gauche à droite, l'horizontal avant le vertical, l'extérieur avant l'intérieur.
\end{definition}

\courssub{Savoir-faire 1 : Lire et comprendre un court texte chinois sur l'environnement}

\begin{methode}[Comprendre un texte sur la protection de l'environnement]
\begin{enumerate}
\item \textbf{Première lecture globale} : lis le texte une première fois sans t'arrêter sur les mots inconnus. Repère le thème général (ici : la pollution, le tri des déchets, la protection de la nature) et les mots déjà connus : \zh{环境} (environnement), \zh{保护} (protéger), \zh{垃圾} (déchets), \zh{空气} (air), \zh{水} (eau).
\item \textbf{Repérage des mots-clés du thème} : souligne les mots du champ lexical de l'environnement : \zh{污染} (pollution), \zh{干净} (propre), \zh{脏} (sale), \zh{种树} (planter des arbres), \zh{节约用水} (économiser l'eau).
\item \textbf{Découpage des phrases} : en chinois, cherche le sujet, puis le verbe, puis le complément. La place des compléments de temps et de lieu est fixe : Sujet + temps + lieu + verbe.
\item \textbf{Deuxième lecture analytique} : pour chaque phrase, traduis mentalement en respectant l'ordre chinois. Utilise le contexte pour deviner les mots inconnus (ex. : \zh{回收} près de \zh{垃圾} = recycler).
\item \textbf{Vérification} : relis le texte entier et résume-le en une phrase en français avant de répondre aux questions.
\end{enumerate}
\end{methode}

\begin{exoresolu}[Compréhension de texte : 保护我们的环境]
\textbf{Énoncé.} Lis le texte support de la leçon, puis réponds aux questions en chinois par des phrases complètes.

\textbf{Questions :}
\begin{enumerate}
\item \zh{城市的空气怎么样？}
\item \zh{为了保护环境，我们应该做什么？（说出两个行动）}
\item \zh{去年，学生们在校园里做了什么？}
\end{enumerate}
\tcblower
\textbf{Solution détaillée.}

\textbf{Question 1.} La question porte sur \zh{怎么样} (comment est...?). On cherche dans le texte la phrase qui parle de l'air des villes : \zh{城市的空气不太好}. Réponse rédigée en phrase complète :

\zh{城市的空气不太好。} \emph{Chéngshì de kōngqì bù tài hǎo.} « L'air des villes n'est pas très bon. »

\textbf{Question 2.} Le texte donne les actions après \zh{我们应该} : \zh{少开车，多骑自行车，节约用水}. Il suffit d'en choisir deux et de rédiger :

\zh{为了保护环境，我们应该少开车，多骑自行车。} \emph{Wèile bǎohù huánjìng, wǒmen yīnggāi shǎo kāi chē, duō qí zìxíngchē.} « Pour protéger l'environnement, nous devrions moins conduire et plus faire de vélo. »

Justification grammaticale : \zh{为了} + but, \zh{应该} + verbe (modalité du devoir), \zh{少}/\zh{多} + verbe expriment « moins / plus faire ».

\textbf{Question 3.} On repère \zh{去年} (l'année dernière) puis l'action : \zh{种了一百多棵树}. Attention au classificateur des arbres : \zh{棵} (kē). Réponse :

\zh{去年，学生们在校园里种了一百多棵树。} \emph{Qùnián, xuéshengmen zài xiàoyuán lǐ zhòng le yì bǎi duō kē shù.} « L'année dernière, les élèves ont planté plus de cent arbres dans la cour de l'école. »

\textbf{Conclusion.} La méthode gagnante : localiser dans le texte la phrase qui contient le mot-clé de la question (\zh{空气}, \zh{应该}, \zh{去年}), puis recopier l'information en phrase complète, sans ajouter d'élément absent du texte.
\end{exoresolu}

\courssub{Savoir-faire 2 : Reconnaître, prononcer et employer le vocabulaire du thème}

\begin{methode}[Mémoriser et employer le vocabulaire de l'environnement]
\begin{enumerate}
\item \textbf{Apprends les mots par familles} : le problème (\zh{污染} pollution, \zh{垃圾} déchets, \zh{脏} sale), la nature (\zh{空气} air, \zh{水} eau, \zh{树} arbre, \zh{河} rivière), les solutions (\zh{保护} protéger, \zh{节约} économiser, \zh{回收} recycler, \zh{种树} planter des arbres).
\item \textbf{Travaille les tons à voix haute} : répète chaque mot trois fois en exagérant les tons. Pièges fréquents : \zh{水} shuǐ (3\up{e} ton), \zh{树} shù (4\up{e} ton) — ne pas confondre ; \zh{河} hé (2\up{e} ton) — ne pas dire hě.
\item \textbf{Apprends les collocations} (mots qui vont ensemble), pas les mots isolés : \zh{保护环境} (protéger l'environnement), \zh{节约用水} (économiser l'eau), \zh{扔垃圾} (jeter des déchets), \zh{空气污染} (pollution de l'air).
\item \textbf{Associe chaque nom à son classificateur} : \zh{一棵树} (un arbre), \zh{一条河} (une rivière), \zh{一个问题} (un problème).
\item \textbf{Réemploie} : fabrique une phrase personnelle avec chaque mot nouveau, par exemple sur ta ville (Yaoundé, Douala).
\end{enumerate}
\end{methode}

\begin{exoresolu}[Vocabulaire : compléter et traduire]
\textbf{Énoncé.} Complète les phrases avec le mot qui convient, choisi dans la liste : \zh{保护 / 污染 / 垃圾 / 节约 / 棵}. Puis traduis chaque phrase en français.
\begin{enumerate}
\item \zh{工厂的废气＿＿了空气。}
\item \zh{请不要在地上扔＿＿。}
\item \zh{我们应该＿＿用水。}
\item \zh{我家门前有两＿＿树。}
\end{enumerate}
\tcblower
\textbf{Solution détaillée.}

\textbf{1.} Le sujet est \zh{废气} (gaz d'échappement) et le complément \zh{空气} (air) : il faut un verbe qui relie les deux. \zh{污染} (polluer) convient ; la particule \zh{了} marque l'accomplissement. \zh{工厂的废气污染了空气。} \emph{Gōngchǎng de fèiqì wūrǎn le kōngqì.} « Les gaz d'échappement des usines ont pollué l'air. »

\textbf{2.} Après le verbe \zh{扔} (jeter), on attend un nom d'objet : \zh{垃圾} (déchets). \zh{请不要在地上扔垃圾。} \emph{Qǐng búyào zài dìshang rēng lājī.} « S'il vous plaît, ne jetez pas de déchets par terre. » Structure à retenir : \zh{请不要} + verbe = « prière de ne pas... ».

\textbf{3.} Devant \zh{用水} (utiliser l'eau), le mot qui forme la collocation correcte est \zh{节约} (économiser) : \zh{节约用水}. \zh{我们应该节约用水。} \emph{Wǒmen yīnggāi jiéyuē yòng shuǐ.} « Nous devrions économiser l'eau. »

\textbf{4.} Entre le chiffre \zh{两} et le nom \zh{树}, il faut obligatoirement un classificateur : celui des arbres est \zh{棵}. \zh{我家门前有两棵树。} \emph{Wǒ jiā mén qián yǒu liǎng kē shù.} « Devant ma maison, il y a deux arbres. » Rappel : avec « deux », on dit \zh{两} (et non \zh{二}) devant un classificateur.

\textbf{Conclusion.} Pour réussir ce type d'exercice, raisonne sur la nature du mot attendu (verbe ? nom ? classificateur ?) et sur les collocations apprises.
\end{exoresolu}

\courssub{Savoir-faire 3 : Répondre à des questions de compréhension en phrases complètes}

\begin{methode}[Rédiger une réponse de compréhension en chinois]
\begin{enumerate}
\item \textbf{Analyse le mot interrogatif} : \zh{什么} (quoi), \zh{谁} (qui), \zh{哪儿/哪里} (où), \zh{什么时候} (quand), \zh{为什么} (pourquoi), \zh{怎么样} (comment), \zh{几/多少} (combien).
\item \textbf{Localise la réponse dans le texte} grâce au mot-clé de la question ; ne réponds jamais de mémoire ou avec tes opinions si la question porte sur le texte.
\item \textbf{Rédige une phrase complète} en reprenant la structure de la question : remplace simplement le mot interrogatif par l'information. Exemple : \zh{我们应该做什么？} → \zh{我们应该种树。}
\item \textbf{Respecte l'ordre des mots chinois} : Sujet + (temps) + (lieu) + verbe + complément. Ne caleque jamais l'ordre français.
\item \textbf{Relis} : vérifie les tons du pinyin si on te le demande, les classificateurs et la particule \zh{了} pour les actions passées accomplies.
\end{enumerate}
\end{methode}

\begin{exoresolu}[Questions de compréhension sur un texte]
\textbf{Énoncé.} Texte : \zh{我的家乡在杜阿拉。以前，海边的垃圾很多，水很脏。从今年开始，每个星期六早上，很多年轻人去海边捡垃圾。现在，海边干净多了，来玩儿的人也多了。}

\emph{Wǒ de jiāxiāng zài Dù'ālā. Yǐqián, hǎibiān de lājī hěn duō, shuǐ hěn zāng. Cóng jīnnián kāishǐ, měi ge xīngqīliù zǎoshang, hěn duō niánqīngrén qù hǎibiān jiǎn lājī. Xiànzài, hǎibiān gānjìng duō le, lái wánr de rén yě duō le.}

« Ma ville natale est Douala. Avant, il y avait beaucoup de déchets au bord de la mer, l'eau était très sale. À partir de cette année, chaque samedi matin, beaucoup de jeunes vont ramasser les déchets au bord de la mer. Maintenant, le bord de la mer est bien plus propre, et il y a aussi plus de gens qui viennent s'y promener. »

\textbf{Questions :} 1. \zh{以前，海边怎么样？} 2. \zh{年轻人什么时候去捡垃圾？} 3. \zh{现在，海边有什么变化？}
\tcblower
\textbf{Solution détaillée.}

\textbf{Question 1.} Mot interrogatif : \zh{怎么样} (comment). Le repère temporel est \zh{以前} (avant). Le texte dit : \zh{垃圾很多，水很脏}. Réponse complète : \zh{以前，海边的垃圾很多，水很脏。} \emph{Yǐqián, hǎibiān de lājī hěn duō, shuǐ hěn zāng.} « Avant, il y avait beaucoup de déchets au bord de la mer et l'eau était très sale. »

\textbf{Question 2.} Mot interrogatif : \zh{什么时候} (quand). Repère : \zh{每个星期六早上}. Réponse : \zh{他们每个星期六早上去海边捡垃圾。} \emph{Tāmen měi ge xīngqīliù zǎoshang qù hǎibiān jiǎn lājī.} « Ils vont ramasser les déchets au bord de la mer chaque samedi matin. » Justification : en chinois, le complément de temps \zh{每个星期六早上} se place entre le sujet et le verbe, jamais à la fin comme en français.

\textbf{Question 3.} La question porte sur le changement (\zh{变化}). Le texte donne deux résultats : \zh{干净多了} et \zh{来玩儿的人也多了}. Réponse : \zh{现在，海边干净多了，来玩儿的人也多了。} \emph{Xiànzài, hǎibiān gānjìng duō le, lái wánr de rén yě duō le.} « Maintenant, le bord de la mer est bien plus propre et il y a plus de visiteurs. » Structure à noter : adjectif + \zh{多了} = « beaucoup plus... ».

\textbf{Conclusion.} Une bonne réponse reprend les mots du texte, dans une phrase complète, avec l'ordre des mots chinois.
\end{exoresolu}

\courssub{Savoir-faire 4 : Traduire un message sur l'environnement (version et thème)}

\begin{methode}[Traduire sans calquer le français]
\begin{enumerate}
\item \textbf{Identifie d'abord le sujet et le verbe} de la phrase française, puis reconstruis la phrase selon le schéma chinois : Sujet + temps + lieu + verbe + complément.
\item \textbf{Place les compléments de temps et de lieu avant le verbe} : « Nous plantons des arbres à l'école le samedi » → \zh{我们星期六在学校种树。}
\item \textbf{Choisis le bon classificateur} pour chaque nom : \zh{一棵树}, \zh{一条河}, \zh{一个问题}.
\item \textbf{Utilise les mots-outils correctement} : \zh{了} pour une action accomplie, \zh{的} entre un déterminant et un nom, \zh{应该} pour le devoir, \zh{为了} pour le but.
\item \textbf{Relis en chinois uniquement} : la phrase doit sonner chinoise, pas être du français mot à mot.
\end{enumerate}
\end{methode}

\begin{exoresolu}[Thème : traduire en chinois]
\textbf{Énoncé.} Traduis en chinois (caractères + pinyin) :
\begin{enumerate}
\item « Pour protéger l'environnement, nous devons économiser l'eau. »
\item « Hier, les élèves ont planté dix arbres dans la cour de l'école. »
\item « L'air de la ville est plus sale qu'avant. »
\end{enumerate}
\tcblower
\textbf{Solution détaillée.}

\textbf{1.} Le but se traduit par \zh{为了} + proposition, placé en tête ; « devons » = \zh{应该} ; « économiser l'eau » = collocation \zh{节约用水}. \zh{为了保护环境，我们应该节约用水。} \emph{Wèile bǎohù huánjìng, wǒmen yīnggāi jiéyuē yòng shuǐ.}

\textbf{2.} « Hier » = \zh{昨天}, placé après le sujet (ou en tête de phrase) ; l'action est accomplie, donc \zh{了} après le verbe \zh{种} ; « dix arbres » exige le classificateur \zh{棵} : \zh{十棵树} ; le lieu \zh{在校园里} se place avant le verbe. \zh{昨天，同学们在校园里种了十棵树。} \emph{Zuótiān, tóngxuémen zài xiàoyuán lǐ zhòng le shí kē shù.} Erreur à éviter : ne pas mettre \zh{在校园里} après le verbe comme en français.

\textbf{3.} La comparaison « plus... que » se construit avec \zh{比} : A + \zh{比} + B + adjectif. « Avant » = \zh{以前}. \zh{城市的空气比以前脏。} \emph{Chéngshì de kōngqì bǐ yǐqián zāng.} On peut renforcer : \zh{比以前脏多了} (« bien plus sale qu'avant »).

\textbf{Conclusion.} En traduction, la priorité absolue est l'ordre des mots : temps et lieu avant le verbe, comparaison avec \zh{比}, classificateur entre le nombre et le nom.
\end{exoresolu}

\courssub{Savoir-faire 5 : Produire un court message écrit sur l'environnement}

\begin{methode}[Rédiger un court texte de sensibilisation]
\begin{enumerate}
\item \textbf{Structure en trois temps} : (a) constat du problème avec \zh{现在……是一个大问题} ; (b) actions proposées avec \zh{我们应该……还应该……} ; (c) conclusion avec un slogan : \zh{保护环境，人人有责。}
\item \textbf{Utilise 4 à 6 mots du glossaire} imposés ou appris : \zh{污染}, \zh{垃圾}, \zh{保护}, \zh{节约}, \zh{种树}, \zh{干净}.
\item \textbf{Phrases courtes} : une idée par phrase, 8 à 15 caractères. Mieux vaut cinq phrases simples correctes que deux phrases longues fautives.
\item \textbf{Connecteurs utiles} : \zh{因为……所以……} (parce que... donc...), \zh{还} (aussi), \zh{但是} (mais), \zh{如果……就……} (si... alors...).
\item \textbf{Relecture ciblée} : vérifie chaque classificateur, chaque \zh{了}, et l'ordre Sujet + temps + lieu + verbe.
\end{enumerate}
\end{methode}

\begin{exoresolu}[Production guidée : un message pour la journée de l'environnement]
\textbf{Énoncé.} Pour la journée de l'environnement de ton lycée à Yaoundé, rédige en chinois un court message (4 à 6 phrases) en employant au moins quatre mots du glossaire : \zh{污染、垃圾、保护、节约、种树、干净}.
\tcblower
\textbf{Solution détaillée (modèle de rédaction).}

\zh{现在，环境污染是一个大问题。我们学校的旁边有很多垃圾，空气也不太好。为了保护环境，我们应该少扔垃圾，节约用水。这个星期六，我们班要在校园里种树。如果每个人都努力，我们的城市就会更干净。保护环境，人人有责！}

\emph{Xiànzài, huánjìng wūrǎn shì yí ge dà wèntí. Wǒmen xuéxiào de pángbiān yǒu hěn duō lājī, kōngqì yě bù tài hǎo. Wèile bǎohù huánjìng, wǒmen yīnggāi shǎo rēng lājī, jiéyuē yòng shuǐ. Zhège xīngqīliù, wǒmen bān yào zài xiàoyuán lǐ zhòng shù. Rúguǒ měi ge rén dōu nǔlì, wǒmen de chéngshì jiù huì gèng gānjìng. Bǎohù huánjìng, rén rén yǒu zé!}

« Aujourd'hui, la pollution de l'environnement est un grand problème. À côté de notre école, il y a beaucoup de déchets, et l'air n'est pas très bon. Pour protéger l'environnement, nous devrions jeter moins de déchets et économiser l'eau. Ce samedi, notre classe va planter des arbres dans la cour de l'école. Si chacun fait des efforts, notre ville sera plus propre. Protéger l'environnement, c'est l'affaire de tous ! »

\textbf{Analyse de la copie.} (1) Le message suit le plan constat → actions → slogan. (2) Les six mots du glossaire sont employés. (3) Structures grammaticales correctes : \zh{为了} + but, \zh{应该} + verbe, \zh{如果……就……} pour l'hypothèse, \zh{更} + adjectif pour « plus... ». (4) Ordre des mots respecté : \zh{这个星期六} (temps) et \zh{在校园里} (lieu) sont placés avant le verbe \zh{种树}.

\textbf{Conclusion.} Un message réussi est court, suit un plan clair, réutilise le vocabulaire du thème et se termine par un slogan mémorable.
\end{exoresolu}

\begin{savaistu}[Le tri des déchets en Chine et au Cameroun]
Depuis quelques années, les grandes villes chinoises comme Shanghai ont rendu le tri des déchets obligatoire : on sépare les déchets recyclables (\zh{可回收垃圾} kě huíshōu lājī), les déchets de cuisine et les autres déchets. Au Cameroun, des associations de jeunes organisent régulièrement des journées de nettoyage (\zh{捡垃圾} jiǎn lājī, « ramasser les déchets ») dans les quartiers de Yaoundé et de Douala, et des entreprises transforment les déchets plastiques en pavés pour les rues. En Chine comme au Cameroun, on retrouve la même idée : \zh{保护环境，人人有责} — protéger l'environnement, c'est l'affaire de tous.
\end{savaistu}

\begin{attention}[Les 5 erreurs qui coûtent des points à l'examen]
\begin{enumerate}
\item \textbf{Confondre les tons de mots proches} : \zh{水} shuǐ (3\up{e} ton, eau) et \zh{树} shù (4\up{e} ton, arbre) ; \zh{河} hé (2\up{e} ton, rivière) prononcé hě est une faute. À l'écrit, ne confonds pas les caractères \zh{境} (dans \zh{环境}) et \zh{镜} (miroir), ni \zh{垃} et \zh{拉}.
\item \textbf{Oublier le classificateur} : on n'écrit jamais \zh{两树} ni \zh{一百树}. Il faut \zh{两棵树}, \zh{一百棵树}, \zh{一条河}, \zh{一个问题}. Le classificateur est obligatoire entre le nombre (ou le démonstratif) et le nom.
\item \textbf{Calquer l'ordre des mots français} : « Nous plantons des arbres à l'école le samedi » ne se traduit pas mot à mot. En chinois, temps et lieu précèdent le verbe : \zh{我们星期六在学校种树。} et non \zh{我们种树在学校星期六}.
\item \textbf{Mal employer les mots-outils} : \zh{了} marque l'action accomplie (\zh{我们种了十棵树。}) mais ne s'emploie pas avec les verbes d'état ni les phrases habituelles ; \zh{的} relie un déterminant au nom (\zh{城市的空气}) ; la comparaison exige \zh{比} (\zh{现在比以前干净多了}), jamais un calque de « plus... que ».
\item \textbf{Traduire les collocations mot à mot} : « économiser l'eau » se dit \zh{节约用水} (et non \zh{节约水} dans un style soigné) ; « protéger l'environnement » = \zh{保护环境} ; « jeter des déchets » = \zh{扔垃圾}. Apprends ces blocs entiers : ce sont eux qui sont attendus à l'examen.
\end{enumerate}
\end{attention}

\begin{exercice}[Pour s'entraîner]
\begin{enumerate}
\item Recopie les caractères \zh{河}, \zh{树}, \zh{境}, \zh{保} cinq fois chacun en respectant l'ordre des traits, et entoure la clé de chacun.
\item Traduis en français : \zh{我们应该回收垃圾，保护环境。}
\item Traduis en chinois (caractères + pinyin) : « Chaque samedi, nous ramassons les déchets dans la cour de l'école. »
\item Réponds en chinois par une phrase complète : \zh{你的城市干净吗？为什么？}
\end{enumerate}
\end{exercice}

\begin{corrige}[Exercice « Pour s'entraîner »]
\textbf{1.} Clés : \zh{河} → \zh{氵} (eau) ; \zh{树} → \zh{木} (bois) ; \zh{境} → \zh{土} (terre) ; \zh{保} → \zh{亻} (homme). Écriture : de gauche à droite (radical d'abord), de haut en bas.

\textbf{2.} \zh{我们应该回收垃圾，保护环境。} \emph{Wǒmen yīnggāi huíshōu lājī, bǎohù huánjìng.} « Nous devrions recycler les déchets et protéger l'environnement. »

\textbf{3.} \zh{每个星期六，我们在校园里捡垃圾。} \emph{Měi ge xīngqīliù, wǒmen zài xiàoyuán lǐ jiǎn lājī.} Le temps \zh{每个星期六} et le lieu \zh{在校园里} précèdent le verbe \zh{捡}. Pas de \zh{了} : il s'agit d'une habitude, pas d'une action accomplie ponctuelle.

\textbf{4.} Modèle de réponse : \zh{我的城市不太干净，因为地上的垃圾很多。} \emph{Wǒ de chéngshì bù tài gānjìng, yīnwèi dìshang de lājī hěn duō.} « Ma ville n'est pas très propre, parce qu'il y a beaucoup de déchets par terre. » Toute réponse complète avec \zh{因为} et le vocabulaire du thème est acceptée.
\end{corrige}

\newpage

\coursec{Exercices}

\courssub{Je vérifie mes connaissances}

\begin{exercice}[QCM : le bon mot]
Entoure la bonne réponse.\par
1. « Protéger l'environnement » se dit :\par
a) \zh{保护环境} \quad b) \zh{污染环境} \quad c) \zh{打扫环境}\par
2. \zh{垃圾} signifie :\par
a) l'eau \quad b) les ordures \quad c) l'air\par
3. « Économiser l'eau » se dit :\par
a) \zh{喝水} \quad b) \zh{买水} \quad c) \zh{节约用水}\par
4. \zh{空气} signifie :\par
a) la terre \quad b) l'arbre \quad c) l'air\par
5. « Planter des arbres » se dit :\par
a) \zh{砍树} \quad b) \zh{种树} \quad c) \zh{买树}
\end{exercice}

\begin{exercice}[Vrai ou faux ? Justifie]
Dis si l'affirmation est vraie ou fausse, puis justifie en français à partir du texte support de la leçon.\par
1. \zh{地球是我们的家。} (Dìqiú shì wǒmen de jiā.) — La Terre est notre maison.\par
2. \zh{乱扔垃圾可以保护环境。} (Luàn rēng lājī kěyǐ bǎohù huánjìng.) — Jeter les ordures n'importe où protège l'environnement.\par
3. \zh{我们应该节约用水。} (Wǒmen yīnggāi jiéyuē yòng shuǐ.) — Nous devons économiser l'eau.\par
4. \zh{种树对环境没有好处。} (Zhòng shù duì huánjìng méiyǒu hǎochù.) — Planter des arbres n'est pas bon pour l'environnement.
\end{exercice}

\begin{exercice}[Vocabulaire : problème ou solution ?]
Associe chaque mot chinois à sa traduction française, puis classe-le dans la colonne « Problèmes » ou « Solutions ».\par
\begin{center}
\begin{tabularx}{\linewidth}{|X|X|X|}\hline
\rowcolor{popblueL} Mot chinois & Traduction française & Problème / Solution \\ \hline
\zh{污染} (wūrǎn) & ................ & ................ \\ \hline
\zh{保护} (bǎohù) & ................ & ................ \\ \hline
\zh{垃圾} (lājī) & ................ & ................ \\ \hline
\zh{节约} (jiéyuē) & ................ & ................ \\ \hline
\zh{废气} (fèiqì) & ................ & ................ \\ \hline
\zh{种树} (zhòng shù) & ................ & ................ \\ \hline
\zh{扔} (rēng) & ................ & ................ \\ \hline
\zh{干净} (gānjìng) & ................ & ................ \\ \hline
\zh{脏} (zāng) & ................ & ................ \\ \hline
\zh{回收} (huíshōu) & ................ & ................ \\ \hline
\end{tabularx}
\end{center}
Traductions à placer : propre ; sale ; recycler ; la pollution ; protéger ; les ordures ; économiser ; les gaz d'échappement ; jeter ; planter des arbres.
\end{exercice}

\begin{exercice}[Pinyin et caractères]
Relie chaque caractère ou mot à son pinyin.\par
\begin{center}
\begin{tabularx}{\linewidth}{|X|X|}\hline
\rowcolor{popblueL} Caractères & Pinyin \\ \hline
1. \zh{环境} & a) shù \\ \hline
2. \zh{保护} & b) huánjìng \\ \hline
3. \zh{污染} & c) jiéyuē \\ \hline
4. \zh{树} & d) bǎohù \\ \hline
5. \zh{节约} & e) wūrǎn \\ \hline
6. \zh{垃圾} & f) lājī \\ \hline
\end{tabularx}
\end{center}
\end{exercice}

\courssub{Je m'entraîne}

\begin{exercice}[Phrases à trous]
Complète les phrases avec les mots suivants : \zh{保护} (bǎohù), \zh{污染} (wūrǎn), \zh{节约} (jiéyuē), \zh{扔} (rēng), \zh{种} (zhòng).\par
1. 我们应该\_\_\_\_用水。\par
2. 请不要乱\_\_\_\_垃圾。\par
3. 工厂的废气\_\_\_\_了空气。\par
4. 我们要\_\_\_\_环境，保护地球。\par
5. 明天我们去公园\_\_\_\_树。
\end{exercice}

\begin{exercice}[Remise en ordre]
Remets les mots dans l'ordre pour former une phrase correcte, puis traduis-la en français.\par
1. 我们 / 应该 / 环境 / 保护\par
2. 水 / 和 / 空气 / 被 / 污染了\par
3. 垃圾 / 不要 / 请 / 乱扔\par
4. 越来越 / 天气 / 热了 / 热\par
5. 种树 / 大家 / 一起 / 去 / 我们
\end{exercice}

\begin{exercice}[Transformation : 很 → 越来越]
Transforme chaque phrase en remplaçant \zh{很} par la structure \zh{越来越} + adjectif + \zh{了}, puis traduis les deux phrases en français.\par
Exemple : \zh{天气很热。} → \zh{天气越来越热了。} (Tiānqì yuè lái yuè rè le. — Il fait de plus en plus chaud.)\par
1. \zh{城市很脏。} (Chéngshì hěn zāng.)\par
2. \zh{河水很干净。} (Hé shuǐ hěn gānjìng.)\par
3. \zh{空气污染很严重。} (Kōngqì wūrǎn hěn yánzhòng.)\par
4. \zh{我们的校园很美。} (Wǒmen de xiàoyuán hěn měi.)
\end{exercice}

\begin{exercice}[Traduction : français → chinois]
Traduis les phrases suivantes en chinois (caractères + pinyin).\par
1. L'air et l'eau sont pollués.\par
2. Nous devons protéger la Terre.\par
3. Ne jetez pas les ordures par terre.\par
4. À Yaoundé, les élèves plantent des arbres dans la cour de l'école.\par
5. Tout le monde doit économiser l'eau et l'électricité.
\end{exercice}

\courssub{Je me prépare au Bac}

\begin{exercice}[Compréhension écrite type Bac]
Lis le texte suivant, puis réponds aux questions.\par
\begin{textebox}[La ville de Lánshān]
蓝山市以前很干净，空气很好，河水也很清。\\
Lánshān shì yǐqián hěn gānjìng, kōngqì hěn hǎo, hé shuǐ yě hěn qīng.\\
« Autrefois, la ville de Lanshan était très propre, l'air était bon et l'eau de la rivière était claire. »\\
可是现在，工厂越来越多，汽车也越来越多，空气和河水都被污染了。\\
Kěshì xiànzài, gōngchǎng yuè lái yuè duō, qìchē yě yuè lái yuè duō, kōngqì hé hé shuǐ dōu bèi wūrǎn le.\\
« Mais maintenant, il y a de plus en plus d'usines et de voitures ; l'air et l'eau de la rivière sont tous pollués. »\\
市政府决定保护环境：他们种了很多树，也请大家不要乱扔垃圾。\\
Shì zhèngfǔ juédìng bǎohù huánjìng : tāmen zhòng le hěn duō shù, yě qǐng dàjiā bú yào luàn rēng lājī.\\
« La mairie a décidé de protéger l'environnement : ils ont planté beaucoup d'arbres et demandent à tout le monde de ne pas jeter les ordures n'importe où. »\\
人们还应该节约用水、节约用电。\\
Rénmen hái yīnggāi jiéyuē yòng shuǐ, jiéyuē yòng diàn.\\
« Les gens doivent aussi économiser l'eau et l'électricité. »\\
大家都希望蓝山市越来越干净、越来越美。\\
Dàjiā dōu xīwàng Lánshān shì yuè lái yuè gānjìng, yuè lái yuè měi.\\
« Tout le monde espère que Lanshan deviendra de plus en plus propre et de plus en plus belle. »
\end{textebox}
\textbf{Questions de compréhension}\par
A. Vrai ou faux ? Justifie par une phrase du texte (en chinois).\par
1. \zh{蓝山市以前很脏。}\par
2. \zh{现在蓝山市的空气被污染了。}\par
3. \zh{市政府种了很多树。}\par
B. Réponds en chinois, par une phrase complète.\par
4. \zh{为什么蓝山市的空气和河水被污染了？}\par
5. \zh{为了保护环境，人们应该做什么？}（Donne deux actions.）\par
\textbf{Questions de langue}\par
6. Releve dans le texte deux phrases avec la structure \zh{越来越} et traduis-les.\par
7. Releve la phrase passive avec \zh{被} et explique sa construction.
\end{exercice}

\begin{exercice}[Thème et version]
\textbf{Thème} — Traduis en chinois (caractères + pinyin) :\par
1. Notre ville est de plus en plus sale.\par
2. Nous ne devons pas jeter les ordures dans la rivière.\par
3. Les élèves du lycée de Douala plantent des arbres chaque année.\par
\textbf{Version} — Traduis en français :\par
4. \zh{地球是我们共同的家，我们要爱护它。} (Dìqiú shì wǒmen gòngtóng de jiā, wǒmen yào àihù tā.)\par
5. \zh{如果每个人都保护环境，世界会越来越美。} (Rúguǒ měi ge rén dōu bǎohù huánjìng, shìjiè huì yuè lái yuè měi.)
\end{exercice}

\begin{exercice}[Expression écrite type Bac]
Sujet : \zh{我们小区的环境} (Wǒmen xiǎoqū de huánjìng) — « L'environnement de mon quartier ».\par
Rédige un court texte de \textbf{5 à 8 phrases} (environ 60 à 90 caractères chinois) sur les gestes écologiques dans ton quartier au Cameroun.\par
Consignes :\par
\begin{itemize}
\item Utilise au moins \textbf{4 mots du glossaire} de la leçon (par exemple : \zh{保护}, \zh{垃圾}, \zh{节约}, \zh{种树}, \zh{干净}).
\item Utilise au moins une fois la structure \zh{越来越}.
\item Utilise au moins une fois \zh{应该} ou \zh{不要}.
\item Écris les caractères avec soin, dans l'ordre correct des traits ; ajoute le pinyin sous chaque phrase.
\end{itemize}
\end{exercice}

\newpage

\coursec{Corrigés des exercices}

\coursec{Corrigés détaillés — Leçon 15 : Protection de l'environnement}

\courssub{Exercices de vocabulaire, grammaire, compréhension et expression}

\begin{corrige}[Exercice 1 — QCM : le bon mot]
1. \textbf{b) est faux ; la bonne réponse est a)} \zh{保护环境} (bǎohù huánjìng) — « protéger l'environnement ». \zh{污染} (wūrǎn) signifie « polluer » et \zh{打扫} (dǎsǎo) signifie « balayer, nettoyer » ; aucun des deux ne convient.\par
2. \textbf{b)} \zh{垃圾} (lājī) signifie « les ordures, les déchets ». L'eau se dit \zh{水} (shuǐ) et l'air \zh{空气} (kōngqì).\par
3. \textbf{c)} \zh{节约用水} (jiéyuē yòng shuǐ) — « économiser l'eau ». \zh{喝水} (hē shuǐ) = « boire de l'eau » ; \zh{买水} (mǎi shuǐ) = « acheter de l'eau ».\par
4. \textbf{c)} \zh{空气} (kōngqì) signifie « l'air ». La terre se dit \zh{地} ou \zh{土地} (tǔdì) ; l'arbre se dit \zh{树} (shù).\par
5. \textbf{b)} \zh{种树} (zhòng shù) — « planter des arbres ». \zh{砍树} (kǎn shù) = « couper des arbres » ; \zh{买树} (mǎi shù) = « acheter des arbres ».\par
\textbf{Point de vigilance :} ne pas confondre \zh{保护} (bǎohù, protéger) et \zh{污染} (wūrǎn, polluer) : ce sont des contraires fréquents dans les QCM.
\end{corrige}

\begin{corrige}[Exercice 2 — Vrai ou faux ? Justifie]
1. \textbf{Vrai.} Le texte dit : \zh{地球是我们的家。} (Dìqiú shì wǒmen de jiā.) — « La Terre est notre maison. » C'est l'idée de départ du texte : nous devons en prendre soin comme de notre maison.\par
2. \textbf{Faux.} Jeter les ordures n'importe où \emph{détruit} l'environnement, il ne le protège pas. Le texte dit le contraire : \zh{请不要乱扔垃圾。} (Qǐng bú yào luàn rēng lājī.) — « Ne jetez pas les ordures n'importe où. »\par
3. \textbf{Vrai.} Le texte affirme : \zh{我们应该节约用水。} (Wǒmen yīnggāi jiéyuē yòng shuǐ.) — « Nous devons économiser l'eau. » L'eau est une ressource précieuse.\par
4. \textbf{Faux.} Planter des arbres est \emph{bon} pour l'environnement : les arbres rendent l'air plus propre. Le texte dit : \zh{种树可以保护环境。} (Zhòng shù kěyǐ bǎohù huánjìng.) — « Planter des arbres peut protéger l'environnement. »\par
\textbf{Point de vigilance :} dans la phrase 2, \zh{乱扔} (luàn rēng) = « jeter n'importe où » ; l'adverbe \zh{乱} (luàn, « en désordre ») se place \emph{avant} le verbe.
\end{corrige}

\begin{corrige}[Exercice 3 — Vocabulaire : problème ou solution ?]
\begin{center}
\begin{tabularx}{\linewidth}{|X|X|X|}\hline
\rowcolor{popblueL} Mot chinois & Traduction française & Problème / Solution \\ \hline
\zh{污染} (wūrǎn) & la pollution & Problème \\ \hline
\zh{保护} (bǎohù) & protéger & Solution \\ \hline
\zh{垃圾} (lājī) & les ordures & Problème \\ \hline
\zh{节约} (jiéyuē) & économiser & Solution \\ \hline
\zh{废气} (fèiqì) & les gaz d'échappement & Problème \\ \hline
\zh{种树} (zhòng shù) & planter des arbres & Solution \\ \hline
\zh{扔} (rēng) & jeter & Problème \\ \hline
\zh{干净} (gānjìng) & propre & Solution \\ \hline
\zh{脏} (zāng) & sale & Problème \\ \hline
\zh{回收} (huíshōu) & recycler & Solution \\ \hline
\end{tabularx}
\end{center}
\textbf{Point de vigilance :} \zh{干净} (gānjìng, propre) et \zh{脏} (zāng, sale) sont des adjectifs contraires ; attention au ton de \zh{脏} : premier ton zāng (sale), à ne pas confondre avec zàng (quatrième ton) dans \zh{心脏} (xīnzàng, le cœur).
\end{corrige}

\begin{corrige}[Exercice 4 — Pinyin et caractères]
1. \zh{环境} → \textbf{b)} huánjìng\par
2. \zh{保护} → \textbf{d)} bǎohù\par
3. \zh{污染} → \textbf{e)} wūrǎn\par
4. \zh{树} → \textbf{a)} shù\par
5. \zh{节约} → \textbf{c)} jiéyuē\par
6. \zh{垃圾} → \textbf{f)} lājī\par
\textbf{Point de vigilance :} dans \zh{垃圾} (lājī), les deux caractères portent la clé de la terre \zh{土} ; le pinyin est lājī (premier ton + premier ton), et non « lèsè » (transcription approximative du minnan, le mandarin standard est lājī).
\end{corrige}

\begin{corrige}[Exercice 5 — Phrases à trous]
1. 我们应该\zh{节约}用水。 (Wǒmen yīnggāi jiéyuē yòng shuǐ.) — « Nous devons économiser l'eau. »\par
2. 请不要乱\zh{扔}垃圾。 (Qǐng bú yào luàn rēng lājī.) — « Ne jetez pas les ordures n'importe où, s'il vous plaît. »\par
3. 工厂的废气\zh{污染}了空气。 (Gōngchǎng de fèiqì wūrǎn le kōngqì.) — « Les gaz d'échappement des usines ont pollué l'air. »\par
4. 我们要\zh{保护}环境，保护地球。 (Wǒmen yào bǎohù huánjìng, bǎohù Dìqiú.) — « Nous devons protéger l'environnement, protéger la Terre. »\par
5. 明天我们去公园\zh{种}树。 (Míngtiān wǒmen qù gōngyuán zhòng shù.) — « Demain, nous allons planter des arbres au parc. »\par
\textbf{Justification :} dans la phrase 3, \zh{污染} est ici un \emph{verbe} (« polluer ») suivi de \zh{了} qui marque l'action accomplie. Dans la phrase 5, \zh{种} se prononce zhòng (quatrième ton) quand il signifie « planter » ; zhǒng (troisième ton) signifie « graine, sorte ».
\end{corrige}

\begin{corrige}[Exercice 6 — Remise en ordre]
1. 我们应该保护环境。 (Wǒmen yīnggāi bǎohù huánjìng.) — « Nous devons protéger l'environnement. » Ordre : Sujet + auxiliaire \zh{应该} + verbe + complément.\par
2. 水和空气被污染了。 (Shuǐ hé kōngqì bèi wūrǎn le.) — « L'eau et l'air sont pollués. » Structure passive : Sujet (subi) + \zh{被} + verbe + \zh{了}.\par
3. 请不要乱扔垃圾。 (Qǐng bú yào luàn rēng lājī.) — « Ne jetez pas les ordures n'importe où, s'il vous plaît. » \zh{请} en tête pour la politesse, puis la négation \zh{不要}.\par
4. 天气越来越热了。 (Tiānqì yuè lái yuè rè le.) — « Il fait de plus en plus chaud. » Structure : Sujet + \zh{越来越} + adjectif + \zh{了}.\par
5. 我们大家一起去种树。 (Wǒmen dàjiā yìqǐ qù zhòng shù.) — « Nous allons tous ensemble planter des arbres. » Les adverbes \zh{都/一起} se placent avant le verbe \zh{去}.\par
\textbf{Point de vigilance :} en chinois, l'ordre est rigide : Sujet → adverbes/auxiliaires → verbe → complément. Ne jamais placer le complément d'objet avant le verbe (sauf structures spéciales comme \zh{把}, non exigée ici).
\end{corrige}

\begin{corrige}[Exercice 7 — Transformation : 很 → 越来越]
1. \zh{城市很脏。} (Chéngshì hěn zāng. — La ville est très sale.) → \zh{城市越来越脏了。} (Chéngshì yuè lái yuè zāng le. — La ville devient de plus en plus sale.)\par
2. \zh{河水很干净。} (Hé shuǐ hěn gānjìng. — L'eau de la rivière est très propre.) → \zh{河水越来越干净了。} (Hé shuǐ yuè lái yuè gānjìng le. — L'eau de la rivière devient de plus en plus propre.)\par
3. \zh{空气污染很严重。} (Kōngqì wūrǎn hěn yánzhòng. — La pollution de l'air est très grave.) → \zh{空气污染越来越严重了。} (Kōngqì wūrǎn yuè lái yuè yánzhòng le. — La pollution de l'air devient de plus en plus grave.)\par
4. \zh{我们的校园很美。} (Wǒmen de xiàoyuán hěn měi. — Notre campus est très beau.) → \zh{我们的校园越来越美了。} (Wǒmen de xiàoyuán yuè lái yuè měi le. — Notre campus devient de plus en plus beau.)\par
\textbf{Justification grammaticale :} la structure \zh{越来越} + adjectif + \zh{了} exprime un changement progressif. Deux règles : (1) on supprime \zh{很}, qui est incompatible avec \zh{越来越} ; (2) le \zh{了} final est obligatoire, il marque le nouveau état constaté.
\end{corrige}

\begin{corrige}[Exercice 8 — Traduction : français → chinois]
1. 空气和水都被污染了。 (Kōngqì hé shuǐ dōu bèi wūrǎn le.) — Passif avec \zh{被} ; \zh{都} résume les deux éléments.\par
2. 我们应该保护地球。 (Wǒmen yīnggāi bǎohù Dìqiú.)\par
3. 请不要把垃圾扔在地上。 (Qǐng bú yào bǎ lājī rēng zài dìshang.) — Variante plus simple acceptée : \zh{请不要乱扔垃圾。} (Qǐng bú yào luàn rēng lājī.)\par
4. 在雅温得，学生们在学校里种树。 (Zài Yǎwēndé, xuéshengmen zài xuéxiào lǐ zhòng shù.) — Le complément de lieu \zh{在雅温得} se place en tête ou après le sujet, jamais à la fin.\par
5. 大家都应该节约用水、节约用电。 (Dàjiā dōu yīnggāi jiéyuē yòng shuǐ, jiéyuē yòng diàn.)\par
\textbf{Points de vigilance :} (1) « par terre » se rend par \zh{在地上} (zài dìshang) placé \emph{après} le verbe ; (2) « l'électricité » = \zh{电} (diàn) ; (3) la virgule de coordination chinoise \zh{、} s'emploie entre deux éléments parallèles.
\end{corrige}

\begin{corrige}[Exercice 9 — Compréhension écrite type Bac]
\textbf{Barème indicatif (sur 20) :} A : 3 × 2 pts (1 pt réponse + 1 pt justification) ; B : 2 × 3 pts ; questions de langue : 6 : 4 pts, 7 : 4 pts.\par
\textbf{A. Vrai ou faux}\par
1. \textbf{Faux.} Le texte dit : \zh{蓝山市以前很干净。} — « Autrefois, Lanshan était très propre », et non sale.\par
2. \textbf{Vrai.} Le texte dit : \zh{空气和河水都被污染了。} — « L'air et l'eau de la rivière sont tous pollués. »\par
3. \textbf{Vrai.} Le texte dit : \zh{他们种了很多树。} — « Ils ont planté beaucoup d'arbres » (\zh{他们} désigne la mairie).\par
\textbf{B. Réponses en phrases complètes}\par
4. \zh{因为现在工厂越来越多，汽车也越来越多，所以空气和河水都被污染了。} (Yīnwèi xiànzài gōngchǎng yuè lái yuè duō, qìchē yě yuè lái yuè duō, suǒyǐ kōngqì hé hé shuǐ dōu bèi wūrǎn le.) — « Parce que maintenant il y a de plus en plus d'usines et de voitures, l'air et l'eau sont pollués. »\par
5. \zh{人们应该种树，还应该节约用水、节约用电。} (Rénmen yīnggāi zhòng shù, hái yīnggāi jiéyuē yòng shuǐ, jiéyuē yòng diàn.) — « Les gens doivent planter des arbres et aussi économiser l'eau et l'électricité. » (On accepte aussi : \zh{不要乱扔垃圾}.)\par
\textbf{Questions de langue}\par
6. Deux phrases avec \zh{越来越} : \zh{工厂越来越多，汽车也越来越多。} — « Il y a de plus en plus d'usines et de plus en plus de voitures. » ; \zh{大家都希望蓝山市越来越干净、越来越美。} — « Tout le monde espère que Lanshan deviendra de plus en plus propre et de plus en plus belle. »\par
7. Phrase passive : \zh{空气和河水都被污染了。} Construction : Sujet subissant l'action (\zh{空气和河水}) + \zh{被} (+ agent, ici sous-entendu) + verbe (\zh{污染}) + \zh{了}. Le sujet ne fait pas l'action, il la subit ; l'agent (« par les usines ») peut être omis quand il est évident.\par
\textbf{Points de vigilance :} répondre par une phrase complète (sujet + verbe), recopier la justification sans la modifier, et ne pas oublier le \zh{了} final dans les phrases avec \zh{越来越}.
\end{corrige}

\begin{corrige}[Exercice 10 — Thème et version]
\textbf{Thème}\par
1. 我们的城市越来越脏了。 (Wǒmen de chéngshì yuè lái yuè zāng le.) — « De plus en plus sale » impose la structure \zh{越来越} + adjectif + \zh{了}.\par
2. 我们不应该把垃圾扔到河里。 (Wǒmen bù yīnggāi bǎ lājī rēng dào hé lǐ.) — Variante acceptée : \zh{我们不应该往河里扔垃圾。} (Wǒmen bù yīnggāi wǎng hé lǐ rēng lājī.)\par
3. 杜阿拉高中的学生每年种树。 (Dù'ālā gāozhōng de xuésheng měi nián zhòng shù.) — Le complément de temps \zh{每年} (chaque année) se place après le sujet, avant le verbe.\par
\textbf{Version}\par
4. \zh{地球是我们共同的家，我们要爱护它。} — « La Terre est notre maison commune, nous devons en prendre soin (la chérir). » \zh{共同} (gòngtóng) = « commun » ; \zh{爱护} (àihù) = « chérir, prendre soin de » ; \zh{它} (tā) = « elle » (pour les choses).\par
5. \zh{如果每个人都保护环境，世界会越来越美。} — « Si chacun protège l'environnement, le monde deviendra de plus en plus beau. » Structure hypothétique \zh{如果……，……会……} (si…, alors…) ; \zh{每个人} (měi ge rén) = « chaque personne ».\par
\textbf{Barème indicatif :} 2 pts par phrase (1 pt exactitude des mots, 1 pt ordre des mots et particules).\par
\textbf{Points de vigilance :} en thème, méfie-toi de la place des compléments de temps et de lieu (avant le verbe) ; en version, rends \zh{越来越} par « de plus en plus » et non par « très ».
\end{corrige}

\begin{corrige}[Exercice 11 — Expression écrite type Bac]
\textbf{Barème indicatif (sur 20) :} respect du sujet et de la longueur (5–8 phrases, 60–90 caractères) : 4 pts ; vocabulaire de la leçon (au moins 4 mots) : 4 pts ; structures exigées (\zh{越来越}, \zh{应该}/\zh{不要}) : 4 pts ; correction grammaticale et ordre des mots : 4 pts ; écriture des caractères et pinyin : 4 pts.\par
\textbf{Proposition de rédaction modèle :}\par
我住在杜阿拉的一个小区。以前我们小区有点脏，垃圾很多。\\
(Wǒ zhù zài Dù'ālā de yí ge xiǎoqū. Yǐqián wǒmen xiǎoqū yǒudiǎn zāng, lājī hěn duō.)\\
« J'habite dans un quartier de Douala. Avant, notre quartier était un peu sale, il y avait beaucoup d'ordures. »\\
现在，大家都想保护环境。\\
(Xiànzài, dàjiā dōu xiǎng bǎohù huánjìng.)\\
« Maintenant, tout le monde veut protéger l'environnement. »\\
我们不应该乱扔垃圾，也应该节约用水。\\
(Wǒmen bù yīnggāi luàn rēng lājī, yě yīnggāi jiéyuē yòng shuǐ.)\\
« Nous ne devons pas jeter les ordures n'importe où, et nous devons aussi économiser l'eau. »\\
上个月，我们在小区里种了很多树。\\
(Shàng ge yuè, wǒmen zài xiǎoqū lǐ zhòng le hěn duō shù.)\\
« Le mois dernier, nous avons planté beaucoup d'arbres dans le quartier. »\\
现在我们的小区越来越干净、越来越美了。\\
(Xiànzài wǒmen de xiǎoqū yuè lái yuè gānjìng, yuè lái yuè měi le.)\\
« Maintenant, notre quartier devient de plus en plus propre et de plus en plus beau. »\\
我爱我的小区！\\
(Wǒ ài wǒ de xiǎoqū!)\\
« J'aime mon quartier ! »\par
(6 phrases, environ 85 caractères ; mots du glossaire utilisés : \zh{小区}, \zh{脏}, \zh{垃圾}, \zh{保护}, \zh{环境}, \zh{节约}, \zh{种}, \zh{干净}.)\par
\textbf{Points de vigilance :} (1) \zh{小区} (xiǎoqū) = « quartier résidentiel » ; (2) le complément de lieu \zh{在小区里} se place avant le verbe ; (3) ne pas écrire \zh{很越来越} : les deux sont incompatibles ; (4) vérifier les tons du pinyin : zāng (1er ton), měi (3e ton), zhòng (4e ton) ; (5) terminer par une phrase d'ouverture personnelle (\zh{我爱……}) pour donner une conclusion naturelle.
\end{corrige}

\newpage

\coursec{Fiche bilan}

\begin{popfigure}
\centering
\begin{tikzpicture}[mindmap, grow cyclic, every node/.style={concept, font=\small}, root concept/.append style={concept color=popgreen, font=\bfseries}, level 1/.append style={level distance=3.2cm, sibling angle=51}, text=white]
\node[root concept, align=center]{保护环境\\ \emph{bǎohù huánjìng}}
  child[concept color=popblue]{ node[align=center]{自然\\ zìrán\\ nature} }
  child[concept color=poporange]{ node[align=center]{污染\\ wūrǎn\\ pollution} }
  child[concept color=poppurple]{ node[align=center]{行动\\ xíngdòng\\ agir} }
  child[concept color=poppink]{ node[align=center]{垃圾\\ lājī\\ déchets} }
  child[concept color=popgold]{ node[align=center]{能源\\ néngyuán\\ énergie} }
  child[concept color=popgreen!80!black]{ node[align=center]{比较句\\ bǐjiào jù\\ comparatif} }
  child[concept color=popblue!70!black]{ node[align=center]{越来越\\ yuè lái yuè\\ de plus en plus} };
\end{tikzpicture}
\legende{Carte mentale de la leçon 15 : la protection de l'environnement}
\end{popfigure}

\begin{bilan}
\end{bilan}

\courssub{Lexique clé à connaître par cœur}
\begin{itemize}
  \item \zh{环境} \emph{huánjìng} : environnement ; \zh{保护} \emph{bǎohù} : protéger ; \zh{污染} \emph{wūrǎn} : pollution / polluer.
  \item \zh{空气} \emph{kōngqì} : air ; \zh{水} \emph{shuǐ} : eau ; \zh{树} \emph{shù} : arbre ; \zh{森林} \emph{sēnlín} : forêt.
  \item \zh{垃圾} \emph{lājī} : déchets ; \zh{扔} \emph{rēng} : jeter ; \zh{回收} \emph{huíshōu} : recycler.
  \item \zh{干净} \emph{gānjìng} : propre ; \zh{脏} \emph{zāng} : sale ; \zh{节约} \emph{jiéyuē} : économiser.
  \item \zh{地球} \emph{dìqiú} : la Terre ; \zh{能源} \emph{néngyuán} : énergie ; \zh{问题} \emph{wèntí} : problème.
\end{itemize}

\courssub{Structures grammaticales à maîtriser}
\begin{center}
\begin{tabularx}{\linewidth}{|X|X|X|}
\hline
\rowcolor{popblueL} Structure & Exemple & Traduction \\
\hline
A + 比 + B + Adj. & \zh{城市比农村脏。} \emph{Chéngshì bǐ nóngcūn zāng.} & La ville est plus sale que la campagne. \\
\hline
越来越 + Adj. & \zh{天气越来越热。} \emph{Tiānqì yuè lái yuè rè.} & Il fait de plus en plus chaud. \\
\hline
应该 + Verbe & \zh{我们应该保护环境。} \emph{Wǒmen yīnggāi bǎohù huánjìng.} & Nous devons protéger l'environnement. \\
\hline
把 + Nom + Verbe & \zh{把垃圾扔进垃圾桶。} \emph{Bǎ lājī rēng jìn lājī tǒng.} & Jeter les déchets dans la poubelle. \\
\hline
为了 + but & \zh{为了保护地球，我们种树。} \emph{Wèile bǎohù dìqiú, wǒmen zhòng shù.} & Pour protéger la Terre, nous plantons des arbres. \\
\hline
\end{tabularx}
\end{center}

\courssub{Méthodes express}
\begin{enumerate}
  \item \textbf{Lire un texte} : d'abord le titre et la traduction, puis repérer les mots connus, enfin deviner les mots nouveaux grâce au contexte.
  \item \textbf{Mémoriser un caractère} : identifier la clé (radical), compter les traits, écrire de haut en bas et de gauche à droite.
  \item \textbf{Répondre aux questions} : citer la phrase du texte qui justifie la réponse, répondre en phrase complète.
  \item \textbf{Comparer} : toujours placer \zh{比} entre les deux termes comparés, jamais avant le sujet.
\end{enumerate}


\begin{aretenir}[Checklist avant le Bac]
\begin{itemize}
  \item $\square$ Je sais lire et écrire \zh{环境}, \zh{保护}, \zh{污染}, \zh{垃圾} sans regarder le modèle.
  \item $\square$ Je connais la clé de \zh{河} (l'eau 氵) et celle de \zh{树} (le bois 木).
  \item $\square$ Je prononce correctement les tons de \zh{huánjìng} et \zh{wūrǎn}.
  \item $\square$ Je sais construire une phrase comparative avec \zh{比}.
  \item $\square$ Je sais utiliser \zh{越来越} pour exprimer une évolution.
  \item $\square$ Je sais exprimer un devoir avec \zh{应该}.
  \item $\square$ Je sais exprimer un but avec \zh{为了}.
  \item $\square$ Je sais employer la construction en \zh{把} avec un complément de résultat.
  \item $\square$ Je peux résumer le texte support en trois phrases simples.
  \item $\square$ Je peux citer deux actions concrètes de protection de l'environnement au Cameroun et en Chine.
  \item $\square$ Je respecte l'ordre des mots : sujet + verbe + complément, sans inversion à la française.
  \item $\square$ Je relis toujours mes tons avant de rendre ma copie.
\end{itemize}
\end{aretenir}

\findecours

\end{document}
